% =====================================================================
%  CAPITULO 3 - ASSOCIACAO DE RESISTORES E TRANSFORMACOES
%  (versao revisada e ampliada: 26 -> 48 questoes)
% =====================================================================
\section{Associação de Resistores e Transformações}

\begin{conceito}[Antes de começar]
\textbf{Série} --- a corrente é a \emph{mesma} em todos os elementos e as tensões
se somam:
\[
\Req = R_1+R_2+\dots+R_n ,\qquad U = U_1+U_2+\dots+U_n .
\]
\textbf{Paralelo} --- a tensão é a \emph{mesma} em todos os ramos e as correntes
se somam:
\[
\frac{1}{\Req}=\frac{1}{R_1}+\dots+\frac{1}{R_n},\qquad
\text{dois resistores: }\ \Req=\frac{R_1R_2}{R_1+R_2},\qquad
\text{$n$ iguais: }\ \Req=\frac{R}{n}.
\]
\textbf{Transformação $\Delta \to Y$} (resistores do triângulo no denominador comum
$S=R_{AB}+R_{BC}+R_{CA}$):
\[
R_A=\frac{R_{AB}\,R_{CA}}{S},\qquad
R_B=\frac{R_{AB}\,R_{BC}}{S},\qquad
R_C=\frac{R_{BC}\,R_{CA}}{S}.
\]
\textbf{Transformação $Y \to \Delta$} (com $P=R_AR_B+R_BR_C+R_CR_A$):
\[
R_{AB}=\frac{P}{R_C},\qquad R_{BC}=\frac{P}{R_A},\qquad R_{CA}=\frac{P}{R_B}.
\]
\end{conceito}

\begin{atencao}[Dois atalhos que economizam muito tempo]
\textbf{Curto-circuito:} um fio ideal em paralelo com um resistor \emph{anula} esse
resistor ($\Req = 0$). \quad
\textbf{Circuito aberto:} um ramo interrompido não conduz e pode ser apagado do
desenho. Antes de calcular, procure por esses dois casos.
\end{atencao}

\legendaniveis

% =====================================================================
\subsection{Circuitos em série}
% =====================================================================

\blocofixacao

\begin{questao}[1]{}
Em uma associação de resistores \textbf{em série}, é correto afirmar que:
\begin{alternativas}
\item a tensão é a mesma em todos os resistores.
\item a corrente é a mesma em todos os resistores.
\item a resistência equivalente é menor que a menor das resistências.
\item a corrente se divide proporcionalmente às resistências.
\item a potência dissipada é a mesma em todos os resistores.
\end{alternativas}
\resposta{(b)}
\resolucao{Só existe um caminho para a corrente, logo ela é a mesma em todos os
elementos --- (a) e (d) descrevem o circuito paralelo. A resistência equivalente
é a \emph{soma}, portanto maior que a maior delas, o que elimina (c).
A potência $P=R\,I^2$ depende de $R$: só é igual se os resistores forem iguais,
o que elimina (e).}
\end{questao}

\begin{questao}[1]{}
Três resistores de $\SI{4}{\ohm}$, $\SI{6}{\ohm}$ e $\SI{10}{\ohm}$ são ligados em
série a uma fonte de $\SI{40}{\volt}$ de resistência interna desprezível.
Determine:
\begin{itensq}
\item a resistência equivalente;
\item a corrente do circuito;
\item a tensão em cada resistor.
\end{itensq}
\resposta{(a) \SI{20}{\ohm}; (b) \SI{2}{\ampere}; (c) \SI{8}{\volt}, \SI{12}{\volt} e \SI{20}{\volt}}
\resolucao{(a) $\Req = 4+6+10 = \SI{20}{\ohm}$.\\
(b) $I = \dfrac{U}{\Req} = \dfrac{40}{20} = \SI{2}{\ampere}$.\\
(c) $U_1 = 4\cdot 2 = \SI{8}{\volt}$, $U_2 = 6\cdot 2 = \SI{12}{\volt}$,
$U_3 = 10\cdot 2 = \SI{20}{\volt}$.\\
\textbf{Verificação:} $8+12+20 = \SI{40}{\volt}$ --- confere com a tensão da fonte.
Repare que a maior tensão cai sobre o maior resistor.}
\end{questao}

\begin{questao}[1]{}
Para cada circuito, calcule a resistência equivalente e a corrente $I$.

\circuitoq{%
\subcirc{a}{%
\begin{circuitikz}[scale=0.78,transform shape,line width=0.8pt]
 \draw (0,0) to[battery1,l_={$E=\SI{10}{\volt}$}] (0,3)
       to[R,l={$R_1=\SI{3}{\ohm}$}] (3,3)
       to[R,l={$R_2=\SI{2}{\ohm}$}] (3,0) -- (0,0);
 \draw[->,Icol,line width=1pt] (0.7,3.4) -- (1.7,3.4)
       node[midway,above,font=\footnotesize]{$I$};
\end{circuitikz}}
\hfill
\subcirc{b}{%
\begin{circuitikz}[scale=0.78,transform shape,line width=0.8pt]
 \draw (0,0) to[battery1,l_={$E=\SI{15}{\volt}$}] (0,3)
       to[R,l={$R_1=\SI{4}{\ohm}$}] (3,3)
       to[R,l={$R_2=\SI{6}{\ohm}$}] (3,0) -- (0,0);
 \draw[->,Icol,line width=1pt] (0.7,3.4) -- (1.7,3.4)
       node[midway,above,font=\footnotesize]{$I$};
\end{circuitikz}}

\vspace{1.1em}

\subcirc{c}{%
\begin{circuitikz}[scale=0.78,transform shape,line width=0.8pt]
 \draw (0,0) to[battery1,l_={$E=\SI{25}{\volt}$}] (0,3)
       to[R,l={$R_1=\SI{4.5}{\ohm}$}] (3.4,3)
       to[R,l={$R_2=\SI{3}{\ohm}$}] (3.4,0)
       to[R,l={$R_3=\SI{2.5}{\ohm}$}] (0,0);
 \draw[->,Icol,line width=1pt] (0.7,3.4) -- (1.7,3.4)
       node[midway,above,font=\footnotesize]{$I$};
\end{circuitikz}}
\hfill
\subcirc{d}{%
\begin{circuitikz}[scale=0.78,transform shape,line width=0.8pt]
 \draw (0,0) to[battery1,l_={$E=\SI{30}{\volt}$}] (0,4.4)
       to[R,l={$R_1=\SI{2.5}{\ohm}$}] (3.6,4.4)
       to[R,l={$R_2=\SI{4}{\ohm}$}] (3.6,2.2)
       to[R,l={$R_3=\SI{1}{\ohm}$}] (3.6,0)
       to[R,l={$R_4=\SI{2.5}{\ohm}$}] (0,0);
 \draw[->,Icol,line width=1pt] (0.7,4.8) -- (1.7,4.8)
       node[midway,above,font=\footnotesize]{$I$};
\end{circuitikz}}}

\resposta{(a) \SI{5}{\ohm}, \SI{2}{\ampere}; (b) \SI{10}{\ohm}, \SI{1.5}{\ampere};
(c) \SI{10}{\ohm}, \SI{2.5}{\ampere}; (d) \SI{10}{\ohm}, \SI{3}{\ampere}}
\resolucao{Em todos os casos, $\Req$ é a soma das resistências e
$I = E/\Req$:\\
(a) $\Req = 3+2 = \SI{5}{\ohm}$, $I = 10/5 = \SI{2}{\ampere}$.\\
(b) $\Req = 4+6 = \SI{10}{\ohm}$, $I = 15/10 = \SI{1.5}{\ampere}$.\\
(c) $\Req = 4{,}5+3+2{,}5 = \SI{10}{\ohm}$, $I = 25/10 = \SI{2.5}{\ampere}$.\\
(d) $\Req = 2{,}5+4+1+2{,}5 = \SI{10}{\ohm}$, $I = 30/10 = \SI{3}{\ampere}$.}
\end{questao}

\begin{questao}[1]{}
Uma associação em série de $n$ resistores iguais de $\SI{15}{\ohm}$ apresenta
resistência equivalente de $\SI{105}{\ohm}$. O valor de $n$ é:
\begin{alternativas}[3]
\item 3
\item 5
\item 6
\item 7
\item 9
\end{alternativas}
\resposta{(d)}
\resolucao{Para $n$ resistores iguais em série, $\Req = n\,R$, logo
$n = \dfrac{105}{15} = 7$.}
\end{questao}

\blocoaplicacao

\begin{questao}[2]{}
No circuito da figura, duas fontes ideais estão ligadas em \textbf{oposição}.

\circuitoq{%
\begin{circuitikz}[scale=0.9,transform shape,line width=0.8pt]
 \draw (0,0) to[battery1,l_={$E_1=\SI{30}{\volt}$}] (0,3.6)
       to[R,l={$R_1=\SI{2.5}{\ohm}$}] (4.5,3.6)
       to[battery1,l={$E_2=\SI{10}{\volt}$},invert] (4.5,1.8)
       to[R,l={$R_2=\SI{1.5}{\ohm}$}] (4.5,0)
       to[R,l={$R_3=\SI{1}{\ohm}$}] (0,0);
 \draw[->,Icol,line width=1pt] (1.0,4.0) -- (2.2,4.0)
       node[midway,above,font=\footnotesize]{$I$};
\end{circuitikz}}

Determine a corrente do circuito e a tensão sobre $R_1$.
\resposta{$I = \SI{4}{\ampere}$; $U_{R_1} = \SI{10}{\volt}$}
\resolucao{Percorrendo a malha, $E_1$ impulsiona a corrente no sentido adotado e
$E_2$ se opõe. Pela lei das malhas:
\[
E_1 - E_2 = I\,(R_1+R_2+R_3)
\quad\Rightarrow\quad
I = \frac{30-10}{2{,}5+1{,}5+1} = \frac{20}{5} = \SI{4}{\ampere}.
\]
$U_{R_1} = R_1 I = 2{,}5\cdot 4 = \SI{10}{\volt}$.\\
\textbf{Observação:} se $E_2$ fosse maior que $E_1$, a corrente resultaria
negativa --- o que significa apenas que ela circula no sentido oposto ao adotado.
Nesse caso $E_2$ estaria \emph{sendo carregada} por $E_1$.}
\end{questao}

\begin{questao}[2]{}
Uma bateria de força eletromotriz $\EMF = \SI{12}{\volt}$ e resistência interna
$r = \SI{0.5}{\ohm}$ alimenta um resistor $R = \SI{5.5}{\ohm}$.

\circuitoq{%
\begin{circuitikz}[scale=0.9,transform shape,line width=0.8pt]
 \draw (0,0) to[battery1,l_={$\EMF=\SI{12}{\volt}$}] (0,2)
       to[R,l_={$r=\SI{0.5}{\ohm}$},color=Vcol] (0,3.4) -- (3.6,3.4)
       to[R,l={$R=\SI{5.5}{\ohm}$}] (3.6,0) -- (0,0);
 \draw[dashed,cinza] (-1.15,-0.5) rectangle (0.95,3.9);
 \node[cinza,font=\footnotesize\sffamily,above] at (-0.1,3.95) {bateria real};
 \draw[->,Icol,line width=1pt] (1.2,3.75) -- (2.3,3.75)
       node[midway,above,font=\footnotesize]{$I$};
\end{circuitikz}}

Determine:
\begin{itensq}
\item a corrente do circuito;
\item a tensão nos terminais da bateria;
\item a potência dissipada internamente.
\end{itensq}
\resposta{(a) \SI{2}{\ampere}; (b) \SI{11}{\volt}; (c) \SI{2}{\watt}}
\resolucao{(a) A resistência interna está \emph{em série} com a carga:
\[
I=\frac{\EMF}{R+r}=\frac{12}{5{,}5+0{,}5}=\frac{12}{6}=\SI{2}{\ampere}.
\]
(b) $U = \EMF - r\,I = 12 - 0{,}5\cdot 2 = \SI{11}{\volt}$
(ou, equivalentemente, $U = R\,I = 5{,}5\cdot 2 = \SI{11}{\volt}$).\\
(c) $P_r = r\,I^2 = 0{,}5\cdot 2^2 = \SI{2}{\watt}$ --- energia perdida dentro da
própria bateria, que é o que a faz aquecer em curtos-circuitos.}
\end{questao}

\begin{questao}[2]{}
Seis lâmpadas idênticas, de resistência $\SI{40}{\ohm}$ cada, são ligadas
\textbf{em série} a uma fonte de $\SI{120}{\volt}$.
\begin{itensq}
\item Calcule a corrente e a tensão em cada lâmpada.
\item Uma das lâmpadas queima (filamento aberto). O que acontece com as demais?
      Justifique.
\end{itensq}
\resposta{(a) $I=\SI{0.5}{\ampere}$ e $U=\SI{20}{\volt}$ por lâmpada;
(b) todas apagam}
\resolucao{(a) $\Req = 6\cdot 40 = \SI{240}{\ohm}$;
$I = \dfrac{120}{240} = \SI{0.5}{\ampere}$;
$U = 40\cdot 0{,}5 = \SI{20}{\volt}$ em cada lâmpada
(confira: $6\cdot 20 = \SI{120}{\volt}$).\\
(b) O filamento aberto interrompe o \emph{único} caminho da corrente:
$I = 0$ e todas as lâmpadas apagam. É exatamente esse o inconveniente dos antigos
cordões de luzes de Natal ligados em série.}
\end{questao}

\begin{questao}[2]{}
Três resistores estão em série sob $\SI{120}{\volt}$: $R_1=\SI{20}{\ohm}$,
$R_2$ desconhecido e $R_3=\SI{30}{\ohm}$. Deseja-se que a tensão sobre $R_2$
seja de $\SI{20}{\volt}$. O valor de $R_2$ deve ser:
\begin{alternativas}[3]
\item \SI{5}{\ohm}
\item \SI{10}{\ohm}
\item \SI{15}{\ohm}
\item \SI{20}{\ohm}
\item \SI{25}{\ohm}
\end{alternativas}
\resposta{(b)}
\resolucao{Se $U_2 = \SI{20}{\volt}$, sobram $120-20 = \SI{100}{\volt}$ para
$R_1$ e $R_3$, que somam \SI{50}{\ohm}. Como a corrente é a mesma em toda a série:
\[
I=\frac{100}{50}=\SI{2}{\ampere}
\qquad\Rightarrow\qquad
R_2=\frac{U_2}{I}=\frac{20}{2}=\SI{10}{\ohm}.
\]}
\end{questao}

\blocoaprofundamento

\begin{questao}[3]{}
Uma carga de $\SI{9.6}{\ohm}$ é alimentada por uma fonte de $\SI{24}{\volt}$ por
meio de um cabo de dois condutores, cada um com resistência de
$\SI{0.2}{\ohm}$.

\circuitoq{%
\begin{circuitikz}[scale=0.9,transform shape,line width=0.8pt]
 \draw (0,0) to[battery1,l_={$E=\SI{24}{\volt}$}] (0,3)
       to[R,l=$\SI{0.2}{\ohm}$,color=Vcol] (4.4,3)
       to[R,l={$R_c=\SI{9.6}{\ohm}$}] (4.4,0)
       to[R,l=$\SI{0.2}{\ohm}$,color=Vcol] (0,0);
 \node[Vcol,font=\footnotesize\sffamily] at (2.2,3.85) {condutores do cabo};
 \draw[->,Icol,line width=1pt] (0.5,2.55) -- (1.5,2.55)
       node[midway,below,font=\footnotesize]{$I$};
\end{circuitikz}}

Determine:
\begin{itensq}
\item a corrente na linha;
\item a queda de tensão total no cabo;
\item a tensão efetivamente entregue à carga;
\item o rendimento da transmissão, $\eta = P_{\text{carga}}/P_{\text{total}}$.
\end{itensq}
\resposta{(a) \SI{2.4}{\ampere}; (b) \SI{0.96}{\volt}; (c) \SI{23.04}{\volt};
(d) \SI{96}{\percent}}
\resolucao{Os dois condutores (ida e volta) ficam \emph{em série} com a carga:
\[
\Req = 0{,}2+9{,}6+0{,}2=\SI{10}{\ohm}
\quad\Rightarrow\quad
I=\frac{24}{10}=\SI{2.4}{\ampere}.
\]
(b) $U_{\text{cabo}} = 0{,}4\cdot 2{,}4 = \SI{0.96}{\volt}$.\\
(c) $U_{\text{carga}} = 24-0{,}96 = \SI{23.04}{\volt}$ (ou $9{,}6\cdot2{,}4$).\\
(d) Como a corrente é a mesma em tudo,
$\eta = \dfrac{R_c}{\Req} = \dfrac{9{,}6}{10} = 0{,}96 = \SI{96}{\percent}$.\\
\textbf{Por que isso importa:} a perda cresce com $I^2$. Se a mesma potência fosse
entregue com o dobro da tensão (e metade da corrente), a perda no cabo cairia à
quarta parte --- é o argumento a favor da transmissão em alta tensão.}
\end{questao}

\begin{questao}[3]{}
Dispõe-se de resistores de $\SI{100}{\ohm}$, todos idênticos, e de uma fonte
ideal de $\SI{50}{\volt}$. Deseja-se obter, a partir dessa fonte, uma saída de
exatamente $\SI{20}{\volt}$ em vazio (sem carga conectada à saída), usando apenas
uma associação série desses resistores como divisor de tensão.
\begin{itensq}
\item Qual é o menor número de resistores capaz de fornecer essa tensão?
\item Nessa configuração, quanto vale a corrente que circula pelo divisor?
\end{itensq}
\resposta{(a) 5 resistores (saída sobre 2 deles); (b) \SI{0.1}{\ampere}}
\resolucao{(a) Com $n$ resistores iguais em série, a tensão sobre $k$ deles é
\[
U_k = E\,\frac{k}{n} \quad\Rightarrow\quad 20 = 50\cdot\frac{k}{n}
\quad\Rightarrow\quad \frac{k}{n}=\frac{2}{5}.
\]
A menor solução inteira é $k=2$ e $n=5$: cinco resistores em série, com a saída
tomada sobre dois deles.\\
(b) $\Req = 5\cdot 100 = \SI{500}{\ohm}$ e
$I = \dfrac{50}{500} = \SI{0.1}{\ampere}$.\\
\textbf{Cuidado prático:} essa tensão de \SI{20}{\volt} só se mantém \emph{em
vazio}. Ao ligar uma carga em paralelo com os dois resistores, a tensão de saída
cai --- por isso divisores resistivos não servem como fonte de alimentação.}
\end{questao}

\begin{questao}[3]{}
Uma linha de transmissão de $\SI{2}{\kilo\meter}$ alimenta uma carga de
$\SI{10}{\ohm}$ a partir de uma fonte de $\SI{240}{\volt}$. O cabo é de cobre
($\rho=\SI{1.7e-8}{\ohm\meter}$) e deve-se escolher sua seção $A$ de modo que a
queda de tensão na linha não ultrapasse $\SI{4}{\percent}$ da tensão da fonte.
\begin{itensq}
\item Determine a resistência máxima admissível para o par de condutores.
\item Calcule a seção mínima do cabo.
\item Determine o rendimento da transmissão nessa condição e a potência perdida.
\end{itensq}
\resposta{(a) $R_{\text{linha}} \le \SI{0.417}{\ohm}$;
(b) $A \ge \SI{163}{\milli\meter\squared}$; (c) $\eta = \SI{96}{\percent}$,
$P_{\text{perdida}}\approx\SI{221}{\watt}$}
\resolucao{\textbf{(a)} A linha e a carga estão em série, formando um divisor de
tensão. Para a queda na linha ser no máximo \SI{4}{\percent}:
\[
\frac{R_{\text{linha}}}{R_{\text{linha}}+R_L}\le 0{,}04
\;\Longrightarrow\;
R_{\text{linha}} \le \frac{0{,}04}{0{,}96}\,R_L = \frac{10}{24}
\approx\SI{0.417}{\ohm}.
\]
\textbf{(b)} O comprimento \emph{total} de condutor é o dobro da distância
(ida e volta): $L = \SI{4000}{\meter}$. De $R=\rho L/A$:
\[
A \ge \frac{\rho L}{R_{\text{linha}}}
=\frac{(\pot{1.7}{-8})(4000)}{0{,}417}
\approx\pot{1.63}{-4}\,\si{\meter\squared}=\SI{163}{\milli\meter\squared}.
\]
Comercialmente, adotaria-se $\SI{185}{\milli\meter\squared}$ (valor normalizado
imediatamente superior).

\textbf{(c)} Com $R_{\text{linha}} = \SI{0.417}{\ohm}$:
\[
I=\frac{240}{10{,}417}\approx\SI{23}{\ampere},
\qquad
\eta=\frac{R_L}{R_L+R_{\text{linha}}}=\frac{10}{10{,}417}\approx\SI{96}{\percent},
\]
\[
P_{\text{perdida}}=R_{\text{linha}}I^{2}\approx 0{,}417\cdot23^{2}
\approx\SI{221}{\watt}.
\]
\textbf{Reflexão de projeto:} note o custo do critério. Exigir apenas
\SI{4}{\percent} de queda impôs um cabo de $\SI{163}{\milli\meter\squared}$ ---
muito caro. Se a mesma potência fosse transmitida em $\SI{2400}{\volt}$
(dez vezes mais), a corrente cairia a $1/10$ e a perda ($\propto I^2$) a
$1/100$, permitindo um cabo cem vezes menor. \textbf{É por isso que se transmite
energia em alta tensão} --- a economia está no cobre, não na energia.}
\end{questao}

% BEGIN BANCO COMPLEMENTAR Circuitos em série
% =====================================================================
%  Banco complementar --- Circuitos Serie  (REVISADO)
%  Substitui a versao gerada por template (11 clones em N2, 12 em N3,
%  10 questoes conceituais indevidamente marcadas como N4).
%  Sem sobreposicao com o cap03 (fonte em oposicao, seis lampadas,
%  queda em cabo e linha de transmissao ja estao la).
%  Distribuicao mantida: 6 N1 + 11 N2 + 12 N3 + 10 N4 = 39
% =====================================================================

\subsubsection*{Banco complementar --- Circuitos em série}

\begin{atencao}[Organização do banco]
Toda associação série repousa em dois fatos: a \textbf{corrente é a mesma} em
todos os elementos e as \textbf{quedas se somam}. Deles decorre que a tensão e
a potência se distribuem \emph{proporcionalmente} a $R$. O N3 explora falhas,
temperatura e tolerância; o N4 trata de projeto, sensibilidade, propagação de
incerteza e otimização.
\end{atencao}

\blocofixacao

\begin{questao}[1]{}
Determine a resistência equivalente de $R_1=\SI{1.2}{\kilo\ohm}$,
$R_2=\SI{470}{\ohm}$ e $R_3=\SI{2.2}{\kilo\ohm}$ em série.
\resposta{$R_{eq}=\SI{3.87}{\kilo\ohm}$}
\resolucao{Converta tudo para a mesma unidade \emph{antes} de somar:
\[
R_{eq}=1200+470+2200=\SI{3870}{\ohm}=\SI{3.87}{\kilo\ohm}.
\]
Em série a soma é direta --- não há inversos envolvidos. Misturar
$\si{\kilo\ohm}$ com $\si{\ohm}$ sem converter é o erro mais frequente aqui.}
\end{questao}

\begin{questao}[1]{}
Uma associação série de três resistores tem $R_{eq}=\SI{75}{\ohm}$. Sabendo que
$R_1=\SI{22}{\ohm}$ e $R_2=\SI{33}{\ohm}$, determine $R_3$.
\resposta{$R_3=\SI{20}{\ohm}$}
\resolucao{
\[
R_3=R_{eq}-R_1-R_2=75-22-33=\SI{20}{\ohm}.
\]
A soma série é reversível: conhecidos o total e todas as parcelas menos uma,
a que falta sai por subtração --- diferente do paralelo, em que seria preciso
subtrair \emph{condutâncias}.}
\end{questao}

\begin{questao}[1]{}
Uma fonte ideal de $\SI{24}{\volt}$ alimenta $R_1=\SI{4}{\ohm}$ em série com
$R_2=\SI{8}{\ohm}$. Determine a corrente do circuito.
\resposta{$I=\SI{2}{\ampere}$}
\resolucao{
\[
R_{eq}=4+8=\SI{12}{\ohm}
\;\Longrightarrow\;
I=\frac{V}{R_{eq}}=\frac{24}{12}=\SI{2}{\ampere}.
\]
Essa corrente é a \emph{mesma} nos dois resistores --- não existe "corrente de
$R_1$" diferente da "corrente de $R_2$" em uma série.}
\end{questao}

\begin{questao}[1]{}
Em uma cadeia série percorrida por $\SI{0.5}{\ampere}$, determine a queda de
tensão em um resistor de $\SI{68}{\ohm}$.
\resposta{$U=\SI{34}{\volt}$}
\resolucao{Lei de Ohm aplicada ao elemento:
\[
U=R\,I=68\cdot0{,}5=\SI{34}{\volt}.
\]
Como $I$ é comum, a queda de cada resistor é \emph{diretamente proporcional} ao
seu valor: o maior resistor sempre fica com a maior parcela da tensão.}
\end{questao}

\begin{questao}[1]{}
Acrescentando-se mais um resistor \textbf{em série} a um circuito alimentado
por fonte ideal, a corrente total:
\begin{alternativas}[2]
\item aumenta
\item diminui
\item não se altera
\item pode aumentar ou diminuir, conforme o valor
\end{alternativas}
\resposta{(b)}
\resolucao{$R_{eq}$ só pode crescer ao se acrescentar um resistor em série
($R_{eq}=\sum R_k$, com todas as parcelas positivas). Com $V$ fixo,
$I=V/R_{eq}$ necessariamente diminui. Consequência prática: em uma cadeia de
lâmpadas, acrescentar mais uma faz \emph{todas} brilharem menos.}
\end{questao}

\begin{questao}[1]{}
Dois resistores de $\SI{15}{\ohm}$ e $\SI{45}{\ohm}$ são ligados em série a uma
fonte. Trocando-se a \textbf{ordem} deles no circuito, o que muda?
\begin{alternativas}[2]
\item a resistência equivalente
\item a corrente do circuito
\item a queda em cada resistor
\item nada
\end{alternativas}
\resposta{(d)}
\resolucao{A soma $R_{eq}=15+45=\SI{60}{\ohm}$ é comutativa, a corrente
$I=V/R_{eq}$ depende só de $R_{eq}$, e cada queda $U_k=R_kI$ depende só do
próprio resistor e de $I$. Nada muda.\\
O que muda é o \emph{potencial} de cada ponto intermediário em relação ao terra
--- relevante apenas se algum ponto do circuito estiver aterrado ou se houver
medição referida a um nó específico.}
\end{questao}

\blocoaplicacao

\begin{questao}[2]{}
$R_1=\SI{5}{\ohm}$, $R_2=\SI{10}{\ohm}$ e $R_3=\SI{15}{\ohm}$ estão em série
sob $\SI{60}{\volt}$. Determine a corrente e as três quedas de tensão.
\resposta{$I=\SI{2}{\ampere}$; $U_1=\SI{10}{\volt}$, $U_2=\SI{20}{\volt}$,
$U_3=\SI{30}{\volt}$}
\resolucao{
\[
R_{eq}=5+10+15=\SI{30}{\ohm}
\;\Longrightarrow\;
I=\frac{60}{30}=\SI{2}{\ampere}.
\]
\[
U_1=5\cdot2=10,\quad U_2=10\cdot2=20,\quad U_3=15\cdot2=\SI{30}{\volt}.
\]
\textbf{Verificação (LKT):} $10+20+30=\SI{60}{\volt}$ --- fecha com a fonte.
Note a proporcionalidade: as quedas estão entre si como $5:10:15=1:2:3$.}
\end{questao}

\begin{questao}[2]{}
No circuito da questão anterior, calcule a potência dissipada em cada resistor
e confronte com a potência fornecida pela fonte.
\resposta{$\SI{20}{\watt}$, $\SI{40}{\watt}$, $\SI{60}{\watt}$; total
$\SI{120}{\watt}$}
\resolucao{Com $I=\SI{2}{\ampere}$ comum, use $P=RI^{2}$:
\[
P_1=5\cdot4=20,\quad P_2=10\cdot4=40,\quad P_3=15\cdot4=\SI{60}{\watt}.
\]
Fonte: $P=VI=60\cdot2=\SI{120}{\watt}=20+40+60$ \checkmark.\\
\textbf{Regra da série:} como $I$ é comum, $P\propto R$. O maior resistor
dissipa mais --- exatamente o oposto do que ocorre no paralelo, onde $V$ é
comum e $P\propto 1/R$.}
\end{questao}

\begin{questao}[2]{}
Em uma cadeia série de três resistores, mediu-se $\SI{18}{\volt}$ sobre o
resistor de $\SI{90}{\ohm}$. Os outros dois valem $\SI{30}{\ohm}$ e
$\SI{60}{\ohm}$. Determine a tensão da fonte.
\resposta{$V=\SI{36}{\volt}$}
\resolucao{A medida sobre um único elemento revela a corrente da cadeia
inteira:
\[
I=\frac{18}{90}=\SI{0.2}{\ampere}.
\]
\[
R_{eq}=30+60+90=\SI{180}{\ohm}
\;\Longrightarrow\;
V=R_{eq}I=180\cdot0{,}2=\SI{36}{\volt}.
\]
Este é o procedimento padrão de bancada: medir a queda em um resistor conhecido
transforma-o em um \emph{sensor de corrente} (resistor \emph{shunt}), sem
abrir o circuito para inserir o amperímetro.}
\end{questao}

\begin{questao}[2]{}
Duas lâmpadas incandescentes, de resistências $\SI{240}{\ohm}$ e
$\SI{960}{\ohm}$, são ligadas \textbf{em série} a $\SI{120}{\volt}$. Qual
brilha mais e por quê?
\resposta{A de $\SI{960}{\ohm}$; $P=\SI{9.6}{\watt}$ contra $\SI{2.4}{\watt}$}
\resolucao{
\[
I=\frac{120}{240+960}=\frac{120}{1200}=\SI{0.1}{\ampere}
\]
\[
P_{240}=240\cdot(0{,}1)^{2}=\SI{2.4}{\watt},\qquad
P_{960}=960\cdot(0{,}1)^{2}=\SI{9.6}{\watt}.
\]
Como a corrente é comum, $P\propto R$: a lâmpada de \emph{maior} resistência
brilha mais. Isso contraria a intuição de quem associa "mais potência" a "menos
resistência" --- válido apenas quando a \emph{tensão} é comum, isto é, no
paralelo. Uma lâmpada de $\SI{100}{\watt}$ e outra de $\SI{25}{\watt}$ ligadas
em série: a de $\SI{25}{\watt}$ é que brilha mais.}
\end{questao}

\begin{questao}[2]{}
Deseja-se limitar a $\SI{150}{\milli\ampere}$ a corrente de uma carga de
$\SI{60}{\ohm}$ alimentada por $\SI{24}{\volt}$. Determine o resistor em série
necessário e sua potência.
\resposta{$R_s=\SI{100}{\ohm}$; $P_s=\SI{2.25}{\watt}$}
\resolucao{
\[
R_{eq}=\frac{V}{I}=\frac{24}{0{,}150}=\SI{160}{\ohm}
\;\Longrightarrow\;
R_s=160-60=\SI{100}{\ohm}.
\]
\[
P_s=R_sI^{2}=100\cdot(0{,}15)^{2}=\SI{2.25}{\watt}.
\]
Especifique pelo menos $\SI{5}{\watt}$ (fator $2\times$ de folga). Observe que o
resistor limitador dissipa $\SI{2.25}{\watt}$ dos $\SI{3.6}{\watt}$ totais ---
$62\%$ da energia vira calor. Limitar corrente com resistor é simples, mas
nunca eficiente.}
\end{questao}

\begin{questao}[2]{}
Duas montagens têm a mesma $R_{eq}=\SI{100}{\ohm}$ sob $\SI{50}{\volt}$:
(A) dois resistores de $\SI{50}{\ohm}$; (B) um de $\SI{10}{\ohm}$ e outro de
$\SI{90}{\ohm}$. Compare corrente total, potência total e a potência do
resistor mais solicitado.
\resposta{Corrente e potência totais iguais ($\SI{0.5}{\ampere}$,
$\SI{25}{\watt}$); resistor mais solicitado: $\SI{12.5}{\watt}$ em (A) e
$\SI{22.5}{\watt}$ em (B)}
\resolucao{Em ambas: $I=50/100=\SI{0.5}{\ampere}$ e
$P_{tot}=50\cdot0{,}5=\SI{25}{\watt}$.\\
(A) $P=50\cdot0{,}25=\SI{12.5}{\watt}$ em cada.\\
(B) $P_{10}=\SI{2.5}{\watt}$ e $P_{90}=\SI{22.5}{\watt}$.\\
\textbf{Conclusão de projeto:} $R_{eq}$ igual \emph{não} significa
especificação igual. A montagem (B) exige um resistor de pelo menos
$\SI{25}{\watt}$, enquanto (A) se resolve com dois de $\SI{15}{\watt}$.
Distribuir a resistência em parcelas iguais distribui também o aquecimento.}
\end{questao}

\begin{questao}[2]{}
Uma cadeia série de $\SI{48}{\ohm}$ é alimentada por $\SI{12}{\volt}$ e
protegida por fusível. Determine a corrente normal de operação e justifique a
escolha de um fusível de $\SI{500}{\milli\ampere}$.
\resposta{$I=\SI{250}{\milli\ampere}$; o fusível deve ficar acima da corrente
normal e abaixo da corrente de falha}
\resolucao{
\[
I=\frac{12}{48}=\SI{0.25}{\ampere}=\SI{250}{\milli\ampere}.
\]
Um fusível de $\SI{500}{\milli\ampere}$ dá folga de $2\times$: não atua em
transitórios normais de energização, mas interrompe bem antes das correntes de
curto. \textbf{Estimativa da falha:} se metade da cadeia entrar em curto,
$R=\SI{24}{\ohm}$ e $I=\SI{500}{\milli\ampere}$ --- exatamente no limiar; um
curto franco ($R\to0$) levaria a corrente a valores muito maiores, e o fusível
atua rapidamente. Um fusível \emph{muito} próximo da corrente nominal produz
desarmes espúrios; muito acima, não protege.}
\end{questao}

\begin{questao}[2]{}
Um reostato de $\SI{0}{\ohm}$ a $\SI{100}{\ohm}$ está em série com uma carga de
$\SI{50}{\ohm}$, sob $\SI{30}{\volt}$. Determine a faixa de variação da
corrente.
\resposta{de $\SI{0.2}{\ampere}$ a $\SI{0.6}{\ampere}$}
\resolucao{
\[
I_{\max}=\frac{30}{50+0}=\SI{0.6}{\ampere},
\qquad
I_{\min}=\frac{30}{50+100}=\SI{0.2}{\ampere}.
\]
A faixa é $3{:}1$ --- \textbf{não} $\infty{:}1$, porque a carga fixa de
$\SI{50}{\ohm}$ impõe um teto à corrente por mais que o reostato seja reduzido.
Para ampliar a faixa seria preciso um reostato de valor maior em relação à
carga. Note também que a relação $I(R)$ é hiperbólica, não linear: o controle é
muito mais "sensível" perto de $R=0$.}
\end{questao}

\begin{questao}[2]{}
Em uma cadeia série sob tensão fixa, se um dos resistores tiver seu valor
\textbf{dobrado}, a queda sobre ele:
\begin{alternativas}[2]
\item dobra
\item aumenta, mas menos que o dobro
\item não se altera
\item diminui
\end{alternativas}
\resposta{(b)}
\resolucao{A queda é $U_k=V\dfrac{R_k}{R_{eq}}$. Dobrar $R_k$ dobra o
numerador, mas \emph{também} aumenta o denominador --- o crescimento é menos
que proporcional.\\
\textbf{Exemplo:} $R_1=R_2=\SI{10}{\ohm}$ sob $\SI{12}{\volt}$ dá
$U_1=\SI{6}{\volt}$. Com $R_1=\SI{20}{\ohm}$: $U_1=12\cdot20/30=\SI{8}{\volt}$
--- aumento de $33\%$, não de $100\%$. No limite $R_1\to\infty$, $U_1\to V$: a
queda satura no valor da fonte, nunca o ultrapassa.}
\end{questao}

\begin{questao}[2]{}
Uma cadeia série dissipa $\SI{300}{\watt}$ e opera $\SI{6}{\hour}$ por dia
durante $\SI{30}{\day}$. Calcule a energia consumida e o custo, a
$\mathrm{R\$}\,0{,}80$ por $\si{\kilo\watt\hour}$.
\resposta{$E=\SI{54}{\kilo\watt\hour}$; custo $=\mathrm{R\$}\,43{,}20$}
\resolucao{
\[
E=P\,t=0{,}300\,\si{\kilo\watt}\times6\times30
=\SI{54}{\kilo\watt\hour}.
\]
\[
\text{Custo}=54\times0{,}80=\mathrm{R\$}\,43{,}20.
\]
A conta não depende de \emph{como} a resistência está distribuída na cadeia ---
só da potência total. A distribuição importa para o dimensionamento térmico de
cada componente, não para a fatura.}
\end{questao}

\begin{questao}[2]{}
Em uma cadeia série sob $\SI{100}{\volt}$ mediram-se as quedas
$\SI{22}{\volt}$, $\SI{35}{\volt}$ e $\SI{41}{\volt}$. A medição é consistente?
Se a corrente for $\SI{50}{\milli\ampere}$, determine os três resistores.
\resposta{Sim ($22+35+41=98\approx100$, dentro de $2\%$);
$\SI{440}{\ohm}$, $\SI{700}{\ohm}$ e $\SI{820}{\ohm}$}
\resolucao{\textbf{Consistência (LKT):} a soma das quedas deve igualar a tensão
da fonte. Aqui $22+35+41=\SI{98}{\volt}$, desvio de $\SI{2}{\volt}$ ($2\%$) ---
compatível com a exatidão típica de um multímetro de bancada. Um desvio de,
digamos, $\SI{20}{\volt}$ indicaria erro de medida ou um quarto elemento não
contabilizado (resistência de contato, por exemplo).\\
\textbf{Resistores:} com $I=\SI{0.050}{\ampere}$ comum a todos,
\[
R_1=\frac{22}{0{,}05}=440,\quad
R_2=\frac{35}{0{,}05}=700,\quad
R_3=\frac{41}{0{,}05}=\SI{820}{\ohm}.
\]}
\end{questao}

\blocoaprofundamento

\begin{questao}[3]{}
$R_1=\SI{10}{\ohm}$ e $R_2=\SI{20}{\ohm}$ estão em série sob
$\SI{30}{\volt}$.
\begin{itensq}
\item Calcule a potência de cada um.
\item Especifique a potência nominal comercial de cada resistor, com fator de
      segurança $2$.
\item Explique por que especificar ambos igualmente seria desperdício ou risco.
\end{itensq}
\resposta{(a) $\SI{10}{\watt}$ e $\SI{20}{\watt}$; (b) $\SI{25}{\watt}$ e
$\SI{50}{\watt}$}
\resolucao{(a) $I=30/30=\SI{1}{\ampere}$; $P_1=10\cdot1=\SI{10}{\watt}$,
$P_2=20\cdot1=\SI{20}{\watt}$.\\
(b) Com fator $2$: $\SI{20}{\watt}$ e $\SI{40}{\watt}$; os valores comerciais
imediatamente superiores são $\SI{25}{\watt}$ e $\SI{50}{\watt}$.\\
(c) Especificar \emph{ambos} como $\SI{25}{\watt}$ deixaria $R_2$ operando a
$80\%$ do nominal --- sem margem para elevação de temperatura ambiente ou
ventilação deficiente. Especificar \emph{ambos} como $\SI{50}{\watt}$ dobraria
volume e custo de $R_1$ sem ganho. \textbf{Regra:} em série, $P\propto R$;
dimensione resistor a resistor.}
\end{questao}

\begin{questao}[3]{}
Um resistor de fio tem $\SI{50}{\ohm}$ a $\SI{20}{\celsius}$ e coeficiente
térmico $\alpha=\SI{0.004}{\per\celsius}$. Ele está em série com um resistor de
filme de $\SI{50}{\ohm}$ e $\alpha\approx0$, sob $\SI{124}{\volt}$.
\begin{itensq}
\item Determine a resistência do resistor de fio a $\SI{80}{\celsius}$.
\item Calcule a corrente a $\SI{20}{\celsius}$ e a $\SI{80}{\celsius}$.
\item Comente a variação percentual.
\end{itensq}
\resposta{(a) $\SI{62}{\ohm}$; (b) $\SI{1.24}{\ampere}$ e
$\SI{1.107}{\ampere}$; (c) queda de $10{,}7\%$}
\resolucao{(a)
$R=R_0[1+\alpha(T-T_0)]=50[1+0{,}004(60)]=50\cdot1{,}24=\SI{62}{\ohm}$.\\
(b) A $\SI{20}{\celsius}$: $R_{eq}=\SI{100}{\ohm}$ e
$I=124/100=\SI{1.24}{\ampere}$.\\
A $\SI{80}{\celsius}$: $R_{eq}=62+50=\SI{112}{\ohm}$ e
$I=124/112=\SI{1.107}{\ampere}$.\\
(c) Redução de $10{,}7\%$. Note que o resistor de fio variou $24\%$, mas a
corrente variou apenas $10{,}7\%$ --- o resistor de filme, estável, "dilui" o
efeito. Essa é a base do item de projeto explorado adiante: combinar
coeficientes para estabilizar a cadeia inteira.}
\end{questao}

\begin{questao}[3]{}
$R_1=\SI{100}{\ohm}\pm5\%$ e $R_2=\SI{200}{\ohm}\pm5\%$ estão em série.
\begin{itensq}
\item Determine $R_{eq}$ e seus limites de pior caso.
\item Expresse a tolerância resultante em porcentagem.
\item Repita supondo $R_2$ com tolerância de $1\%$.
\end{itensq}
\resposta{(a) $\SI{300}{\ohm}$, entre $\SI{285}{\ohm}$ e $\SI{315}{\ohm}$;
(b) $\pm5\%$; (c) $\pm2{,}33\%$}
\resolucao{(a) Em série as incertezas \emph{absolutas} se somam:
$\Delta R=5+10=\SI{15}{\ohm}$, logo $300\pm15$, isto é, de
$\SI{285}{\ohm}$ a $\SI{315}{\ohm}$.\\
(b) $15/300=5\%$ --- \emph{idêntica} às parcelas.\\
(c) Agora $\Delta R=5+2=\SI{7}{\ohm}$, e $7/300=2{,}33\%$.\\
\textbf{Padrão geral:} a tolerância relativa da série é a \emph{média
ponderada} das tolerâncias individuais, com pesos $R_k/R_{eq}$:
\[
t_{eq}=\frac{\sum R_k t_k}{\sum R_k}.
\]
Em (b), $t_1=t_2$ e a média devolve o mesmo valor. Em (c), o resistor
\emph{maior} domina o resultado --- por isso, ao apertar tolerâncias, comece
sempre pela maior parcela da cadeia.}
\end{questao}

\begin{questao}[3]{}
Na cadeia $R_1=\SI{10}{\ohm}$, $R_2=\SI{20}{\ohm}$, $R_3=\SI{30}{\ohm}$ sob
$\SI{60}{\volt}$, o resistor $R_2$ sofre \textbf{curto-circuito}. Compare
corrente e potências antes e depois.
\resposta{$I$: $\SI{1}{\ampere}\to\SI{1.5}{\ampere}$;
$P_1$: $10\to\SI{22.5}{\watt}$; $P_3$: $30\to\SI{67.5}{\watt}$}
\resolucao{\textbf{Antes:} $R_{eq}=\SI{60}{\ohm}$, $I=\SI{1}{\ampere}$,
$P=(10;20;30)\,\si{\watt}$, total $\SI{60}{\watt}$.\\
\textbf{Depois:} $R_{eq}=10+30=\SI{40}{\ohm}$, $I=60/40=\SI{1.5}{\ampere}$.
\[
P_1=10(1{,}5)^{2}=\SI{22.5}{\watt},\qquad P_3=30(1{,}5)^{2}=\SI{67.5}{\watt},
\]
total $\SI{90}{\watt}$ --- $50\%$ a mais consumido da fonte.\\
\textbf{Ponto crítico:} $P_1$ e $P_3$ subiram $125\%$ (fator $1{,}5^{2}=2{,}25$).
Um curto em \emph{um} componente pode destruir os \emph{outros}, que passam a
operar muito acima do nominal. Essa é a razão de as falhas em cadeia serem
frequentemente sequenciais, e não simultâneas.}
\end{questao}

\begin{questao}[3]{}
Na mesma cadeia ($\SI{10}{\ohm}$, $\SI{20}{\ohm}$, $\SI{30}{\ohm}$ sob
$\SI{60}{\volt}$), agora $R_2$ \textbf{abre}. Determine a corrente e a leitura
de um voltímetro ideal colocado sobre cada resistor.
\resposta{$I=0$; $\SI{0}{\volt}$ em $R_1$ e $R_3$; $\SI{60}{\volt}$ em $R_2$}
\resolucao{Com o circuito interrompido, $I=0$ em toda a série. Pela lei de Ohm,
$U=RI=0$ nos resistores íntegros --- mesmo estando eles perfeitamente
conectados à fonte.\\
Pela LKT, a soma das quedas deve continuar valendo $\SI{60}{\volt}$; como as
outras são nulas, \textbf{toda} a tensão da fonte aparece sobre a
descontinuidade: o voltímetro sobre $R_2$ lê $\SI{60}{\volt}$.\\
\textbf{Uso prático:} é assim que se localiza o componente aberto em uma cadeia
--- percorre-se a série com o voltímetro e o defeito é o único ponto onde
aparece tensão. Note que o voltímetro precisa ser de alta impedância; um
instrumento de baixa impedância fecharia o circuito e falsearia a leitura.}
\end{questao}

\begin{questao}[3]{}
Um LED com queda direta de $\SI{2.0}{\volt}$ deve operar com
$\SI{15}{\milli\ampere}$ a partir de uma fonte de $\SI{5.0}{\volt}$.
\begin{itensq}
\item Determine o resistor em série ideal.
\item Escolha o valor comercial (série E12) e recalcule a corrente real.
\item Determine a potência do resistor.
\end{itensq}
\resposta{(a) $\SI{200}{\ohm}$; (b) $\SI{220}{\ohm}\Rightarrow
\SI{13.6}{\milli\ampere}$; (c) $\SI{41}{\milli\watt}$}
\resolucao{(a) O resistor absorve a diferença entre a fonte e a queda do LED:
\[
R=\frac{V_{fonte}-V_{LED}}{I}=\frac{5{,}0-2{,}0}{0{,}015}=\SI{200}{\ohm}.
\]
(b) A E12 oferece $\SI{180}{\ohm}$ e $\SI{220}{\ohm}$. Escolha o
\textbf{maior} --- errar para menos aumenta a corrente e reduz a vida do LED:
\[
I=\frac{3{,}0}{220}=\SI{13.6}{\milli\ampere}.
\]
(c) $P=R I^{2}=220(0{,}0136)^{2}=\SI{41}{\milli\watt}$; um resistor de
$\frac{1}{8}\,\si{\watt}$ já basta com folga generosa.\\
\textbf{Observação:} o LED \emph{não} obedece à lei de Ohm --- sua queda é
aproximadamente constante. Por isso ele entra na conta como uma "fonte de
tensão" a ser subtraída, e não como uma resistência a ser somada.}
\end{questao}

\begin{questao}[3]{}
Uma bateria de $\EMF=\SI{12}{\volt}$ e resistência interna $r=\SI{0.5}{\ohm}$
alimenta uma carga de $\SI{9.5}{\ohm}$.
\begin{itensq}
\item Determine a corrente e a tensão nos terminais.
\item Calcule a regulação de tensão percentual.
\item Determine a fração da potência perdida internamente.
\end{itensq}
\resposta{(a) $\SI{1.2}{\ampere}$, $\SI{11.4}{\volt}$; (b) $5{,}26\%$;
(c) $5\%$}
\resolucao{(a) $I=\dfrac{12}{9{,}5+0{,}5}=\SI{1.2}{\ampere}$;
$V_{term}=12-0{,}5(1{,}2)=\SI{11.4}{\volt}$.\\
(b) Tomando a tensão em vazio como referência de partida:
\[
\text{Reg}=\frac{V_{vazio}-V_{carga}}{V_{carga}}
=\frac{12-11{,}4}{11{,}4}=5{,}26\%.
\]
(c) $P_r=rI^{2}=0{,}5(1{,}44)=\SI{0.72}{\watt}$ contra
$P_{tot}=12(1{,}2)=\SI{14.4}{\watt}$: exatamente $5\%$, que é a razão
$r/(r+R_L)$.\\
\textbf{Compromisso:} reduzir $r$ melhora regulação \emph{e} rendimento
simultaneamente --- por isso $r$ baixo é sempre desejável em fontes. É diferente
da máxima transferência de potência, que exigiria $R_L=r$ e um rendimento de
apenas $50\%$.}
\end{questao}

\begin{questao}[3]{}
Duas cadeias têm $R_{eq}=\SI{120}{\ohm}$ sob $\SI{120}{\volt}$: (A) quatro
resistores de $\SI{30}{\ohm}$; (B) um de $\SI{100}{\ohm}$ e dois de
$\SI{10}{\ohm}$. Compare a potência do componente mais solicitado e a potência
total, e escolha a montagem preferível.
\resposta{Total $\SI{120}{\watt}$ em ambas; pior componente:
$\SI{30}{\watt}$ em (A) e $\SI{100}{\watt}$ em (B) --- (A) é preferível}
\resolucao{$I=120/120=\SI{1}{\ampere}$ nas duas.\\
(A) $P=30\cdot1=\SI{30}{\watt}$ em cada um dos quatro.\\
(B) $P_{100}=\SI{100}{\watt}$; $P_{10}=\SI{10}{\watt}$ (dois).\\
Total: $\SI{120}{\watt}$ nos dois casos --- a fonte não distingue as montagens.\\
\textbf{Decisão:} (A) exige quatro resistores de $\SI{50}{\watt}$ nominais;
(B) exige um de $\SI{200}{\watt}$, componente caro, volumoso e com ponto quente
concentrado. Sempre que possível, \textbf{divida a dissipação} em parcelas
iguais: a área de troca térmica cresce e a temperatura máxima cai.}
\end{questao}

\begin{questao}[3]{}
Um termistor NTC de $\SI{120}{\ohm}$ (frio) e $\SI{5}{\ohm}$ (quente) é
colocado em série com uma carga de $\SI{20}{\ohm}$, sob $\SI{100}{\volt}$.
Determine a corrente no instante da energização e em regime, e explique a
função do componente.
\resposta{$\SI{0.71}{\ampere}$ na partida e $\SI{4.0}{\ampere}$ em regime}
\resolucao{\textbf{Partida (frio):}
$I=\dfrac{100}{120+20}=\SI{0.71}{\ampere}$.\\
\textbf{Regime (aquecido):}
$I=\dfrac{100}{5+20}=\SI{4.0}{\ampere}$.\\
\textbf{Função:} sem o NTC, a corrente de energização seria imediatamente
$100/20=\SI{5}{\ampere}$. O termistor a limita a $\SI{0.71}{\ampere}$ e vai
"saindo do caminho" à medida que aquece --- um \textbf{limitador de corrente de
partida}. É o que protege capacitores de entrada e contatos de relé em fontes
chaveadas.\\
\textbf{Limitação:} em regime o NTC ainda dissipa $5(4)^{2}=\SI{80}{\watt}$, e
precisa de tempo para esfriar antes de um novo religamento --- por isso ele não
protege contra reenergizações rápidas e sucessivas.}
\end{questao}

\begin{questao}[3]{}
Em uma cadeia de três resistores sob $\SI{90}{\volt}$, sabe-se que
$R_2=2R_1$, $R_3=3R_1$ e que a corrente é $\SI{0.5}{\ampere}$. Determine os
três valores e as quedas.
\resposta{$R_1=\SI{30}{\ohm}$, $R_2=\SI{60}{\ohm}$, $R_3=\SI{90}{\ohm}$;
quedas $\SI{15}{\volt}$, $\SI{30}{\volt}$, $\SI{45}{\volt}$}
\resolucao{
\[
R_{eq}=\frac{V}{I}=\frac{90}{0{,}5}=\SI{180}{\ohm}.
\]
Como $R_{eq}=R_1+2R_1+3R_1=6R_1$:
\[
R_1=\frac{180}{6}=\SI{30}{\ohm},\quad R_2=\SI{60}{\ohm},\quad R_3=\SI{90}{\ohm}.
\]
Quedas: $U_k=R_kI$, ou seja $15$, $30$ e $\SI{45}{\volt}$; soma
$=\SI{90}{\volt}$ \checkmark.\\
Como as quedas guardam a mesma proporção $1{:}2{:}3$ dos resistores, poderiam
ter sido obtidas direto: $90$ dividido em $6$ partes de $\SI{15}{\volt}$.}
\end{questao}

\begin{questao}[3]{}
Uma cadeia de $n$ resistores iguais de $\SI{47}{\ohm}$ deve limitar a corrente
a no máximo $\SI{40}{\milli\ampere}$ sob $\SI{24}{\volt}$. Determine o menor
$n$ e a corrente resultante.
\resposta{$n=13$; $I=\SI{39.3}{\milli\ampere}$}
\resolucao{Exige-se $R_{eq}\ge V/I_{\max}=24/0{,}040=\SI{600}{\ohm}$:
\[
47n\ge600 \;\Longrightarrow\; n\ge12{,}77 \;\Longrightarrow\; n=13.
\]
\[
R_{eq}=13\cdot47=\SI{611}{\ohm}
\;\Longrightarrow\;
I=\frac{24}{611}=\SI{39.3}{\milli\ampere}\ \le\ \SI{40}{\milli\ampere}\ \checkmark
\]
\textbf{Cuidado com o arredondamento:} $n=12$ daria
$R_{eq}=\SI{564}{\ohm}$ e $I=\SI{42.6}{\milli\ampere}$ --- \emph{acima} do
limite. Em restrições do tipo "no máximo", arredonde sempre \textbf{para cima},
mesmo que a fração seja pequena.}
\end{questao}

\begin{questao}[3]{}
Em uma cadeia série de dois resistores sob tensão fixa $V$, demonstre que a
soma das potências é $P=\dfrac{V^{2}}{R_1+R_2}$ e explique por que aumentar
\emph{qualquer} resistor reduz a potência total, embora possa \emph{aumentar} a
potência do outro.
\resposta{Aumentar $R_k$ reduz $P_{tot}$; a potência do outro sempre diminui
--- é a do próprio $R_k$ que pode subir}
\resolucao{Com $I=V/(R_1+R_2)$:
\[
P=VI=\frac{V^{2}}{R_1+R_2}.
\]
Como $R_1+R_2$ cresce, $P$ diminui --- \emph{sempre}.\\
Para o outro resistor, $P_2=R_2I^{2}=\dfrac{V^{2}R_2}{(R_1+R_2)^{2}}$: com
$R_2$ fixo e $R_1$ crescendo, $P_2$ só pode \textbf{diminuir}.\\
Já a potência do \emph{próprio} $R_1$,
$P_1=\dfrac{V^{2}R_1}{(R_1+R_2)^{2}}$, cresce enquanto $R_1<R_2$, atinge um
máximo em $R_1=R_2$ e depois decresce. \textbf{Conclusão:} o enunciado
intuitivo "mais resistência, mais aquecimento" é falso em série --- o
aquecimento de um resistor é máximo quando ele iguala a soma dos demais, e cai
para os dois lados. Este resultado é retomado quantitativamente no bloco de
desafio.}
\end{questao}

\blocodesafio

\begin{questao}[4]{}
\textbf{(Projeto.)} Uma cadeia série de três resistores deve dividir
$\SI{24}{\volt}$ em quedas de $\SI{6}{\volt}$, $\SI{8}{\volt}$ e
$\SI{10}{\volt}$, com corrente de $\SI{20}{\milli\ampere}$.
\begin{itensq}
\item Determine os três resistores e suas potências.
\item Escolha valores comerciais E24 e recalcule as quedas reais.
\item Avalie se o erro é aceitável para alimentar cargas com tolerância de
      $\pm5\%$.
\end{itensq}
\resposta{(a) $\SI{300}{\ohm}$, $\SI{400}{\ohm}$, $\SI{500}{\ohm}$;
$\SI{0.12}{\watt}$, $\SI{0.16}{\watt}$, $\SI{0.20}{\watt}$;
(b) $300/390/510\ \si{\ohm}$ $\Rightarrow$ $5{,}98/7{,}77/10{,}16\ \si{\volt}$;
(c) sim, erro máximo de $2{,}9\%$}
\resolucao{(a) Como a corrente é comum, $R_k=U_k/I$:
\[
R_1=\frac{6}{0{,}02}=300,\quad R_2=\frac{8}{0{,}02}=400,\quad
R_3=\frac{10}{0{,}02}=\SI{500}{\ohm}.
\]
Potências: $P_k=U_kI$ $\Rightarrow$ $0{,}12$; $0{,}16$;
$\SI{0.20}{\watt}$. Um resistor de $\frac{1}{2}\,\si{\watt}$ em cada posição
oferece folga confortável.\\
(b) A E24 tem $300$ e $510$, mas não $400$ --- o vizinho é $\SI{390}{\ohm}$.
Com $R_{eq}=300+390+510=\SI{1200}{\ohm}$, a corrente real é
$24/1200=\SI{20.0}{\milli\ampere}$ (coincidência favorável), e as quedas ficam
\[
U_1=\SI{6.00}{\volt},\quad U_2=\SI{7.80}{\volt},\quad U_3=\SI{10.20}{\volt}.
\]
(c) Desvios: $0\%$, $-2{,}5\%$ e $+2{,}0\%$ --- todos dentro de $\pm5\%$.
\textbf{Ressalva essencial:} este projeto só vale \emph{em vazio}. Qualquer
carga real conectada em paralelo com uma das quedas altera a corrente da cadeia
e derruba todas as tensões. Um divisor resistivo puro serve para referências de
altíssima impedância, nunca para alimentar cargas.}
\end{questao}

\begin{questao}[4]{}
\textbf{(Propagação de incerteza.)} Uma cadeia é formada por $N$ resistores
nominalmente iguais, cada um com tolerância relativa $t$.
\begin{itensq}
\item Obtenha a tolerância de $R_{eq}$ no critério de \textbf{pior caso}.
\item Obtenha-a no critério \textbf{estatístico} (erros independentes).
\item Compare para $N=5$ e $t=5\%$, e discuta qual usar.
\end{itensq}
\resposta{(a) $t$; (b) $t/\sqrt{N}$; (c) $5\%$ contra $2{,}24\%$}
\resolucao{(a) Pior caso: todos desviam para o mesmo lado.
$\Delta R_{eq}=N(tR)$ e $R_{eq}=NR$, logo
\[
\frac{\Delta R_{eq}}{R_{eq}}=\frac{NtR}{NR}=t.
\]
A tolerância relativa \textbf{não melhora nem piora} --- resultado que
surpreende, mas decorre de a série somar tanto os valores quanto os desvios.\\
(b) Desvios independentes combinam-se em quadratura:
\[
u(R_{eq})=\sqrt{N}\,(tR)
\;\Longrightarrow\;
\frac{u(R_{eq})}{R_{eq}}=\frac{\sqrt{N}\,tR}{NR}=\frac{t}{\sqrt{N}}.
\]
(c) Para $N=5$, $t=5\%$: pior caso $5\%$; estatístico
$5/\sqrt5=2{,}24\%$.\\
\textbf{Qual usar:} o \emph{pior caso} para projetos de segurança e para lotes
pequenos, em que os desvios podem ser correlacionados (resistores do mesmo
rolo, mesma batelada, mesma temperatura). O \emph{estatístico} para produção em
série com componentes de origens independentes, onde o pior caso seria
excessivamente conservador. Note que a hipótese de independência é quase sempre
otimista na prática.}
\end{questao}

\begin{questao}[4]{}
\textbf{(Projeto com compensação térmica.)} Deseja-se uma associação série cuja
resistência total seja \emph{independente} da temperatura, combinando um
resistor de $R_1=\SI{100}{\ohm}$ com $\alpha_1=+4000\,\mathrm{ppm/^{\circ}C}$ e
outro com $\alpha_2=-1000\,\mathrm{ppm/^{\circ}C}$.
\begin{itensq}
\item Deduza a condição de compensação.
\item Determine $R_2$ e a resistência total.
\item Discuta a limitação prática dessa compensação.
\end{itensq}
\resposta{(a) $R_1\alpha_1+R_2\alpha_2=0$; (b) $R_2=\SI{400}{\ohm}$,
$R_{eq}=\SI{500}{\ohm}$; (c) só é exata em primeira ordem e numa temperatura}
\resolucao{(a) Para pequenas variações,
$R_k(T)=R_k[1+\alpha_k\Delta T]$, logo
\[
R_{eq}(T)=(R_1+R_2)+(R_1\alpha_1+R_2\alpha_2)\Delta T.
\]
Anular o termo em $\Delta T$ exige
\[
R_1\alpha_1+R_2\alpha_2=0.
\]
(b) $R_2=-\dfrac{R_1\alpha_1}{\alpha_2}
=\dfrac{100\cdot4000}{1000}=\SI{400}{\ohm}$, e
$R_{eq}=100+400=\SI{500}{\ohm}$.\\
\textbf{Verificação a $\Delta T=\SI{50}{\celsius}$:}
$R_1\to100(1{,}20)=\SI{120}{\ohm}$ e $R_2\to400(0{,}95)=\SI{380}{\ohm}$; soma
$=\SI{500}{\ohm}$ \checkmark.\\
(c) Três limitações. \textbf{(i)} A compensação é de \emph{primeira ordem};
$\alpha$ real varia com $T$, e o cancelamento se degrada longe da temperatura
de projeto. \textbf{(ii)} Os dois resistores precisam estar \emph{à mesma
temperatura} --- exige acoplamento térmico, o que é difícil quando as
dissipações diferem (aqui, $P_2=4P_1$). \textbf{(iii)} O valor de $R_2$ fica
\emph{amarrado} pela razão $\alpha_1/\alpha_2$: não se pode escolher
livremente $R_{eq}$ \emph{e} exigir compensação. Materiais dedicados
(manganina, constantan, com $\alpha$ naturalmente próximo de zero) resolvem o
problema de forma mais direta.}
\end{questao}

\begin{questao}[4]{}
\textbf{(Projeto sob incerteza.)} Um LED deve ser alimentado por uma fonte
entre $\SI{9}{\volt}$ e $\SI{12}{\volt}$. Sua queda direta varia entre
$\SI{1.8}{\volt}$ e $\SI{2.2}{\volt}$, e a corrente não pode exceder
$\SI{20}{\milli\ampere}$.
\begin{itensq}
\item Identifique a combinação de pior caso e determine o menor $R$ admissível.
\item Determine a corrente mínima resultante e avalie o impacto no brilho.
\item Especifique a potência do resistor.
\end{itensq}
\resposta{(a) $\SI{12}{\volt}$ com $\SI{1.8}{\volt}$ $\Rightarrow$
$R\ge\SI{510}{\ohm}$; (b) $\SI{13.3}{\milli\ampere}$;
(c) $\SI{0.204}{\watt}\Rightarrow\frac{1}{2}\,\si{\watt}$}
\resolucao{(a) A corrente é máxima quando a fonte está no \textbf{máximo} e a
queda do LED no \textbf{mínimo} --- é preciso combinar os extremos que atuam no
mesmo sentido, e não os extremos "nominais":
\[
R\ge\frac{12-1{,}8}{0{,}020}=\SI{510}{\ohm}.
\]
O valor $\SI{510}{\ohm}$ existe na série E24.\\
(b) O outro extremo (fonte mínima, queda máxima):
\[
I_{\min}=\frac{9-2{,}2}{510}=\SI{13.3}{\milli\ampere}.
\]
A corrente varia $13{,}3$--$\SI{20}{\milli\ampere}$, razão de $1{,}5{:}1$. Como
o fluxo luminoso é aproximadamente proporcional à corrente, o brilho varia
cerca de $33\%$ --- perceptível, porém aceitável para sinalização. Para brilho
constante seria preciso uma fonte de \emph{corrente}, não um resistor.\\
(c) $P=RI_{\max}^{2}=510(0{,}020)^{2}=\SI{0.204}{\watt}$; especifique
$\frac{1}{2}\,\si{\watt}$.\\
\textbf{Método:} em projeto com faixas, nunca use valores nominais --- enumere
os extremos e identifique qual combinação é crítica para \emph{cada} requisito.}
\end{questao}

\begin{questao}[4]{}
\textbf{(Sensibilidade.)} Em uma cadeia $I=\dfrac{V}{R_1+R_2}$, com
$V=\SI{12}{\volt}$, $R_1=\SI{100}{\ohm}$ e $R_2=\SI{200}{\ohm}$.
\begin{itensq}
\item Obtenha $\partial I/\partial R_1$ e avalie-a numericamente.
\item Obtenha a sensibilidade \emph{relativa}
      $S=\dfrac{\partial I/I}{\partial R_1/R_1}$.
\item Interprete o resultado para $R_1\ll R_2$ e $R_1\gg R_2$.
\end{itensq}
\resposta{(a) $-\pot{1.33}{-4}\,\si{\ampere\per\ohm}$;
(b) $S=-R_1/(R_1+R_2)=-1/3$; (c) $S\to0$ e $S\to-1$}
\resolucao{(a)
\[
\frac{\partial I}{\partial R_1}=-\frac{V}{(R_1+R_2)^{2}}
=-\frac{12}{300^{2}}=-\pot{1.33}{-4}\,\si{\ampere\per\ohm}.
\]
Um aumento de $\SI{1}{\ohm}$ em $R_1$ reduz a corrente em
$\SI{0.133}{\milli\ampere}$ (de $\SI{40}{\milli\ampere}$).\\
(b)
\[
S=\frac{\partial I}{\partial R_1}\cdot\frac{R_1}{I}
=-\frac{V}{(R_1+R_2)^{2}}\cdot\frac{R_1(R_1+R_2)}{V}
=-\frac{R_1}{R_1+R_2}=-\frac13.
\]
Ou seja: $1\%$ de erro em $R_1$ produz apenas $0{,}33\%$ de erro em $I$.\\
(c) Se $R_1\ll R_2$, $S\to0$: a corrente é governada por $R_2$ e $R_1$ torna-se
irrelevante --- é o caso do resistor \emph{shunt} de medição, deliberadamente
pequeno para não perturbar o circuito. Se $R_1\gg R_2$, $S\to-1$: $R_1$ passa a
determinar sozinho a corrente, e toda a sua imprecisão é transferida
integralmente. \textbf{Regra de projeto:} coloque a tolerância apertada no
elemento cuja sensibilidade relativa se aproxima de $1$; nos demais, use
componentes baratos.}
\end{questao}

\begin{questao}[4]{}
\textbf{(Análise de falhas.)} Dez resistores de $\SI{10}{\ohm}$ formam uma
cadeia sob $\SI{100}{\volt}$.
\begin{itensq}
\item Determine corrente, queda e potência por resistor em operação normal.
\item Repita supondo que \emph{um} resistor entre em curto.
\item Compare com o caso de um resistor \emph{aberto} e discuta qual falha é
      mais perigosa.
\end{itensq}
\resposta{(a) $\SI{1}{\ampere}$, $\SI{10}{\volt}$, $\SI{10}{\watt}$;
(b) $\SI{1.11}{\ampere}$, $\SI{11.1}{\volt}$, $\SI{12.3}{\watt}$;
(c) o curto é mais perigoso}
\resolucao{(a) $R_{eq}=\SI{100}{\ohm}$, $I=\SI{1}{\ampere}$,
$U=\SI{10}{\volt}$, $P=\SI{10}{\watt}$ cada; total $\SI{100}{\watt}$.\\
(b) $R_{eq}=\SI{90}{\ohm}$, $I=100/90=\SI{1.111}{\ampere}$ ($+11\%$),
$U=\SI{11.1}{\volt}$, $P=10(1{,}111)^{2}=\SI{12.35}{\watt}$ ($+23{,}5\%$);
total $\SI{111}{\watt}$.\\
(c) O \textbf{aberto} zera a corrente: o circuito para de funcionar, mas nada
sofre --- é uma falha \emph{segura} (\emph{fail-safe}), fácil de localizar com
voltímetro.\\
O \textbf{curto} é insidioso: o circuito continua operando, aparentemente
normal, mas cada resistor remanescente passa a dissipar $23{,}5\%$ a mais. Se
já operavam próximos do nominal, entram em sobrecarga --- e cada novo curto
agrava o quadro de forma \emph{acelerada}. Com dois curtos,
$P=\SI{15.6}{\watt}$ ($+56\%$); com três, $\SI{20.4}{\watt}$ ($+104\%$).
\textbf{Conclusão:} falhas que \emph{reduzem} a resistência de um circuito em
série tendem a ser progressivas e devem ser detectadas por monitoração de
corrente, não por observação funcional.}
\end{questao}

\begin{questao}[4]{}
\textbf{(Método experimental.)} Para determinar $\EMF$ e $r$ de uma bateria,
mediu-se a tensão de terminal em duas condições de carga:
$(\SI{1}{\ampere};\ \SI{11.5}{\volt})$ e
$(\SI{3}{\ampere};\ \SI{10.5}{\volt})$.
\begin{itensq}
\item Obtenha $r$ e $\EMF$.
\item Explique a vantagem de usar dois pontos em vez de medir a tensão em
      vazio.
\item Estime a corrente de curto-circuito e comente sua validade.
\end{itensq}
\resposta{(a) $r=\SI{0.5}{\ohm}$, $\EMF=\SI{12}{\volt}$;
(c) $I_{cc}=\SI{24}{\ampere}$, valor apenas indicativo}
\resolucao{(a) O modelo é $V=\EMF-rI$ --- uma reta. Dois pontos a determinam:
\[
r=-\frac{\Delta V}{\Delta I}=-\frac{10{,}5-11{,}5}{3-1}=\SI{0.5}{\ohm},
\]
\[
\EMF=V+rI=11{,}5+0{,}5(1)=\SI{12}{\volt}.
\]
(b) A medida em vazio exigiria um voltímetro de impedância \emph{infinita}; com
impedância finita, circula uma pequena corrente e a leitura já não é $\EMF$.
Além disso, baterias reais apresentam recuperação de tensão em repouso, o que
falseia a leitura em vazio. O ajuste por dois (ou mais) pontos de carga
\textbf{extrapola} para $I=0$ sem precisar chegar lá, e ainda entrega $r$ de
brinde. Com três ou mais pontos, uma regressão linear também revela se o modelo
é adequado.\\
(c) $I_{cc}=\EMF/r=12/0{,}5=\SI{24}{\ampere}$. \textbf{Ressalva:} o valor
pressupõe $r$ constante, hipótese que falha em correntes altas --- $r$ cresce
com a corrente e com a descarga, e há aquecimento. O número serve como
\emph{ordem de grandeza} para dimensionar proteção, jamais como previsão
precisa.}
\end{questao}

\begin{questao}[4]{}
\textbf{(Projeto com componentes disponíveis.)} Deseja-se uma resistência de
$\SI{1}{\kilo\ohm}$ capaz de dissipar $\SI{2}{\watt}$, dispondo apenas de
resistores de $\SI{200}{\ohm}$ e $\frac{1}{2}\,\si{\watt}$.
\begin{itensq}
\item Proponha a associação e verifique o valor.
\item Determine a potência de cada resistor e a margem de segurança.
\item Calcule a tensão e a corrente nominais da associação.
\end{itensq}
\resposta{(a) cinco em série; (b) $\SI{0.4}{\watt}$ cada, margem de $25\%$;
(c) $\SI{44.7}{\volt}$ e $\SI{44.7}{\milli\ampere}$}
\resolucao{(a) $5\times\SI{200}{\ohm}=\SI{1000}{\ohm}$ \checkmark.\\
(b) Em série a corrente é comum e os resistores são iguais, logo a potência se
reparte igualmente:
\[
P_k=\frac{2}{5}=\SI{0.4}{\watt}\ <\ \SI{0.5}{\watt}\ \checkmark
\]
Margem: $(0{,}5-0{,}4)/0{,}4=25\%$. É uma folga modesta --- em ambiente quente
ou sem ventilação, considere sete unidades de $\SI{150}{\ohm}$
($\SI{1050}{\ohm}$) ou resistores de $\SI{1}{\watt}$.\\
(c)
\[
V=\sqrt{PR}=\sqrt{2\cdot1000}=\SI{44.7}{\volt},
\qquad
I=\sqrt{P/R}=\sqrt{0{,}002}=\SI{44.7}{\milli\ampere}.
\]
\textbf{Por que a série resolve:} ela multiplica \emph{simultaneamente} o valor
e a capacidade de dissipação por $N$. Uma associação em paralelo de cinco
unidades também daria $\SI{2.5}{\watt}$, mas $\SI{40}{\ohm}$ --- resolveria a
potência e destruiria o valor.}
\end{questao}

\begin{questao}[4]{}
\textbf{(Instrumentação.)} Cinco sensores resistivos de $\SI{1}{\kilo\ohm}$
estão em série, alimentados por uma fonte de \textbf{corrente constante} de
$\SI{1}{\milli\ampere}$. Mede-se a tensão de cada nó em relação ao terra.
\begin{itensq}
\item Determine as cinco tensões nodais em condição nominal.
\item Um sensor passa a $\SI{1.2}{\kilo\ohm}$. Mostre como identificar
      \emph{qual} deles mudou.
\item Explique por que a fonte de corrente é preferível a uma fonte de tensão.
\end{itensq}
\resposta{(a) $1$, $2$, $3$, $4$ e $\SI{5}{\volt}$; (b) pela quebra do degrau
uniforme de $\SI{1}{\volt}$; (c) desacopla os sensores entre si}
\resolucao{(a) Com $I=\SI{1}{\milli\ampere}$, cada sensor cai
$\SI{1}{\volt}$. Acumulando a partir do terra: $1$, $2$, $3$, $4$,
$\SI{5}{\volt}$.\\
(b) Se o \emph{terceiro} sensor for a $\SI{1.2}{\kilo\ohm}$, as tensões passam
a $1$, $2$, $3{,}2$, $4{,}2$, $\SI{5.2}{\volt}$. A \textbf{diferença entre nós
consecutivos} vale $1;\,1;\,1{,}2;\,1;\,1$ --- o degrau anômalo aponta
diretamente o sensor defeituoso, e sua magnitude dá o novo valor:
$\Delta R=\Delta U/I=0{,}2/0{,}001=\SI{200}{\ohm}$.\\
(c) Com \textbf{corrente constante}, a queda de cada sensor é $U_k=IR_k$ ---
depende \emph{só} do próprio sensor. A variação de um não afeta a leitura dos
demais, e a relação é rigorosamente linear.\\
Com fonte de \textbf{tensão}, a corrente seria $V/\sum R_k$: a mudança de um
sensor alteraria a corrente e, portanto, \emph{todas} as leituras
simultaneamente --- o sistema ficaria acoplado, com resposta não linear e sem
identificação direta do sensor. É exatamente por isso que sistemas de RTD e
extensômetros são excitados por corrente.}
\end{questao}

\begin{questao}[4]{}
\textbf{(Otimização.)} Uma cadeia série sob tensão fixa $V$ contém um resistor
ajustável $R$ e um bloco fixo de resistência total $R_f$.
\begin{itensq}
\item Obtenha $P_R(R)$ e mostre que ela é máxima em $R=R_f$.
\item Calcule $P_{R,\max}$ e o rendimento nesse ponto.
\item Compare com a condição de máxima potência \emph{total} e explique a
      diferença.
\end{itensq}
\resposta{(a) $R=R_f$; (b) $P_{\max}=V^{2}/(4R_f)$, rendimento $50\%$;
(c) a potência total é máxima em $R\to0$}
\resolucao{(a)
\[
P_R=RI^{2}=\frac{V^{2}R}{(R+R_f)^{2}}.
\]
\[
\frac{\mathrm{d}P_R}{\mathrm{d}R}
=V^{2}\frac{(R+R_f)^{2}-R\cdot2(R+R_f)}{(R+R_f)^{4}}
=V^{2}\frac{R_f-R}{(R+R_f)^{3}}=0
\;\Longrightarrow\;
R=R_f.
\]
A derivada é positiva para $R<R_f$ e negativa para $R>R_f$: máximo confirmado.\\
(b) $P_{\max}=\dfrac{V^{2}R_f}{(2R_f)^{2}}=\dfrac{V^{2}}{4R_f}$. A potência
total nesse ponto é $\dfrac{V^{2}}{2R_f}$, de modo que $R$ recebe exatamente
\textbf{metade} --- rendimento de $50\%$.\\
(c) A potência \emph{total} $P=\dfrac{V^{2}}{R+R_f}$ é monotonicamente
decrescente: ela é máxima quando $R\to0$, e aí $P\to V^{2}/R_f$ --- o dobro do
valor anterior, porém integralmente entregue ao bloco fixo.\\
\textbf{Distinção que importa:} "máxima potência \emph{no elemento}" e "máxima
potência \emph{do circuito}" são objetivos diferentes e incompatíveis. O
primeiro é o critério de casamento (relevante quando o bloco fixo é a
resistência interna de uma fonte); o segundo é o critério de eficiência. Este
mesmo antagonismo reaparece no teorema da máxima transferência de potência.}
\end{questao}

% END BANCO COMPLEMENTAR

\gabarito[3.1 Circuitos em série]

% =====================================================================
\subsection{Circuitos em paralelo}
% =====================================================================

\blocofixacao

\begin{questao}[1]{}
Sobre uma associação de resistores \textbf{em paralelo}, é correto afirmar que:
\begin{alternativas}
\item a corrente é a mesma em todos os ramos.
\item a resistência equivalente é sempre maior que a maior das resistências.
\item a tensão é a mesma em todos os ramos.
\item a tensão se divide proporcionalmente às resistências.
\item a resistência equivalente é a soma das resistências.
\end{alternativas}
\resposta{(c)}
\resolucao{Todos os ramos estão ligados ao mesmo par de nós, logo compartilham a
mesma tensão. A corrente, essa sim, se divide --- e de forma \emph{inversamente}
proporcional às resistências. A resistência equivalente do paralelo é sempre
\textbf{menor} que a menor das resistências associadas, o que elimina (b) e (e).}
\end{questao}

\begin{questao}[1]{}
Para cada circuito, calcule a resistência equivalente, as correntes de ramo e a
corrente total $I_t$.

\circuitoq{%
\subcirc{a}{%
\begin{circuitikz}[scale=0.72,transform shape,line width=0.8pt]
 \draw (0,0) to[battery1,l_={$E=\SI{30}{\volt}$}] (0,3.4);
 \draw (0,3.4) -- (9,3.4);  \draw (0,0) -- (9,0);
 \draw (2.2,3.4) to[R,l=$\SI{4}{\ohm}$,*-*] (2.2,0);
 \draw (4.5,3.4) to[R,l=$\SI{3}{\ohm}$,*-*] (4.5,0);
 \draw (6.8,3.4) to[R,l=$\SI{6}{\ohm}$,*-*] (6.8,0);
 \draw (9,3.4)   to[R,l=$\SI{12}{\ohm}$] (9,0);
 \draw[->,Icol,line width=1pt] (0.5,3.75) -- (1.6,3.75)
       node[midway,above,font=\footnotesize]{$I_t$};
\end{circuitikz}}
\hfill
\subcirc{b}{%
\begin{circuitikz}[scale=0.72,transform shape,line width=0.8pt]
 \draw (0,0) to[battery1,l_={$E=\SI{24}{\volt}$}] (0,3.4);
 \draw (0,3.4) -- (6.6,3.4);  \draw (0,0) -- (6.6,0);
 \draw (2.2,3.4) to[R,l=$\SI{12}{\ohm}$,*-*] (2.2,0);
 \draw (4.4,3.4) to[R,l=$\SI{12}{\ohm}$,*-*] (4.4,0);
 \draw (6.6,3.4) to[R,l=$\SI{12}{\ohm}$] (6.6,0);
 \draw[->,Icol,line width=1pt] (0.5,3.75) -- (1.6,3.75)
       node[midway,above,font=\footnotesize]{$I_t$};
\end{circuitikz}}

\vspace{1.1em}

\subcirc{c}{%
\begin{circuitikz}[scale=0.72,transform shape,line width=0.8pt]
 \draw (0,0) to[battery1,l_={$E=\SI{15}{\volt}$}] (0,3.4);
 \draw (0,3.4) -- (6.6,3.4);  \draw (0,0) -- (6.6,0);
 \draw (2.2,3.4) to[R,l=$\SI{12}{\ohm}$,*-*] (2.2,0);
 \draw (4.4,3.4) to[R,l=$\SI{12}{\ohm}$,*-*] (4.4,0);
 \draw (6.6,3.4) to[R,l=$\SI{3}{\ohm}$] (6.6,0);
 \draw[->,Icol,line width=1pt] (0.5,3.75) -- (1.6,3.75)
       node[midway,above,font=\footnotesize]{$I_t$};
\end{circuitikz}}
\hfill
\subcirc{d}{%
\begin{circuitikz}[scale=0.72,transform shape,line width=0.8pt]
 \draw (0,0) to[battery1,l_={$E=\SI{9}{\volt}$}] (0,3.4);
 \draw (0,3.4) -- (4.4,3.4);  \draw (0,0) -- (4.4,0);
 \draw (2.2,3.4) to[R,l=$\SI{30}{\ohm}$,*-*] (2.2,0);
 \draw (4.4,3.4) to[R,l=$\SI{60}{\ohm}$] (4.4,0);
 \draw[->,Icol,line width=1pt] (0.5,3.75) -- (1.6,3.75)
       node[midway,above,font=\footnotesize]{$I_t$};
\end{circuitikz}}}

\resposta{(a) \SI{1.2}{\ohm}, $I_t=\SI{25}{\ampere}$;
(b) \SI{4}{\ohm}, $I_t=\SI{6}{\ampere}$;
(c) \SI{2}{\ohm}, $I_t=\SI{7.5}{\ampere}$;
(d) \SI{20}{\ohm}, $I_t=\SI{0.45}{\ampere}$}
\resolucao{Em todos os casos a tensão da fonte aparece \emph{inteira} sobre cada
ramo, então cada corrente sai direto da lei de Ohm.\\
\textbf{(a)} $I_1=30/4=\SI{7.5}{\ampere}$, $I_2=30/3=\SI{10}{\ampere}$,
$I_3=30/6=\SI{5}{\ampere}$, $I_4=30/12=\SI{2.5}{\ampere}$;
$I_t=\SI{25}{\ampere}$ e $\Req=30/25=\SI{1.2}{\ohm}$.\\
\textbf{(b)} Três resistores iguais: $\Req = 12/3 = \SI{4}{\ohm}$;
cada ramo com $24/12=\SI{2}{\ampere}$ e $I_t = \SI{6}{\ampere}$.\\
\textbf{(c)} $\dfrac{1}{\Req}=\dfrac{1}{12}+\dfrac{1}{12}+\dfrac{1}{3}
=\dfrac{1+1+4}{12}=\dfrac{6}{12}$, logo $\Req=\SI{2}{\ohm}$.
$I_1=I_2=\SI{1.25}{\ampere}$, $I_3=\SI{5}{\ampere}$, $I_t=\SI{7.5}{\ampere}$.\\
\textbf{(d)} $\Req=\dfrac{30\cdot 60}{30+60}=\SI{20}{\ohm}$;
$I_1=\SI{0.3}{\ampere}$, $I_2=\SI{0.15}{\ampere}$, $I_t=\SI{0.45}{\ampere}$.\\
Note em (a) que $\Req$ ficou menor que o menor resistor (\SI{3}{\ohm}) --- sempre
é assim no paralelo.}
\end{questao}

\begin{questao}[1]{}
$n$ resistores iguais de $\SI{60}{\ohm}$ são associados em paralelo, resultando
em $\SI{10}{\ohm}$. O valor de $n$ é:
\begin{alternativas}[3]
\item 2
\item 3
\item 4
\item 6
\item 10
\end{alternativas}
\resposta{(d)}
\resolucao{Para $n$ resistores iguais em paralelo, $\Req = \dfrac{R}{n}$, logo
$n = \dfrac{60}{10} = 6$.}
\end{questao}

\blocoaplicacao

\begin{questao}[2]{}
Um resistor de $\SI{20}{\ohm}$ é associado em paralelo com um resistor $R$
desconhecido, resultando em uma resistência equivalente de $\SI{12}{\ohm}$.
Determine $R$.
\resposta{$R = \SI{30}{\ohm}$}
\resolucao{Da definição de paralelo,
\[
\frac{1}{12}=\frac{1}{20}+\frac{1}{R}
\quad\Rightarrow\quad
\frac{1}{R}=\frac{1}{12}-\frac{1}{20}=\frac{5-3}{60}=\frac{2}{60}
\quad\Rightarrow\quad R=\SI{30}{\ohm}.
\]
\textbf{Atalho útil:} com dois resistores,
$R = \dfrac{R_1\,\Req}{R_1-\Req} = \dfrac{20\cdot 12}{20-12} = \SI{30}{\ohm}$.
Note que $\Req$ precisa necessariamente ser menor que \SI{20}{\ohm}; se o problema
pedisse $\Req = \SI{25}{\ohm}$, não haveria solução com resistor positivo.}
\end{questao}

\begin{questao}[2]{}
Uma corrente total de $\SI{12}{\ampere}$ chega a um nó e se divide entre dois
resistores em paralelo, de $\SI{4}{\ohm}$ e $\SI{12}{\ohm}$. As correntes em cada
resistor valem, respectivamente:
\begin{alternativas}[2]
\item \SI{3}{\ampere} e \SI{9}{\ampere}
\item \SI{9}{\ampere} e \SI{3}{\ampere}
\item \SI{6}{\ampere} e \SI{6}{\ampere}
\item \SI{4}{\ampere} e \SI{8}{\ampere}
\item \SI{8}{\ampere} e \SI{4}{\ampere}
\end{alternativas}
\resposta{(b)}
\resolucao{A corrente se divide de forma \emph{inversamente} proporcional às
resistências: o menor resistor recebe a maior corrente. Pelo divisor de corrente,
\[
I_1 = I_t\,\frac{R_2}{R_1+R_2}=12\cdot\frac{12}{16}=\SI{9}{\ampere},
\qquad
I_2 = 12-9 = \SI{3}{\ampere}.
\]
\textbf{Verificação alternativa:} $\Req = \dfrac{4\cdot12}{16}=\SI{3}{\ohm}$ e
$U = 3\cdot 12 = \SI{36}{\volt}$; daí $I_1 = 36/4 = \SI{9}{\ampere}$ e
$I_2 = 36/12 = \SI{3}{\ampere}$. A armadilha da alternativa (a) é aplicar a
proporção direta, como se fosse divisor de tensão.}
\end{questao}

\begin{questao}[2]{}
Três lâmpadas incandescentes de $\SI{60}{\watt}$/$\SI{120}{\volt}$ são ligadas
\textbf{em paralelo} à rede de $\SI{120}{\volt}$.
\begin{itensq}
\item Calcule a resistência de cada lâmpada e a resistência equivalente.
\item Determine a corrente total fornecida pela rede.
\item Se uma das lâmpadas queimar, o que acontece com o brilho das outras duas?
\end{itensq}
\resposta{(a) \SI{240}{\ohm} cada, $\Req=\SI{80}{\ohm}$;
(b) \SI{1.5}{\ampere}; (c) nada muda}
\resolucao{(a) De $P = \dfrac{U^2}{R}$:
$R = \dfrac{120^2}{60} = \SI{240}{\ohm}$ por lâmpada.
Três iguais em paralelo: $\Req = \dfrac{240}{3} = \SI{80}{\ohm}$.\\
(b) $I_t = \dfrac{120}{80} = \SI{1.5}{\ampere}$
(ou $3\times\SI{0.5}{\ampere}$ por lâmpada).\\
(c) As demais continuam com os mesmos \SI{120}{\volt} nos seus terminais, logo o
brilho \emph{não muda}: cada ramo é independente. É por isso que as instalações
residenciais são feitas em paralelo, e não em série.}
\end{questao}

\begin{questao}[2]{}
No trecho da figura, um fio ideal (sem resistência) é ligado em paralelo com o
resistor $R_2$.

\circuitoq{%
\begin{circuitikz}[scale=0.95,transform shape,line width=0.8pt]
 \draw (0,0) to[battery1,l_={$E=\SI{20}{\volt}$}] (0,3)
       to[R,l={$R_1=\SI{10}{\ohm}$}] (3,3)
       to[R,l={$R_2=\SI{20}{\ohm}$}] (6,3) -- (6,0) -- (0,0);
 \draw[Vcol,line width=1.1pt] (3,3) -- (3,4.2) -- (6,4.2) -- (6,3);
 \node[Vcol,font=\footnotesize\sffamily,above] at (4.5,4.25) {fio ideal};
 \fill (3,3) circle (0.06); \fill (6,3) circle (0.06);
\end{circuitikz}}

A corrente fornecida pela fonte vale:
\begin{alternativas}[3]
\item \SI{0.67}{\ampere}
\item \SI{1}{\ampere}
\item \SI{1.5}{\ampere}
\item \SI{2}{\ampere}
\item \SI{3}{\ampere}
\end{alternativas}
\resposta{(d)}
\resolucao{O fio ideal curto-circuita $R_2$: os dois terminais de $R_2$ ficam no
\emph{mesmo potencial}, logo $U_{R_2} = 0$ e nenhuma corrente passa por ele.
Formalmente, $R_2 \parallel 0 = 0$. Resta apenas $R_1$:
\[
I=\frac{20}{10}=\SI{2}{\ampere}.
\]
\textbf{Erro comum:} calcular $10 + 20 = \SI{30}{\ohm}$ e responder
\SI{0.67}{\ampere}, ignorando o curto.}
\end{questao}

\blocoaprofundamento

\begin{questao}[3]{}
Um resistor de $\SI{100}{\ohm}$ suporta no máximo $\SI{2}{\watt}$ de dissipação.
Deseja-se construir, com resistores desse mesmo tipo, uma associação
\textbf{em paralelo} de $\SI{25}{\ohm}$.
\begin{itensq}
\item Quantos resistores são necessários?
\item Qual é a máxima tensão que pode ser aplicada à associação?
\item Qual é a potência máxima que a associação suporta?
\end{itensq}
\resposta{(a) 4 resistores; (b) $\approx\SI{14.1}{\volt}$;
(c) \SI{8}{\watt}}
\resolucao{(a) $\Req = R/n \Rightarrow n = 100/25 = 4$ resistores.\\
(b) Todos ficam sob a \emph{mesma} tensão, então basta respeitar o limite de um
deles:
\[
P_{\max}=\frac{U^2}{R}\ \Rightarrow\ U_{\max}=\sqrt{P_{\max}R}=\sqrt{2\cdot100}
=\sqrt{200}\approx\SI{14.1}{\volt}.
\]
(c) Como os 4 resistores trabalham no limite simultaneamente,
$P_{\text{total}} = 4\cdot 2 = \SI{8}{\watt}$ --- que também pode ser conferido
por $P = U^2/\Req = 200/25 = \SI{8}{\watt}$.\\
\textbf{Ideia geral:} associar resistores iguais em paralelo (ou em série)
multiplica por $n$ a capacidade de dissipação do conjunto.}
\end{questao}

\begin{questao}[3]{}
Determine a resistência equivalente entre os terminais A e B do circuito da
figura, no qual dois trechos são interligados por fios ideais.

\circuitoq{%
\begin{circuitikz}[scale=0.85,transform shape,line width=0.8pt]
 \draw (0,0) node[left,font=\bfseries] {A}
       to[R,l=$\SI{15}{\ohm}$,o-*] (3,0)
       to[R,l=$\SI{30}{\ohm}$,-*] (6,0)
       to[R,l=$\SI{30}{\ohm}$,-o] (9,0)
       node[right,font=\bfseries] {B};
 \draw[Vcol,line width=1.1pt] (3,0) -- (3,1.8) -- (6,1.8) -- (6,0);
 \node[Vcol,font=\footnotesize\sffamily,above] at (4.5,1.85) {fio ideal};
\end{circuitikz}}

\resposta{$R_{AB} = \SI{45}{\ohm}$}
\resolucao{O fio ideal liga os dois terminais do resistor de \SI{30}{\ohm}
central, curto-circuitando-o. Ele deixa de conduzir e pode ser desconsiderado no circuito equivalente. Sobram
apenas os resistores de \SI{15}{\ohm} e \SI{30}{\ohm} em série:
\[
R_{AB}=15+30=\SI{45}{\ohm}.
\]
\textbf{Método seguro:} antes de calcular, redesenhe o circuito nomeando os nós.
Dois pontos ligados por fio ideal são \emph{o mesmo nó} --- qualquer componente
entre eles está curto-circuitado.}
\end{questao}

\blocodesafio

\begin{questao}[4]{}
$n$ resistores \emph{diferentes} $R_1, R_2, \dots, R_n$ são associados em
paralelo.
\begin{itensq}
\item Demonstre que a resistência equivalente é sempre \textbf{menor} que a menor
      das resistências.
\item Mostre que, se um dos resistores tende a zero, $\Req \to 0$, e interprete
      fisicamente.
\item Dois resistores em paralelo dão $\SI{6}{\ohm}$; em série, dariam
      $\SI{25}{\ohm}$. Determine seus valores.
\end{itensq}
\resposta{(a) e (b) demonstrações; (c) $\SI{10}{\ohm}$ e $\SI{15}{\ohm}$}
\resolucao{\textbf{(a)} Seja $R_{\min}$ a menor das resistências. Da definição,
\[
\frac{1}{\Req}=\frac{1}{R_1}+\dots+\frac{1}{R_n}
> \frac{1}{R_{\min}},
\]
pois a soma inclui o termo $\dfrac{1}{R_{\min}}$ \emph{mais} outras parcelas
positivas. Invertendo a desigualdade (ambos os lados positivos):
$\Req < R_{\min}$. \checkmark\\
Interpretação: acrescentar um caminho \emph{sempre} facilita a passagem de
corrente.

\textbf{(b)} Se $R_k \to 0$, então $\dfrac{1}{R_k}\to\infty$, e a soma das
condutâncias diverge, levando $\Req \to 0$. Fisicamente, esse ramo é um
\textbf{curto-circuito}: ele desvia toda a corrente e anula a tensão sobre o
conjunto, tornando irrelevantes os demais resistores.

\textbf{(c)} Temos o sistema
\[
\begin{cases}
R_1+R_2 = 25\\[2pt]
\dfrac{R_1R_2}{R_1+R_2}=6 \;\Rightarrow\; R_1R_2 = 6\cdot25 = 150.
\end{cases}
\]
Soma e produto conhecidos: $R_1$ e $R_2$ são as raízes de
$x^2-25x+150=0$, isto é,
\[
x=\frac{25\pm\sqrt{625-600}}{2}=\frac{25\pm5}{2}
\;\Longrightarrow\;
R_1=\SI{15}{\ohm},\quad R_2=\SI{10}{\ohm}.
\]
\textbf{Verificação:} $15+10=25$ \checkmark\ e
$\dfrac{15\cdot10}{25}=\SI{6}{\ohm}$ \checkmark\\
\textbf{Observação:} o problema só tem solução real se
$(R_s)^2 \ge 4R_sR_p$, ou seja, $R_s \ge 4R_p$. Aqui, $25 \ge 24$ --- a condição é satisfeita por pequena margem. Se o enunciado pedisse $\Req=\SI{7}{\ohm}$ em paralelo com a mesma soma
$\SI{25}{\ohm}$, não existiriam resistores reais que satisfizessem.}
\end{questao}

% BEGIN BANCO COMPLEMENTAR Circuitos em paralelo
% =====================================================================
%  Banco complementar --- Circuitos Paralelo  (REVISADO)
%  Substitui a versao gerada por template (11 clones em N2, 13 em N3,
%  9 questoes conceituais indevidamente marcadas como N4).
%  Sem sobreposicao com o cap03 (n resistores iguais, R desconhecido a
%  partir de Req, divisor 12 A entre 4 e 12 ohm, tres lampadas de 60 W,
%  fio ideal em paralelo, banco de 100 ohm/2 W e a demonstracao de que
%  Req < menor resistor ja estao la).
%  Distribuicao mantida: 7 N1 + 11 N2 + 13 N3 + 9 N4 = 40
% =====================================================================

\subsubsection*{Banco complementar --- Circuitos paralelo}

\begin{atencao}[Organização do banco]
No paralelo, a \textbf{tensão é comum} e as \textbf{correntes se somam}. Daí
decorre que corrente e potência se distribuem \emph{inversamente} a $R$ ---
exatamente o oposto da série. Trabalhe com \textbf{condutâncias}: elas somam
diretamente e tornam quase toda conta trivial. O N3 explora falhas de ramo,
fuga, deriva térmica e instrumentação; o N4 trata de projeto de \emph{shunt},
sensibilidade, propagação de incerteza e seletividade de proteção.
\end{atencao}

\blocofixacao

\begin{questao}[1]{}
Determine a resistência equivalente de $R_1=\SI{30}{\ohm}$ e
$R_2=\SI{60}{\ohm}$ em paralelo.
\resposta{$R_{eq}=\SI{20}{\ohm}$}
\resolucao{Para dois resistores, a regra do produto pela soma é a mais rápida:
\[
R_{eq}=\frac{R_1R_2}{R_1+R_2}=\frac{30\cdot60}{90}=\SI{20}{\ohm}.
\]
\textbf{Conferência instantânea:} o resultado ($\SI{20}{\ohm}$) é menor que o
menor dos dois ($\SI{30}{\ohm}$) --- se não for, há erro. Cuidado: a regra
produto/soma vale \emph{apenas} para dois resistores.}
\end{questao}

\begin{questao}[1]{}
Três resistores de $\SI{4}{\ohm}$, $\SI{8}{\ohm}$ e $\SI{8}{\ohm}$ estão em
paralelo. Calcule a condutância total e $R_{eq}$.
\resposta{$G=\SI{0.5}{\siemens}$; $R_{eq}=\SI{2}{\ohm}$}
\resolucao{
\[
G=\frac14+\frac18+\frac18=0{,}25+0{,}125+0{,}125=\SI{0.5}{\siemens}
\;\Longrightarrow\;
R_{eq}=\frac{1}{0{,}5}=\SI{2}{\ohm}.
\]
Com três ou mais ramos, some as condutâncias e inverta \emph{uma única vez} no
final. Inverter parcialmente a cada passo é a origem mais comum de erro.}
\end{questao}

\begin{questao}[1]{}
Três resistores iguais de $\SI{90}{\ohm}$ são associados em paralelo. Determine
$R_{eq}$.
\resposta{$R_{eq}=\SI{30}{\ohm}$}
\resolucao{Para $n$ resistores \emph{iguais},
\[
R_{eq}=\frac{R}{n}=\frac{90}{3}=\SI{30}{\ohm}.
\]
O atalho vale só quando todos são iguais; com valores diferentes é obrigatório
somar condutâncias.}
\end{questao}

\begin{questao}[1]{}
Acrescentando-se mais um ramo \textbf{em paralelo} a um circuito alimentado por
fonte ideal de tensão, a corrente total:
\begin{alternativas}[2]
\item aumenta
\item diminui
\item não se altera
\item pode aumentar ou diminuir, conforme o valor
\end{alternativas}
\resposta{(a)}
\resolucao{Cada novo ramo acrescenta uma condutância positiva:
$G_{eq}=\sum G_k$ só pode crescer, logo $R_{eq}$ só pode diminuir. Com $V$
fixo, $I=V/R_{eq}$ aumenta.\\
Note o contraste com a série, em que acrescentar elemento \emph{reduz} a
corrente. E note também que os ramos já existentes continuam com \emph{a mesma}
corrente --- quem cresce é o total.}
\end{questao}

\begin{questao}[1]{}
Um ramo de $\SI{48}{\ohm}$ pertence a uma associação paralela ligada a
$\SI{24}{\volt}$. Determine a corrente desse ramo.
\resposta{$I=\SI{0.5}{\ampere}$}
\resolucao{
\[
I=\frac{V}{R}=\frac{24}{48}=\SI{0.5}{\ampere}.
\]
No paralelo, a tensão é \emph{comum} a todos os ramos: a corrente de cada um
depende só do próprio resistor. Não é preciso conhecer os demais ramos nem
$R_{eq}$.}
\end{questao}

\begin{questao}[1]{}
Uma associação paralela sob $\SI{40}{\volt}$ tem ramos que conduzem
$\SI{2}{\ampere}$, $\SI{3}{\ampere}$ e $\SI{5}{\ampere}$. Determine a corrente
total e $R_{eq}$.
\resposta{$I_t=\SI{10}{\ampere}$; $R_{eq}=\SI{4}{\ohm}$}
\resolucao{Pela LKC, as correntes de ramo se somam:
$I_t=2+3+5=\SI{10}{\ampere}$.
\[
R_{eq}=\frac{V}{I_t}=\frac{40}{10}=\SI{4}{\ohm}.
\]
Observe que não foi preciso conhecer os resistores individualmente --- embora
eles saiam de imediato: $20$, $13{,}3$ e $\SI{8}{\ohm}$.}
\end{questao}

\begin{questao}[1]{}
Em uma associação paralela, o ramo que conduz a \textbf{maior} corrente é:
\begin{alternativas}[2]
\item o de maior resistência
\item o de menor resistência
\item o mais próximo da fonte
\item todos conduzem igualmente
\end{alternativas}
\resposta{(b)}
\resolucao{Com $V$ comum, $I_k=V/R_k$: a corrente é \emph{inversamente}
proporcional à resistência. O ramo de menor $R$ é o caminho de menor oposição e
leva a maior corrente.\\
A alternativa (c) é sedutora, mas falsa: em circuitos de parâmetros
concentrados, a posição física dos ramos é irrelevante --- só importa a
topologia.}
\end{questao}

\blocoaplicacao

\begin{questao}[2]{}
$R_1=\SI{10}{\ohm}$, $R_2=\SI{20}{\ohm}$ e $R_3=\SI{20}{\ohm}$ estão em
paralelo sob $\SI{60}{\volt}$. Determine as correntes de ramo, a corrente total
e $R_{eq}$.
\resposta{$\SI{6}{\ampere}$, $\SI{3}{\ampere}$, $\SI{3}{\ampere}$;
$I_t=\SI{12}{\ampere}$; $R_{eq}=\SI{5}{\ohm}$}
\resolucao{Tensão comum, então cada ramo é independente:
\[
I_1=\frac{60}{10}=6,\quad I_2=I_3=\frac{60}{20}=\SI{3}{\ampere}.
\]
\[
I_t=6+3+3=\SI{12}{\ampere}
\;\Longrightarrow\;
R_{eq}=\frac{60}{12}=\SI{5}{\ohm}.
\]
\textbf{Conferência:} $G=0{,}1+0{,}05+0{,}05=\SI{0.2}{\siemens}$ e
$1/0{,}2=\SI{5}{\ohm}$ \checkmark.}
\end{questao}

\begin{questao}[2]{}
No circuito da questão anterior, calcule a potência de cada ramo e confronte
com a potência total.
\resposta{$\SI{360}{\watt}$, $\SI{180}{\watt}$, $\SI{180}{\watt}$; total
$\SI{720}{\watt}$}
\resolucao{Com $V$ comum, use $P=V^{2}/R$:
\[
P_1=\frac{3600}{10}=360,\qquad P_2=P_3=\frac{3600}{20}=\SI{180}{\watt}.
\]
Fonte: $P=VI_t=60\cdot12=\SI{720}{\watt}=360+180+180$ \checkmark.\\
\textbf{Regra do paralelo:} $P\propto1/R$ --- o \emph{menor} resistor dissipa
mais. Compare com a série, em que $P\propto R$. Trocar essas duas regras é o
erro conceitual mais frequente em associações.}
\end{questao}

\begin{questao}[2]{}
Uma corrente de $\SI{15}{\ampere}$ chega a um nó e se divide entre
$R_1=\SI{8}{\ohm}$ e $R_2=\SI{12}{\ohm}$. Determine as correntes de ramo pela
fórmula do divisor.
\resposta{$I_1=\SI{9}{\ampere}$, $I_2=\SI{6}{\ampere}$}
\resolucao{Para \emph{dois} ramos, a corrente de um deles é proporcional à
resistência do \textbf{outro}:
\[
I_1=I_t\frac{R_2}{R_1+R_2}=15\cdot\frac{12}{20}=\SI{9}{\ampere},
\qquad
I_2=15\cdot\frac{8}{20}=\SI{6}{\ampere}.
\]
\textbf{Verificação:} $9+6=\SI{15}{\ampere}$ \checkmark, e a maior corrente ficou
no \emph{menor} resistor \checkmark. O "cruzamento" dos índices é o ponto que
mais confunde --- ele existe porque $I_k\propto G_k$, e escrever a fórmula em
condutâncias ($I_1=I_tG_1/(G_1+G_2)$) elimina a inversão.}
\end{questao}

\begin{questao}[2]{}
Três resistores em paralelo sob $\SI{100}{\volt}$ drenam $\SI{5}{\ampere}$ no
total. Dois deles valem $\SI{50}{\ohm}$ e $\SI{100}{\ohm}$. Determine o
terceiro.
\resposta{$R_3=\SI{50}{\ohm}$}
\resolucao{Trabalhe em condutâncias, onde a subtração é direta:
\[
G_{eq}=\frac{I_t}{V}=\frac{5}{100}=\SI{0.05}{\siemens},
\qquad
G_1+G_2=0{,}02+0{,}01=\SI{0.03}{\siemens}.
\]
\[
G_3=0{,}05-0{,}03=\SI{0.02}{\siemens}
\;\Longrightarrow\;
R_3=\SI{50}{\ohm}.
\]
Tentar subtrair \emph{resistências} aqui levaria a um resultado sem sentido:
no paralelo, o que é aditivo é $G$, não $R$.}
\end{questao}

\begin{questao}[2]{}
Duas lâmpadas de resistências $\SI{240}{\ohm}$ e $\SI{960}{\ohm}$ são ligadas
\textbf{em paralelo} a $\SI{120}{\volt}$. Qual brilha mais? Compare com o
resultado da mesma dupla em série.
\resposta{Em paralelo, a de $\SI{240}{\ohm}$ ($\SI{60}{\watt}$ contra
$\SI{15}{\watt}$) --- inverte-se em relação à série}
\resolucao{
\[
P_{240}=\frac{120^{2}}{240}=\SI{60}{\watt},
\qquad
P_{960}=\frac{120^{2}}{960}=\SI{15}{\watt}.
\]
Em paralelo, $V$ é comum e $P\propto1/R$: brilha mais a de \emph{menor}
resistência. Em série (mesma dupla), $I$ é comum e $P\propto R$: brilharia mais
a de $\SI{960}{\ohm}$.\\
\textbf{Conclusão que vale guardar:} não existe "a lâmpada mais potente" em
termos absolutos --- a potência depende da \emph{ligação}. É por isso que
lâmpadas residenciais, sempre em paralelo, obedecem à intuição usual
($\SI{100}{\watt}$ brilha mais que $\SI{25}{\watt}$), mas guirlandas em série
se comportam ao contrário.}
\end{questao}

\begin{questao}[2]{}
Um galvanômetro tem resistência interna $\SI{100}{\ohm}$ e deflexão total com
$\SI{1}{\milli\ampere}$. Determine o resistor \emph{shunt} que amplia a escala
para $\SI{11}{\milli\ampere}$.
\resposta{$R_{sh}=\SI{10}{\ohm}$}
\resolucao{O shunt fica em paralelo com o galvanômetro e desvia o excedente:
$I_{sh}=11-1=\SI{10}{\milli\ampere}$. Como a tensão é comum,
\[
R_{sh}I_{sh}=R_gI_g
\;\Longrightarrow\;
R_{sh}=\frac{R_gI_g}{I_{sh}}=\frac{100\cdot1}{10}=\SI{10}{\ohm}.
\]
A regra geral é $R_{sh}=\dfrac{R_g}{n-1}$, com $n=I_{total}/I_g$ o fator de
ampliação --- aqui $n=11$ e $R_{sh}=100/10$ \checkmark.}
\end{questao}

\begin{questao}[2]{}
Um circuito residencial de $\SI{220}{\volt}$ alimenta, em paralelo, cargas de
$\SI{4400}{\watt}$, $\SI{1100}{\watt}$ e $\SI{550}{\watt}$. Determine a
corrente total e escolha o disjuntor.
\resposta{$I_t=\SI{27.5}{\ampere}$; disjuntor de $\SI{32}{\ampere}$}
\resolucao{Cada carga drena $I=P/V$:
\[
I_1=\frac{4400}{220}=20,\quad
I_2=\frac{1100}{220}=5,\quad
I_3=\frac{550}{220}=\SI{2.5}{\ampere}.
\]
Total: $\SI{27.5}{\ampere}$ (as correntes se somam porque as cargas estão em
paralelo).\\
O disjuntor comercial imediatamente superior é o de $\SI{32}{\ampere}$ ---
$\SI{25}{\ampere}$ seria insuficiente. \textbf{Atenção:} escolhido o disjuntor,
a seção do condutor deve suportar \emph{a corrente do disjuntor}, não a da
carga; caso contrário, o cabo vira o elo fraco da proteção.}
\end{questao}

\begin{questao}[2]{}
Dois condutores idênticos de $\SI{0.8}{\ohm}$ cada são ligados em paralelo para
alimentar uma carga. Determine a resistência resultante e explique a prática.
\resposta{$\SI{0.4}{\ohm}$ --- metade}
\resolucao{$R_{eq}=0{,}8/2=\SI{0.4}{\ohm}$.\\
Dobrar o número de condutores idênticos equivale a dobrar a seção: a
resistência cai à metade, e com ela a queda de tensão e as perdas
$RI^{2}$ na linha.\\
\textbf{Condição prática:} os condutores devem ter mesmo comprimento, seção e
material. Se um for mais curto, ele terá menor resistência e assumirá corrente
desproporcional --- podendo sobreaquecer enquanto o outro fica ocioso. Por isso
condutores em paralelo são sempre especificados com o mesmo lance.}
\end{questao}

\begin{questao}[2]{}
Um resistor de $\SI{1}{\mega\ohm}$ é ligado em paralelo com um de
$\SI{100}{\ohm}$. A resistência equivalente:
\begin{alternativas}[2]
\item fica praticamente igual a $\SI{100}{\ohm}$
\item fica praticamente igual a $\SI{1}{\mega\ohm}$
\item vale aproximadamente $\SI{500}{\kilo\ohm}$
\item vale exatamente $\SI{50}{\ohm}$
\end{alternativas}
\resposta{(a)}
\resolucao{$G=0{,}01+10^{-6}\approx\SI{0.010001}{\siemens}$, logo
$R_{eq}\approx\SI{99.99}{\ohm}$ --- diferença de $0{,}01\%$.\\
\textbf{Princípio geral:} no paralelo, quem manda é o \emph{menor} resistor;
ramos muito maiores são praticamente circuitos abertos. É exatamente por isso
que um voltímetro de alta impedância "não carrega" o circuito sob medição.}
\end{questao}

\begin{questao}[2]{}
Um chuveiro possui duas resistências de $\SI{22}{\ohm}$ que podem ser ligadas
isoladamente ou em paralelo, sob $\SI{220}{\volt}$. Determine a potência em
cada caso.
\resposta{$\SI{2200}{\watt}$ com uma; $\SI{4400}{\watt}$ com as duas}
\resolucao{Uma resistência:
$P=\dfrac{220^{2}}{22}=\SI{2200}{\watt}$.\\
Duas em paralelo: $R_{eq}=\SI{11}{\ohm}$ e
$P=\dfrac{220^{2}}{11}=\SI{4400}{\watt}$.\\
A potência \emph{dobra} porque a resistência caiu à metade sob a mesma tensão.
É assim que funcionam as chaves inverno/verão. \textbf{Consequência para a
instalação:} a corrente também dobra ($10\to\SI{20}{\ampere}$), e o
dimensionamento de disjuntor e condutor deve ser feito para a posição de maior
potência.}
\end{questao}

\begin{questao}[2]{}
Em uma associação paralela mediram-se as correntes de ramo $\SI{4}{\ampere}$,
$\SI{1}{\ampere}$ e $\SI{1}{\ampere}$, com corrente total de
$\SI{6}{\ampere}$ e tensão de $\SI{20}{\volt}$. Verifique a consistência e
determine as três resistências e a potência total.
\resposta{Consistente; $\SI{5}{\ohm}$, $\SI{20}{\ohm}$, $\SI{20}{\ohm}$;
$P=\SI{120}{\watt}$}
\resolucao{\textbf{LKC:} $4+1+1=\SI{6}{\ampere}$ \checkmark.\\
\textbf{Resistências:} $R_k=V/I_k$ $\Rightarrow$ $20/4=5$ e
$20/1=\SI{20}{\ohm}$ (duas vezes).\\
\textbf{Potência:} $P=VI_t=20\cdot6=\SI{120}{\watt}$; por ramo,
$80+20+20=\SI{120}{\watt}$ \checkmark.\\
Dois testes independentes --- LKC e balanço de potência --- fecham. Se apenas o
primeiro fechasse, o erro estaria na leitura da tensão.}
\end{questao}

\blocoaprofundamento

\begin{questao}[3]{}
Uma corrente total de $\SI{6}{\ampere}$ divide-se entre $\SI{5}{\ohm}$,
$\SI{20}{\ohm}$ e $\SI{20}{\ohm}$ em paralelo.
\begin{itensq}
\item Determine $R_{eq}$ e a tensão da associação.
\item Obtenha as três correntes pelo divisor em condutâncias.
\item Verifique por dois caminhos independentes.
\end{itensq}
\resposta{$R_{eq}=\SI{3.33}{\ohm}$, $V=\SI{20}{\volt}$;
$\SI{4}{\ampere}$, $\SI{1}{\ampere}$, $\SI{1}{\ampere}$}
\resolucao{(a) $G=0{,}2+0{,}05+0{,}05=\SI{0.3}{\siemens}$,
$R_{eq}=1/0{,}3=\SI{3.33}{\ohm}$ e $V=I_tR_{eq}=6\cdot3{,}33=\SI{20}{\volt}$.\\
(b) Com três ou mais ramos, a forma correta do divisor usa condutâncias:
\[
I_k=I_t\frac{G_k}{G_{eq}}
\;\Rightarrow\;
I_1=6\cdot\frac{0{,}2}{0{,}3}=4,\quad
I_2=I_3=6\cdot\frac{0{,}05}{0{,}3}=\SI{1}{\ampere}.
\]
\textbf{Alerta:} a fórmula $I_1=I_tR_2/(R_1+R_2)$ \emph{não} se generaliza para
três ramos --- ela é um caso particular que só funciona com dois.\\
(c) \emph{LKC:} $4+1+1=6$ \checkmark. \emph{Lei de Ohm por ramo:}
$20/5=4$ e $20/20=1$ \checkmark.}
\end{questao}

\begin{questao}[3]{}
Três ramos de $\SI{10}{\ohm}$, $\SI{20}{\ohm}$ e $\SI{20}{\ohm}$ operam sob
fonte \textbf{ideal} de $\SI{60}{\volt}$. O ramo de $\SI{10}{\ohm}$ se
\textbf{abre}. Determine o que muda nas correntes de ramo, na corrente total e
em $R_{eq}$.
\resposta{Ramos restantes inalterados ($\SI{3}{\ampere}$ cada);
$I_t$: $12\to\SI{6}{\ampere}$; $R_{eq}$: $5\to\SI{10}{\ohm}$}
\resolucao{Com fonte ideal, $V=\SI{60}{\volt}$ permanece fixo. Como cada ramo
depende \emph{só} da tensão comum e do próprio resistor, os dois ramos íntegros
continuam com $60/20=\SI{3}{\ampere}$ --- \textbf{nada muda para eles}.\\
O que muda é o agregado: $I_t=3+3=\SI{6}{\ampere}$ (queda de $50\%$) e
$R_{eq}=\SI{10}{\ohm}$.\\
\textbf{Por isso a ligação residencial é em paralelo:} a queima de uma lâmpada
não afeta as demais. Em série, uma falha derruba o circuito inteiro.\\
\emph{Ressalva importante:} isso depende de a fonte ser ideal. Com resistência
interna, $V$ sobe quando um ramo abre --- caso tratado no bloco de desafio.}
\end{questao}

\begin{questao}[3]{}
Um dos ramos de uma associação paralela sofre \textbf{curto-circuito}. Analise
$R_{eq}$, a corrente total e o efeito sobre os demais ramos, e justifique a
necessidade de proteção.
\resposta{$R_{eq}\to0$; $I_t\to\infty$ (limitada só pela fonte e pela fiação);
os demais ramos ficam sem tensão}
\resolucao{Um ramo em curto tem $R=0$, portanto $G\to\infty$. Como
$G_{eq}=\sum G_k$, a condutância total explode e $R_{eq}\to0$:
\[
R_{eq}=\frac{1}{\infty+G_2+G_3}\to0.
\]
Com fonte ideal, $I_t=V/R_{eq}\to\infty$. Na prática, a corrente é limitada
pela resistência interna da fonte e pela impedância da fiação --- mas ainda
assim atinge dezenas ou centenas de ampères.\\
\textbf{Efeito sobre os demais:} o curto derruba a tensão do barramento a
praticamente zero. Os outros ramos, embora íntegros, \emph{param de funcionar}
--- eles são vítimas, não causa.\\
\textbf{Proteção:} é este o cenário que o disjuntor combate. Diferentemente do
ramo aberto (falha benigna e localizada), o curto em um ramo compromete o
sistema inteiro e libera energia suficiente para incendiar a instalação. Daí a
assimetria de tratamento: monitoram-se \emph{sobrecorrentes}, não
subcorrentes.}
\end{questao}

\begin{questao}[3]{}
Dispõe-se de resistores de $\SI{220}{\ohm}$ com dissipação máxima de
$\SI{1}{\watt}$. Projete uma associação paralela de $\SI{55}{\ohm}$ capaz de
dissipar pelo menos $\SI{4}{\watt}$, e determine a tensão máxima aplicável.
\resposta{Quatro em paralelo; $P_{\max}=\SI{4}{\watt}$;
$V_{\max}=\SI{14.8}{\volt}$}
\resolucao{\textbf{Valor:} $220/n=55\Rightarrow n=4$.\\
\textbf{Potência:} sendo iguais e sob a mesma tensão, os quatro dividem a carga
térmica igualmente, então $P_{tot}=4\times1=\SI{4}{\watt}$ \checkmark.\\
\textbf{Tensão máxima:} limita-se pelo componente individual:
\[
V_{\max}=\sqrt{P_1R_1}=\sqrt{1\cdot220}=\SI{14.8}{\volt}.
\]
\textbf{Conferência:} $P_{tot}=V^{2}/R_{eq}=220/55=\SI{4}{\watt}$ \checkmark.\\
\textbf{Observação de projeto:} o paralelo multiplica a capacidade térmica
\emph{e} divide o valor por $n$; a série multiplica ambos. Se o objetivo exige
valor \emph{maior} com mais potência, use série; valor \emph{menor}, paralelo.
Para manter o valor e aumentar a potência, é preciso uma matriz
série--paralelo de $n\times n$ unidades.}
\end{questao}

\begin{questao}[3]{}
Uma rede tem cinco ramos em paralelo: $\SI{10}{\ohm}$, $\SI{20}{\ohm}$,
$\SI{20}{\ohm}$, $\SI{40}{\ohm}$ e $\SI{40}{\ohm}$. Calcule $R_{eq}$ e explique
por que o cálculo em condutâncias é preferível.
\resposta{$G=\SI{0.25}{\siemens}$; $R_{eq}=\SI{4}{\ohm}$}
\resolucao{
\[
G=0{,}100+0{,}050+0{,}050+0{,}025+0{,}025=\SI{0.25}{\siemens}
\;\Longrightarrow\;
R_{eq}=\SI{4}{\ohm}.
\]
\textbf{Três vantagens da condutância.} \emph{(i)} A operação é uma soma
simples, com um único inverso no fim --- contra quatro aplicações sucessivas da
regra produto/soma. \emph{(ii)} A distribuição de corrente sai de graça:
$I_k/I_t=G_k/G_{eq}$, aqui $40\%$, $20\%$, $20\%$, $10\%$ e $10\%$.
\emph{(iii)} Para incerteza, a propagação em $G$ é \emph{linear}, ao passo que
em $R$ envolve derivadas de uma função racional.\\
É também a razão pela qual a análise nodal --- construída sobre a LKC --- adota
condutâncias como grandeza natural.}
\end{questao}

\begin{questao}[3]{}
Um galvanômetro de $\SI{99}{\ohm}$ e $\SI{1}{\milli\ampere}$ de fundo de escala
deve medir até $\SI{100}{\milli\ampere}$.
\begin{itensq}
\item Determine o shunt.
\item Calcule a resistência do amperímetro resultante.
\item Explique por que essa resistência precisa ser pequena.
\end{itensq}
\resposta{(a) $R_{sh}=\SI{1}{\ohm}$; (b) $\SI{0.99}{\ohm}$}
\resolucao{(a) O shunt conduz $100-1=\SI{99}{\milli\ampere}$:
\[
R_{sh}=\frac{R_gI_g}{I_{sh}}=\frac{99\cdot1}{99}=\SI{1}{\ohm}.
\]
(b) $R_A=R_g\parallel R_{sh}=\dfrac{99\cdot1}{100}=\SI{0.99}{\ohm}$.\\
(c) O amperímetro é inserido \textbf{em série} no ramo sob medição, somando-se
à resistência do circuito. Se $R_A$ não for desprezível diante do circuito, a
própria medição altera a corrente que se pretende medir --- o chamado
\emph{erro de inserção}.\\
\textbf{Exemplo:} medindo um ramo de $\SI{100}{\ohm}$, o erro é
$0{,}99/100\approx1\%$; medindo um de $\SI{1}{\ohm}$, seria de $50\%$.\\
Note a dualidade: o \emph{amperímetro} deve ter resistência quase nula e vai em
série; o \emph{voltímetro}, resistência quase infinita e vai em paralelo.}
\end{questao}

\begin{questao}[3]{}
Uma carga de $\SI{100}{\ohm}$ opera a $\SI{200}{\volt}$. A isolação
degradada cria um caminho de fuga de $\SI{10}{\kilo\ohm}$ em paralelo.
\begin{itensq}
\item Determine a corrente de fuga e a corrente total.
\item Calcule $R_{eq}$ e o erro relativo cometido ao ignorar a fuga.
\item Comente o risco.
\end{itensq}
\resposta{(a) $\SI{20}{\milli\ampere}$ e $\SI{2.02}{\ampere}$;
(b) $\SI{99.01}{\ohm}$, erro de $1\%$; (c) risco de choque}
\resolucao{(a) $I_{carga}=200/100=\SI{2}{\ampere}$;
$I_{fuga}=200/10000=\SI{20}{\milli\ampere}$; total $\SI{2.02}{\ampere}$.\\
(b) $R_{eq}=\dfrac{100\cdot10000}{10100}=\SI{99.01}{\ohm}$ --- desvio de
$\approx1\%$ em relação a $\SI{100}{\ohm}$.\\
(c) Do ponto de vista \emph{elétrico}, $1\%$ é irrelevante: nenhum instrumento
de painel acusaria a anomalia. Do ponto de vista da \textbf{segurança}, é
gravíssimo: $\SI{20}{\milli\ampere}$ atravessando um corpo humano já provoca
tetanização muscular, e a partir de $\SI{30}{\milli\ampere}$ o risco de
fibrilação é real.\\
\textbf{Por isso existe o DR (dispositivo diferencial residual):} ele não mede
a corrente total, e sim o \emph{desequilíbrio} entre fase e neutro --- que é
exatamente a corrente de fuga. Um DR de $\SI{30}{\milli\ampere}$ detecta o que
os $\SI{2}{\ampere}$ da carga mascaram completamente.}
\end{questao}

\begin{questao}[3]{}
$R_1=\SI{100}{\ohm}\pm5\%$ e $R_2=\SI{400}{\ohm}\pm5\%$ estão em paralelo.
Determine $R_{eq}$ nominal e seus extremos de pior caso.
\resposta{$R_{eq}=\SI{80}{\ohm}$, entre $\SI{76}{\ohm}$ e $\SI{84}{\ohm}$
($\pm5\%$)}
\resolucao{Nominal: $R_{eq}=\dfrac{100\cdot400}{500}=\SI{80}{\ohm}$.\\
Como $R_{eq}$ cresce monotonicamente com $R_1$ e com $R_2$, os extremos ocorrem
com \emph{ambos} nos extremos:
\[
R_{eq}^{\min}=\frac{95\cdot380}{475}=\SI{76}{\ohm},
\qquad
R_{eq}^{\max}=\frac{105\cdot420}{525}=\SI{84}{\ohm}.
\]
Ou seja, $80\pm4$, isto é, $\pm5\%$ --- exatamente a tolerância das parcelas.\\
\textbf{Resultado notável:} o paralelo \emph{também} preserva a tolerância
relativa, assim como a série. Isso não é coincidência, e a razão é demonstrada
no bloco de desafio.}
\end{questao}

\begin{questao}[3]{}
$R_1=\SI{100}{\ohm}$ com $\alpha=\SI{0.004}{\per\celsius}$ está em paralelo com
$R_2=\SI{100}{\ohm}$ de coeficiente desprezível. Determine $R_{eq}$ a
$\SI{20}{\celsius}$ e a $\SI{80}{\celsius}$, e calcule a variação percentual.
\resposta{$\SI{50}{\ohm}$ e $\SI{55.36}{\ohm}$; variação de $+10{,}7\%$}
\resolucao{A $\SI{20}{\celsius}$: $R_{eq}=100/2=\SI{50}{\ohm}$.\\
A $\SI{80}{\celsius}$: $R_1=100[1+0{,}004(60)]=\SI{124}{\ohm}$ e
\[
R_{eq}=\frac{124\cdot100}{224}=\SI{55.36}{\ohm}.
\]
Variação: $(55{,}36-50)/50=+10{,}7\%$.\\
\textbf{Comparação instrutiva:} $R_1$ variou $24\%$, mas $R_{eq}$ variou
$10{,}7\%$ --- o ramo estável "amortece" a deriva. Curiosamente, o valor é o
mesmo obtido para a \emph{corrente} da associação série equivalente: em ambos
os casos o elemento estável divide o efeito aproximadamente pela metade, porque
as duas resistências são comparáveis.}
\end{questao}

\begin{questao}[3]{}
Um resistor de $\SI{1000}{\ohm}$ está em paralelo com um de $\SI{10}{\ohm}$.
\begin{itensq}
\item Calcule $R_{eq}$.
\item Recalcule supondo que $R_1$ varie $+10\%$.
\item Interprete o resultado em termos de sensibilidade.
\end{itensq}
\resposta{(a) $\SI{9.901}{\ohm}$; (b) $\SI{9.910}{\ohm}$; (c) variação de
apenas $0{,}09\%$}
\resolucao{(a) $R_{eq}=\dfrac{1000\cdot10}{1010}=\SI{9.901}{\ohm}$.\\
(b) Com $R_1=\SI{1100}{\ohm}$:
$R_{eq}=\dfrac{1100\cdot10}{1110}=\SI{9.910}{\ohm}$.\\
(c) Uma variação de $10\%$ em $R_1$ produziu $0{,}09\%$ em $R_{eq}$ ---
atenuação de mais de cem vezes. O ramo dominante ($\SI{10}{\ohm}$) fixa
praticamente sozinho o resultado.\\
\textbf{Uso prático:} é assim que se faz um ajuste \emph{fino}. Colocando um
resistor grande em paralelo com um pequeno, obtém-se um deslocamento pequeno e
controlado --- técnica usada em calibração de shunts e de divisores de
precisão. A quantificação exata dessa sensibilidade aparece no bloco de
desafio.}
\end{questao}

\begin{questao}[3]{}
As cargas da questão residencial ($\SI{27.5}{\ampere}$ totais) são alimentadas
por um circuito cuja fiação tem $\SI{0.2}{\ohm}$ de ida e volta, a partir de um
quadro de $\SI{220}{\volt}$.
\begin{itensq}
\item Determine a queda de tensão na fiação e a tensão nas cargas.
\item Calcule a perda de potência na fiação.
\item Avalie se a queda atende ao limite usual de $4\%$.
\end{itensq}
\resposta{(a) $\SI{5.5}{\volt}$; $\SI{214.5}{\volt}$;
(b) $\SI{151.25}{\watt}$; (c) $2{,}5\%$ --- atende}
\resolucao{(a) $\Delta V=R_fI=0{,}2\cdot27{,}5=\SI{5.5}{\volt}$, logo as cargas
recebem $220-5{,}5=\SI{214.5}{\volt}$.\\
(b) $P_f=R_fI^{2}=0{,}2(27{,}5)^{2}=\SI{151.25}{\watt}$ --- dissipados em puro
aquecimento do cabo.\\
(c) $5{,}5/220=2{,}5\%$, abaixo do limite típico de $4\%$ para circuitos
terminais. \checkmark\\
\textbf{Observação sutil:} a fiação está em \emph{série} com o conjunto, mesmo
que as cargas estejam em paralelo entre si. É por isso que ligar mais uma carga
--- que não afeta as demais em uma rede ideal --- na prática \emph{reduz} a
tensão de todas: a corrente total sobe e a queda na fiação também. O efeito é
visível quando as luzes piscam ao ligar um chuveiro.}
\end{questao}

\begin{questao}[3]{}
Deseja-se obter $\SI{75}{\ohm}$ dispondo de um resistor de $\SI{100}{\ohm}$.
Determine o resistor a ser associado em paralelo e discuta a viabilidade.
\resposta{$R=\SI{300}{\ohm}$}
\resolucao{
\[
\frac{1}{75}=\frac{1}{100}+\frac{1}{R}
\;\Longrightarrow\;
\frac{1}{R}=\frac{1}{75}-\frac{1}{100}=\frac{4-3}{300}
\;\Longrightarrow\;
R=\SI{300}{\ohm}.
\]
\textbf{Viabilidade:} $\SI{300}{\ohm}$ existe na série E24, e o alvo
($\SI{75}{\ohm}$) está a $25\%$ do resistor de partida --- folga confortável.
A precisão do resultado depende principalmente do resistor de
$\SI{100}{\ohm}$, que domina o paralelo.\\
\textbf{Limite do método:} só se pode \emph{reduzir} o valor de partida; alvos
acima de $\SI{100}{\ohm}$ exigem série. E, como se verá no desafio, alvos muito
próximos de $\SI{100}{\ohm}$ tornam a solução impraticável.}
\end{questao}

\begin{questao}[3]{}
Compare a mesma dupla $R_1=\SI{30}{\ohm}$ e $R_2=\SI{60}{\ohm}$ ligada em série
e em paralelo, sob $\SI{90}{\volt}$: resistência, corrente total, distribuição
de corrente e de potência.
\resposta{Série: $\SI{90}{\ohm}$, $\SI{1}{\ampere}$, $P=30$ e
$\SI{60}{\watt}$. Paralelo: $\SI{20}{\ohm}$, $\SI{4.5}{\ampere}$,
$P=270$ e $\SI{135}{\watt}$}
\resolucao{\textbf{Série:} $R=\SI{90}{\ohm}$, $I=\SI{1}{\ampere}$ (comum),
$P_1=30$, $P_2=\SI{60}{\watt}$; total $\SI{90}{\watt}$. O \emph{maior}
resistor dissipa mais.\\
\textbf{Paralelo:} $R=\SI{20}{\ohm}$, $I_1=3$, $I_2=\SI{1.5}{\ampere}$,
$I_t=\SI{4.5}{\ampere}$, $P_1=\SI{270}{\watt}$, $P_2=\SI{135}{\watt}$; total
$\SI{405}{\watt}$. O \emph{menor} resistor dissipa mais.\\
\textbf{Razões:} $R_{s}/R_{p}=90/20=4{,}5$, e a potência total é $4{,}5$ vezes
maior no paralelo --- o mesmo fator, pois $P=V^{2}/R$ com $V$ fixo.\\
\textbf{Síntese:} de dois resistores obtêm-se quatro valores
($30$, $60$, $20$ e $\SI{90}{\ohm}$) e duas distribuições opostas de potência.
A ligação é uma variável de projeto tão importante quanto os componentes.}
\end{questao}

\blocodesafio

\begin{questao}[4]{}
\textbf{(Projeto de instrumentação.)} Um galvanômetro tem $R_g=\SI{99}{\ohm}$ e
fundo de escala $I_g=\SI{1}{\milli\ampere}$. Deseja-se um amperímetro de três
escalas: $\SI{10}{\milli\ampere}$, $\SI{100}{\milli\ampere}$ e
$\SI{1}{\ampere}$.
\begin{itensq}
\item Determine os três shunts na configuração de shunts comutados.
\item Aponte o risco dessa configuração e descreva a alternativa (shunt de
      Ayrton).
\item Calcule a resistência do instrumento em cada escala.
\end{itensq}
\resposta{(a) $\SI{11}{\ohm}$, $\SI{1}{\ohm}$ e $\SI{0.0991}{\ohm}$;
(c) $\SI{9.9}{\ohm}$, $\SI{0.99}{\ohm}$ e $\SI{0.099}{\ohm}$}
\resolucao{(a) Com $n=I_{\max}/I_g$, vale
$R_{sh}=\dfrac{R_g}{n-1}$:
\[
n=10\Rightarrow\frac{99}{9}=\SI{11}{\ohm};\quad
n=100\Rightarrow\frac{99}{99}=\SI{1}{\ohm};\quad
n=1000\Rightarrow\frac{99}{999}=\SI{0.0991}{\ohm}.
\]
(b) \textbf{Risco:} durante a comutação, a chave passa por um instante em que
\emph{nenhum} shunt está conectado. Toda a corrente de linha --- até
$\SI{1}{\ampere}$ --- atravessaria o galvanômetro, projetado para
$\SI{1}{\milli\ampere}$: destruição imediata.\\
\textbf{Shunt de Ayrton:} os resistores ficam permanentemente ligados em série
entre si e o conjunto em paralelo com o galvanômetro; a chave apenas escolhe
\emph{em que ponto} da cadeia a corrente entra. Assim o galvanômetro nunca
fica desprotegido, mesmo durante a comutação.\\
(c) $R_A=R_g\parallel R_{sh}$: $99\parallel11=\SI{9.9}{\ohm}$;
$99\parallel1=\SI{0.99}{\ohm}$; $99\parallel0{,}0991=\SI{0.099}{\ohm}$.\\
Note que $R_A=R_g/n$ exatamente --- escalas maiores exigem instrumentos de
resistência menor, o que é justamente o desejável.}
\end{questao}

\begin{questao}[4]{}
\textbf{(Demonstração --- escalonamento de bancos.)} Considere $N$ resistores
iguais de valor $R$ e potência nominal $P_1$.
\begin{itensq}
\item Mostre que a associação paralela vale $R/N$ e suporta $NP_1$.
\item Mostre que a série vale $NR$ e também suporta $NP_1$.
\item Deduza a matriz $m\times n$ necessária para manter o valor $R$ e obter
      potência $mnP_1$.
\end{itensq}
\resposta{(c) $m$ colunas em paralelo de $n$ unidades em série, com $m=n$}
\resolucao{(a) $G_{eq}=NG\Rightarrow R_{eq}=R/N$. Sob tensão $V$ comum, cada
unidade dissipa $V^{2}/R$; como todas são iguais e submetidas à mesma tensão, o
total é $N$ vezes o individual: $P_{\max}=NP_1$.\\
(b) $R_{eq}=NR$. Com corrente $I$ comum, cada unidade dissipa $RI^{2}$ e o
total é novamente $NP_1$.\\
\textbf{Fato central:} \emph{qualquer} associação de $N$ resistores iguais
suporta $NP_1$, desde que a dissipação se reparta igualmente. O que a topologia
escolhe é o \emph{valor}, não a capacidade térmica.\\
(c) Uma matriz de $n$ unidades em série, repetida em $m$ colunas paralelas:
\[
R_{eq}=\frac{nR}{m},\qquad P_{\max}=mnP_1.
\]
Para preservar o valor, $n=m$; então $R_{eq}=R$ e
$P_{\max}=n^{2}P_1$.\\
\textbf{Exemplo:} para quadruplicar a potência sem alterar o valor, use
$2\times2=4$ unidades. Para $\times9$, use $3\times3=9$. A potência cresce com
o \emph{quadrado} do lado da matriz --- e o custo, também.}
\end{questao}

\begin{questao}[4]{}
\textbf{(Sensibilidade.)} Para $R_{eq}=\dfrac{R_1R_2}{R_1+R_2}$:
\begin{itensq}
\item Obtenha $\dfrac{\partial R_{eq}}{\partial R_1}$ e mostre que vale
      $\left(\dfrac{R_2}{R_1+R_2}\right)^{2}$.
\item Obtenha a sensibilidade relativa e mostre que
      $S_1+S_2=1$.
\item Avalie para $R_1=\SI{100}{\ohm}$ e $R_2=\SI{400}{\ohm}$, e interprete.
\end{itensq}
\resposta{(a) $0{,}64$; (b) $S_1=R_2/(R_1+R_2)$, $S_2=R_1/(R_1+R_2)$;
(c) $S_1=0{,}8$ e $S_2=0{,}2$}
\resolucao{(a) Derivando o quociente em relação a $R_1$:
\[
\frac{\partial R_{eq}}{\partial R_1}
=\frac{R_2(R_1+R_2)-R_1R_2}{(R_1+R_2)^{2}}
=\frac{R_2^{2}}{(R_1+R_2)^{2}}
=\left(\frac{R_2}{R_1+R_2}\right)^{2}.
\]
Com $R_1=100$ e $R_2=400$: $(400/500)^{2}=0{,}64$ --- adimensional, pois é
$\si{\ohm}$ por $\si{\ohm}$.\\
(b)
\[
S_1=\frac{\partial R_{eq}}{\partial R_1}\cdot\frac{R_1}{R_{eq}}
=\frac{R_2^{2}}{(R_1+R_2)^{2}}\cdot\frac{R_1(R_1+R_2)}{R_1R_2}
=\frac{R_2}{R_1+R_2}.
\]
Por simetria, $S_2=\dfrac{R_1}{R_1+R_2}$, e a soma é $1$ --- consequência de
$R_{eq}$ ser \textbf{homogênea de grau 1}.\\
(c) $S_1=0{,}8$ e $S_2=0{,}2$: um erro de $1\%$ em $R_1$ produz $0{,}8\%$ em
$R_{eq}$, e o mesmo erro em $R_2$ produz apenas $0{,}2\%$.\\
\textbf{Regra de projeto:} o resistor \emph{menor} domina o paralelo e deve
receber a tolerância apertada --- exatamente o inverso da série, onde manda o
maior. E se $R_2\gg R_1$, então $S_1\to1$ e $S_2\to0$: o ramo grande vira
ajuste fino, o que confirma quantitativamente o observado na questão do
$\SI{1000}{\ohm}$ com $\SI{10}{\ohm}$.}
\end{questao}

\begin{questao}[4]{}
\textbf{(Divisão de corrente sob tolerância.)} Três shunts iguais de
$\SI{0.01}{\ohm}$ devem dividir igualmente $\SI{30}{\ampere}$.
\begin{itensq}
\item Determine a corrente ideal por shunt e a tensão sobre eles.
\item Se um deles for $2\%$ menor e os outros dois nominais, calcule o
      desequilíbrio de corrente.
\item Discuta a implicação térmica.
\end{itensq}
\resposta{(a) $\SI{10}{\ampere}$ e $\SI{0.1}{\volt}$;
(b) $\SI{10.13}{\ampere}$ contra $\SI{9.93}{\ampere}$ --- desvio de $1{,}3\%$;
(c) o mais "fraco" aquece mais}
\resolucao{(a) $I=30/3=\SI{10}{\ampere}$;
$V=0{,}01\cdot10=\SI{0.1}{\volt}$.\\
(b) Com $R_a=\SI{0.0098}{\ohm}$ e $R_b=R_c=\SI{0.0100}{\ohm}$:
\[
G=\frac{1}{0{,}0098}+\frac{2}{0{,}0100}=102{,}04+200{,}00=\SI{302.04}{\siemens},
\]
\[
I_a=30\cdot\frac{102{,}04}{302{,}04}=\SI{10.13}{\ampere},
\qquad
I_b=I_c=30\cdot\frac{100{,}00}{302{,}04}=\SI{9.93}{\ampere}.
\]
\textbf{Verificação:} $10{,}13+9{,}93+9{,}93=\SI{29.99}{\ampere}$ \checkmark.\\
\textbf{Fato tranquilizador:} um desvio de $2\%$ na resistência produziu apenas
$1{,}3\%$ de desequilíbrio --- a divisão é \emph{mais estável} que os
componentes, porque os outros dois ramos "puxam" a média.\\
(c) $P_a=0{,}0098(10{,}13)^{2}=\SI{1.006}{\watt}$ contra
$P_b=0{,}0100(9{,}93)^{2}=\SI{0.986}{\watt}$: o shunt de menor resistência
dissipa $2\%$ a mais. Se tiver coeficiente térmico positivo, aquecerá,
aumentará sua resistência e \emph{devolverá} corrente aos vizinhos --- uma
realimentação \textbf{estabilizadora}. Com coeficiente negativo ocorreria o
oposto: fuga térmica. É por isso que resistores para divisão de corrente são
feitos em ligas de $\alpha$ positivo e pequeno, como a manganina.}
\end{questao}

\begin{questao}[4]{}
\textbf{(Compensação térmica em paralelo.)} Deseja-se uma associação paralela
com resistência independente da temperatura, usando
$R_1=\SI{100}{\ohm}$ com $\alpha_1=+4000\,\mathrm{ppm/^{\circ}C}$ e um segundo
resistor com $\alpha_2=-1000\,\mathrm{ppm/^{\circ}C}$.
\begin{itensq}
\item Deduza a condição de compensação em termos de condutâncias.
\item Determine $R_2$ e $R_{eq}$.
\item Verifique a compensação e comente sua exatidão.
\end{itensq}
\resposta{(a) $G_1\alpha_1+G_2\alpha_2=0$; (b) $R_2=\SI{25}{\ohm}$,
$R_{eq}=\SI{20}{\ohm}$; (c) exata só em primeira ordem}
\resolucao{(a) Partindo de $G_{eq}=G_1+G_2$ e de
$R_k(T)=R_k(1+\alpha_k\Delta T)$, tem-se
$G_k\approx G_k(1-\alpha_k\Delta T)$ para $\alpha_k\Delta T\ll1$. Então
\[
G_{eq}(T)\approx(G_1+G_2)-(G_1\alpha_1+G_2\alpha_2)\Delta T,
\]
e a condição de estabilidade é
\[
G_1\alpha_1+G_2\alpha_2=0
\quad\Longleftrightarrow\quad
\frac{\alpha_1}{R_1}+\frac{\alpha_2}{R_2}=0.
\]
(b) $R_2=-\dfrac{\alpha_2R_1}{\alpha_1}=\dfrac{1000\cdot100}{4000}
=\SI{25}{\ohm}$, e $R_{eq}=100\parallel25=\SI{20}{\ohm}$.\\
(c) \emph{Para $\Delta T=\SI{1}{\celsius}$:}
$R_1=\SI{100.4}{\ohm}$, $R_2=\SI{24.975}{\ohm}$, $R_{eq}=\SI{20.000}{\ohm}$ ---
compensação praticamente perfeita.\\
\emph{Para $\Delta T=\SI{50}{\celsius}$:} $R_1=\SI{120}{\ohm}$,
$R_2=\SI{23.75}{\ohm}$, $R_{eq}=\SI{19.83}{\ohm}$ --- desvio de $0{,}9\%$.\\
A compensação é de \textbf{primeira ordem}: excelente perto da temperatura de
projeto, degradando-se quadraticamente conforme se afasta.\\
\textbf{Compare com a série}, cuja condição era $R_1\alpha_1+R_2\alpha_2=0$
(ponderada por resistências). Aqui a ponderação é por \emph{condutâncias} ---
mais uma manifestação da dualidade série/paralelo.}
\end{questao}

\begin{questao}[4]{}
\textbf{(Falha de ramo com fonte real.)} Uma fonte de $\SI{24}{\volt}$ com
$r=\SI{2}{\ohm}$ alimenta três ramos de $\SI{6}{\ohm}$ em paralelo.
\begin{itensq}
\item Determine a tensão do barramento, a corrente total e a de cada ramo.
\item Um ramo se abre. Recalcule tudo.
\item Explique por que os ramos remanescentes \emph{mudam}, ao contrário do
      caso com fonte ideal.
\end{itensq}
\resposta{(a) $\SI{12}{\volt}$, $\SI{6}{\ampere}$, $\SI{2}{\ampere}$ por ramo;
(b) $\SI{14.4}{\volt}$, $\SI{4.8}{\ampere}$, $\SI{2.4}{\ampere}$ por ramo}
\resolucao{(a) $R_{eq}=6/3=\SI{2}{\ohm}$;
\[
I_t=\frac{24}{2+2}=\SI{6}{\ampere},\qquad
V_{barra}=24-2(6)=\SI{12}{\volt},\qquad
I_{ramo}=\frac{12}{6}=\SI{2}{\ampere}.
\]
(b) Com dois ramos, $R_{eq}=\SI{3}{\ohm}$:
\[
I_t=\frac{24}{2+3}=\SI{4.8}{\ampere},\qquad
V_{barra}=24-2(4{,}8)=\SI{14.4}{\volt},\qquad
I_{ramo}=\frac{14{,}4}{6}=\SI{2.4}{\ampere}.
\]
(c) Os ramos remanescentes \emph{não mudaram de valor}, mas a tensão que os
alimenta subiu de $12$ para $\SI{14.4}{\volt}$ ($+20\%$), e com ela a corrente
de cada um ($+20\%$).\\
\textbf{Mecanismo:} $r$ está em série com o conjunto. Menos ramos $\Rightarrow$
menor corrente total $\Rightarrow$ menor queda interna $\Rightarrow$ maior
tensão disponível.\\
\textbf{Implicação de projeto:} o isolamento entre ramos paralelos é uma
\emph{idealização}. Com fonte real, ramos interagem via $r$ --- e a perda de
uma carga pode sobrecarregar as demais em $20\%$ ou mais. Fontes reguladas
existem justamente para levar $r$ efetiva a valores próximos de zero e
restaurar o desacoplamento.}
\end{questao}

\begin{questao}[4]{}
\textbf{(Limite do método de ajuste.)} Dispõe-se de um resistor fixo de
$\SI{100}{\ohm}$ e deseja-se obter valores-alvo por associação paralela.
\begin{itensq}
\item Deduza $R(alvo)$ e calcule para alvos de $\SI{90}{\ohm}$,
      $\SI{95}{\ohm}$ e $\SI{99}{\ohm}$.
\item Obtenha a sensibilidade do alvo em relação a $R$ e avalie no caso de
      $\SI{99}{\ohm}$.
\item Conclua sobre a viabilidade prática.
\end{itensq}
\resposta{(a) $\SI{900}{\ohm}$, $\SI{1900}{\ohm}$ e $\SI{9900}{\ohm}$;
(c) inviável para alvos próximos do resistor fixo}
\resolucao{(a) De $\dfrac{1}{R_{alvo}}=\dfrac{1}{R_f}+\dfrac{1}{R}$:
\[
R=\frac{R_fR_{alvo}}{R_f-R_{alvo}}.
\]
\[
90\to\frac{100\cdot90}{10}=\SI{900}{\ohm};\quad
95\to\frac{100\cdot95}{5}=\SI{1900}{\ohm};\quad
99\to\frac{100\cdot99}{1}=\SI{9900}{\ohm}.
\]
O denominador $R_f-R_{alvo}$ tende a zero e $R$ \textbf{diverge}.\\
(b) Pela sensibilidade relativa do paralelo,
$S=\dfrac{R}{R_f+R}$. Para $R=\SI{9900}{\ohm}$:
$S=9900/10000=0{,}99$.\\
Ou seja: $1\%$ de erro em $R$ produz $0{,}99\%$ de erro no alvo. Como o ajuste
pretendido é de apenas $1\%$ (de $100$ para $\SI{99}{\ohm}$), \emph{o erro tem
a mesma ordem de grandeza do ajuste} --- ele o anula.\\
(c) \textbf{Inviável.} Com resistor de $5\%$ de tolerância, o "ajuste" de
$\SI{1}{\ohm}$ vem com incerteza de $\SI{0.99}{\ohm}$: o alvo pode ficar em
qualquer ponto entre $98$ e $\SI{100}{\ohm}$. Além disso,
$\SI{9900}{\ohm}$ não é valor comercial, e a resistência de contato já é
comparável ao efeito buscado.\\
\textbf{Regra prática:} ajuste por paralelo só é confiável para correções de
\emph{dezenas} de por cento. Para correções finas, use potenciômetro de ajuste
ou selecione o componente por medição direta.}
\end{questao}

\begin{questao}[4]{}
\textbf{(Demonstração --- propagação de incerteza.)} Mostre que, se
\emph{todos} os resistores de uma rede puramente resistiva tiverem a mesma
tolerância relativa $t$, então $R_{eq}$ terá exatamente a tolerância $t$,
qualquer que seja a topologia.
\begin{itensq}
\item Demonstre a propriedade de homogeneidade.
\item Verifique em um exemplo série e em um paralelo.
\item Discuta o que muda quando as tolerâncias são diferentes.
\end{itensq}
\resposta{$R_{eq}$ é homogênea de grau 1: $R_{eq}(\lambda R_k)=\lambda
R_{eq}(R_k)$; logo a tolerância relativa se conserva}
\resolucao{(a) Qualquer $R_{eq}$ de rede resistiva é construída por somas
($R_i+R_j$) e por combinações do tipo $\dfrac{R_iR_j}{R_i+R_j}$ --- ambas
\textbf{homogêneas de grau 1}. Multiplicando \emph{todos} os resistores por
$\lambda$:
\[
\lambda R_i+\lambda R_j=\lambda(R_i+R_j),
\qquad
\frac{(\lambda R_i)(\lambda R_j)}{\lambda R_i+\lambda R_j}
=\lambda\frac{R_iR_j}{R_i+R_j}.
\]
Por indução sobre a construção da rede, $R_{eq}(\lambda R_k)=\lambda R_{eq}$.
Fazendo $\lambda=1\pm t$, obtém-se $R_{eq}(1\pm t)$: tolerância relativa
exatamente $t$.\\
(b) \emph{Série:} $100\pm5\%$ e $200\pm5\%$ dão $300\pm\SI{15}{\ohm}=\pm5\%$
\checkmark. \emph{Paralelo:} $100\pm5\%$ e $400\pm5\%$ dão
$80\pm\SI{4}{\ohm}=\pm5\%$ \checkmark. A demonstração explica por que os dois
casos, aparentemente distintos, deram o mesmo resultado.\\
(c) Com tolerâncias \emph{diferentes}, a homogeneidade não se aplica --- os
resistores não escalam pelo mesmo $\lambda$. Vale então a combinação ponderada
pelas sensibilidades:
\[
t_{eq}=\sum_k S_k\,t_k,
\qquad \text{com } \sum_k S_k=1.
\]
Em série, $S_k=R_k/\sum R_j$ (manda o maior); em paralelo,
$S_k=G_k/\sum G_j$ (manda o menor). A restrição $\sum S_k=1$ é justamente a
identidade de Euler para funções homogêneas de grau 1 --- e garante que
$t_{eq}$ jamais ultrapasse a maior das tolerâncias individuais.}
\end{questao}

\begin{questao}[4]{}
\textbf{(Seletividade da proteção.)} Um barramento de $\SI{220}{\volt}$
alimenta três circuitos em paralelo, com correntes nominais de
$\SI{20}{\ampere}$, $\SI{5}{\ampere}$ e $\SI{2.5}{\ampere}$, protegidos
individualmente e por um disjuntor geral.
\begin{itensq}
\item Especifique os disjuntores de ramo e o geral.
\item Explique o requisito de seletividade e o que ocorre se ele falhar.
\item Discuta por que a soma dos disjuntores de ramo pode exceder o geral.
\end{itensq}
\resposta{(a) ramos de $25$, $10$ e $\SI{10}{\ampere}$; geral de
$\SI{32}{\ampere}$; (b) o ramo deve atuar antes do geral}
\resolucao{(a) Cada disjuntor de ramo fica acima da corrente nominal do próprio
circuito: $\SI{25}{\ampere}$, $\SI{10}{\ampere}$ e $\SI{10}{\ampere}$
(valores comerciais). O geral cobre a soma das cargas,
$\SI{27.5}{\ampere}$: disjuntor de $\SI{32}{\ampere}$.\\
(b) \textbf{Seletividade} significa que, diante de uma falta em um ramo,
\emph{apenas} o disjuntor daquele ramo abre --- os demais circuitos continuam
alimentados. Isso exige que a curva tempo--corrente do dispositivo de ramo
esteja inteiramente \emph{abaixo} da do geral.\\
Se a seletividade falhar, um curto em uma tomada derruba a instalação inteira:
além do transtorno, perde-se a informação de \emph{onde} está o defeito, e a
localização passa a exigir religamentos sucessivos.\\
(c) Porque as cargas dificilmente operam \emph{todas} no máximo ao mesmo tempo
--- é o \textbf{fator de demanda}. Aqui $25+10+10=\SI{45}{\ampere}$ contra um
geral de $\SI{32}{\ampere}$: perfeitamente admissível, desde que a demanda
simultânea real permaneça abaixo de $\SI{32}{\ampere}$.\\
\textbf{Consequência do paralelo:} como as correntes de ramo se somam no
alimentador, é o \emph{comportamento estatístico} do conjunto --- e não a soma
aritmética dos máximos --- que dimensiona a entrada. Superdimensionar o geral
para $\SI{45}{\ampere}$ exigiria condutores muito mais caros sem ganho real de
disponibilidade.}
\end{questao}

% END BANCO COMPLEMENTAR

\gabarito[3.2 Circuitos em paralelo]

% =====================================================================
\subsection{Circuitos mistos}
% =====================================================================

\blocofixacao

\begin{questao}[1]{}
A estratégia geral para reduzir um circuito misto é:
\begin{alternativas}
\item somar todas as resistências, independentemente da ligação.
\item resolver primeiro as associações série mais internas, depois os paralelos.
\item identificar os trechos em série e em paralelo e reduzi-los, dos mais internos
      para os mais externos, um de cada vez.
\item aplicar sempre a transformação estrela-triângulo.
\item calcular a corrente antes de reduzir o circuito.
\end{alternativas}
\resposta{(c)}
\resolucao{Reduz-se o circuito por etapas, sempre pelo trecho onde a ligação
(série ou paralelo) é inequívoca --- em geral o mais interno ---, substituindo-o
por um resistor equivalente e redesenhando. Repete-se até restar um único
resistor. A transformação estrela-triângulo (alternativa d) só é necessária
quando \emph{nenhum} trecho é puramente série ou paralelo, como nas pontes.}
\end{questao}

\begin{questao}[1]{}
Para cada circuito, calcule a resistência equivalente vista pela fonte e a
corrente total.

\circuitoq{%
\subcirc{a}{%
\begin{circuitikz}[scale=0.8,transform shape,line width=0.8pt]
 \draw (0,0) to[battery1,l_=$\SI{24}{\volt}$] (0,3)
       to[R,l=$\SI{6}{\ohm}$] (2.6,3) coordinate(n);
 \draw (n) to[R,l=$\SI{6}{\ohm}$] (2.6,0) -- (0,0);
 \draw (n) -- (4.6,3) to[R,l=$\SI{3}{\ohm}$] (4.6,0) -- (2.6,0);
 \draw[->,Icol,line width=1pt] (0.6,3.4)--(1.5,3.4) node[midway,above,font=\footnotesize]{$I$};
\end{circuitikz}}
\hfill
\subcirc{b}{%
\begin{circuitikz}[scale=0.8,transform shape,line width=0.8pt]
 \draw (0,0) to[battery1,l_=$\SI{20}{\volt}$] (0,3) coordinate(t);
 \draw (t) -- (2,3) to[R,l=$\SI{4}{\ohm}$] (2,0) coordinate(b);
 \draw (2,3) -- (4,3) to[R,l=$\SI{4}{\ohm}$] (4,0) -- (2,0);
 \draw (4,3) to[R,l=$\SI{8}{\ohm}$] (6,3) -- (6,0) -- (b) -- (0,0);
 \draw[->,Icol,line width=1pt] (0.5,3.4)--(1.4,3.4) node[midway,above,font=\footnotesize]{$I$};
\end{circuitikz}}}

\resposta{(a) \SI{8}{\ohm}, \SI{3}{\ampere}; (b) \SI{10}{\ohm}, \SI{2}{\ampere}}
\resolucao{\textbf{(a)} Os resistores de \SI{6}{\ohm} e \SI{3}{\ohm} da direita
estão em paralelo: $6\parallel 3 = \SI{2}{\ohm}$. Esse resultado fica em série com
o \SI{6}{\ohm} de entrada: $\Req = 6+2 = \SI{8}{\ohm}$;
$I = 24/8 = \SI{3}{\ampere}$.\\
\textbf{(b)} Os dois \SI{4}{\ohm} em paralelo dão \SI{2}{\ohm}, em série com o
\SI{8}{\ohm}: $\Req = 2+8 = \SI{10}{\ohm}$; $I = 20/10 = \SI{2}{\ampere}$.}
\end{questao}

\blocoaplicacao

\begin{questao}[2]{}
No circuito da figura, determine a resistência equivalente vista pela fonte, a
corrente total e a corrente em cada resistor de \SI{20}{\ohm}.

\circuitoq{%
\begin{circuitikz}[scale=0.9,transform shape,line width=0.8pt]
 \draw (0,0) to[battery1,l_={$E=\SI{40}{\volt}$}] (0,3)
       to[R,l=$\SI{10}{\ohm}$] (3,3) coordinate(n);
 \draw (n) to[R,l=$\SI{20}{\ohm}$] (3,0);
 \draw (n) -- (5.5,3) to[R,l=$\SI{20}{\ohm}$] (5.5,0) -- (3,0) -- (0,0);
 \draw[->,Icol,line width=1pt] (0.6,3.4)--(1.6,3.4) node[midway,above,font=\footnotesize]{$I_t$};
\end{circuitikz}}

\resposta{$\Req=\SI{20}{\ohm}$; $I_t=\SI{2}{\ampere}$; \SI{1}{\ampere} em cada \SI{20}{\ohm}}
\resolucao{Os dois \SI{20}{\ohm} em paralelo equivalem a $20/2 = \SI{10}{\ohm}$,
em série com o \SI{10}{\ohm} de entrada:
$\Req = 10+10 = \SI{20}{\ohm}$; $I_t = 40/20 = \SI{2}{\ampere}$.
Como os ramos são iguais, a corrente total se divide igualmente:
\SI{1}{\ampere} em cada resistor de \SI{20}{\ohm}.\\
\textbf{Confira as tensões:} sobre o \SI{10}{\ohm} caem $10\cdot 2 = \SI{20}{\volt}$;
sobram \SI{20}{\volt} no bloco paralelo, e $20/20 = \SI{1}{\ampere}$ em cada ramo.}
\end{questao}

\begin{questao}[2]{}
Calcule a resistência equivalente entre A e B.

\circuitoq{%
\begin{circuitikz}[scale=0.9,transform shape,line width=0.8pt]
 \draw (0,0) node[left,font=\bfseries]{A} to[short,o-*] (1.2,0) coordinate(a);
 % bloco 1: R=2 em serie
 \draw (a) to[R,l=$\SI{2}{\ohm}$] (3.2,0) coordinate(b);
 % bloco 2: 6 || 3
 \draw (b) to[R,l=$\SI{6}{\ohm}$] (5.4,0) coordinate(c);
 \draw (b) ++(0,1.4) coordinate(b2) to[R,l=$\SI{3}{\ohm}$] ($(c)+(0,1.4)$) coordinate(c2);
 \draw (b) -- (b2);  \draw (c) -- (c2);
 % bloco 3: 10 || 15
 \draw (c) to[R,l=$\SI{10}{\ohm}$] (7.6,0) coordinate(d);
 \draw (c) ++(0,1.4) coordinate(c3) to[R,l=$\SI{15}{\ohm}$] ($(d)+(0,1.4)$) coordinate(d2);
 \draw (c) -- (c3);  \draw (d) -- (d2);
 \draw (d) to[short,-o] (8.8,0) node[right,font=\bfseries]{B};
\end{circuitikz}}

\resposta{$R_{AB} = \SI{10}{\ohm}$}
\resolucao{Há dois blocos em paralelo, separados por um resistor em série:
\[
6\parallel 3=\frac{6\cdot3}{9}=\SI{2}{\ohm},
\qquad
10\parallel 15=\frac{10\cdot15}{25}=\SI{6}{\ohm}.
\]
Os três trechos estão em série (o \SI{2}{\ohm} de entrada, o primeiro bloco e o
segundo bloco):
\[
R_{AB}=2+2+6=\SI{10}{\ohm}.
\]}
\end{questao}

\blocoaprofundamento

\begin{questao}[3]{}
Determine a resistência equivalente entre A e B da rede em escada da figura.

\circuitoq{%
\begin{circuitikz}[scale=0.9,transform shape,line width=0.8pt]
 \draw (0,2) node[left,font=\bfseries]{A} to[short,o-*] (1,2) coordinate(a);
 \draw (a) to[R,l=$\SI{12}{\ohm}$] (1,0) coordinate(b);
 \draw (a) to[R,l=$\SI{4}{\ohm}$] (4,2) coordinate(c);
 \draw (c) to[R,l=$\SI{6}{\ohm}$] (4,0) coordinate(d);
 \draw (c) to[R,l=$\SI{3}{\ohm}$] (6.5,2) coordinate(e);
 \draw (e) -- (6.5,0) -- (d) -- (b);
 \draw (b) to[short,*-o] (0,0) node[left,font=\bfseries]{B};
 \draw (e) -- (7,2) node[right]{};
\end{circuitikz}}

\resposta{$R_{AB} = \SI{4}{\ohm}$}
\resolucao{Resolve-se do trecho mais distante para o mais próximo da entrada:
\begin{align*}
6 \parallel 3 &= \SI{2}{\ohm} && \text{(último ramo)}\\
4 + 2 &= \SI{6}{\ohm} && \text{(em série com o } \SI{4}{\ohm}\text{)}\\
12 \parallel 6 &= \SI{4}{\ohm} && \text{(em paralelo com o } \SI{12}{\ohm}\text{ de entrada)}.
\end{align*}
Logo $R_{AB} = \SI{4}{\ohm}$. O nome "escada" vem justamente desse padrão de
reduzir degrau a degrau.}
\end{questao}

\begin{questao}[3]{}
No circuito da figura, um fio ideal curto-circuita parte da rede. Determine a
resistência equivalente entre A e B.

\circuitoq{%
\begin{circuitikz}[scale=0.9,transform shape,line width=0.8pt]
 \draw (0,1.5) node[left,font=\bfseries]{A} to[short,o-*] (1.5,1.5) coordinate(a);
 \draw (a) to[R,l=$\SI{10}{\ohm}$] (1.5,3.2) -- (5,3.2) coordinate(top);
 \draw (a) to[R,l=$\SI{20}{\ohm}$] (5,1.5) coordinate(m);
 \draw (top) -- (m);
 \draw[Vcol,line width=1.1pt] (a) -- (1.5,0) -- (5,0) -- (m);
 \node[Vcol,font=\footnotesize\sffamily,below] at (3.25,0) {fio ideal};
 \draw (m) to[R,l=$\SI{10}{\ohm}$] (7,1.5) to[short,-o] (7.8,1.5)
       node[right,font=\bfseries]{B};
\end{circuitikz}}

\resposta{$R_{AB} = \SI{10}{\ohm}$}
\resolucao{O fio ideal liga diretamente o nó A ao nó central: eles passam a ser
\emph{o mesmo nó}. Com isso, os resistores de \SI{10}{\ohm} e \SI{20}{\ohm} do topo
ficam ambos ligados entre A e o nó central --- que agora coincidem ---, ou seja,
têm os dois terminais no mesmo ponto e ficam curto-circuitados. Sobra apenas o
resistor de \SI{10}{\ohm} até B:
\[
R_{AB} = \SI{10}{\ohm}.
\]
\textbf{Lição:} identificar nós em curto \emph{antes} de somar resistências evita
contas longas e erradas.}
\end{questao}

\begin{questao}[3]{}
Doze resistores iguais de $\SI{6}{\ohm}$ são soldados formando as arestas de um
\textbf{cubo}. Determine a resistência equivalente entre dois vértices ligados por
uma \emph{aresta} (A e B na figura).

\circuitoq{%
\begin{tikzpicture}[scale=1,line width=0.8pt]
 \def\a{2}
 % vertices do cubo em perspectiva
 \coordinate (A) at (0,0);
 \coordinate (B) at (\a,0);
 \coordinate (C) at (\a,\a);
 \coordinate (D) at (0,\a);
 \coordinate (E) at (0.7,0.6);
 \coordinate (F) at (\a+0.7,0.6);
 \coordinate (G) at (\a+0.7,\a+0.6);
 \coordinate (H) at (0.7,\a+0.6);
 \draw[Rcol] (A)--(B)--(C)--(D)--cycle;
 \draw[Rcol] (E)--(F)--(G)--(H)--cycle;
 \draw[Rcol] (A)--(E); \draw[Rcol] (B)--(F);
 \draw[Rcol] (C)--(G); \draw[Rcol] (D)--(H);
 \fill[Vcol] (A) circle (0.09) node[below left,font=\bfseries]{A};
 \fill[Vcol] (B) circle (0.09) node[below right,font=\bfseries]{B};
 \node[Rcol,font=\footnotesize] at (\a/2,-0.3){$\SI{6}{\ohm}$ (cada aresta)};
\end{tikzpicture}}{Cubo de resistores.}

\resposta{$R_{AB} = \SI{5}{\ohm}$}
\resolucao{Explorando a \emph{simetria} do cubo, injeta-se corrente $I$ em A e
retira-se em B. Por simetria, a corrente se distribui de forma previsível:
o resultado clássico para a resistência entre vértices de uma \emph{aresta} é
\[
R_{AB} = \frac{7}{12}\,R.
\]
Com $R = \SI{6}{\ohm}$: $R_{AB} = \dfrac{7}{12}\cdot6 = \SI{3.5}{\ohm}$.\\
(Observação: os três casos clássicos do cubo são $\tfrac{7}{12}R$ pela aresta,
$\tfrac{3}{4}R$ pela diagonal de face e $\tfrac{5}{6}R$ pela diagonal principal.
Com $R=6$: \SI{3.5}{}, \SI{4.5}{} e \SI{5}{\ohm}, respectivamente.) A técnica-chave
é usar a simetria para agrupar nós de mesmo potencial antes de calcular.}
\end{questao}

\begin{questao}[3]{}
Doze resistores iguais de $\SI{12}{\ohm}$ formam as arestas de um cubo.
Determine a resistência equivalente entre:
\begin{itensq}
\item dois vértices de uma mesma \emph{aresta};
\item dois vértices opostos de uma mesma \emph{face} (diagonal de face);
\item dois vértices \emph{diametralmente opostos} (diagonal principal).
\end{itensq}
\resposta{(a) \SI{7}{\ohm}; (b) \SI{9}{\ohm}; (c) \SI{10}{\ohm}}
\resolucao{A chave é explorar a \textbf{simetria} para identificar nós de mesmo
potencial, que podem ser unidos sem alterar o circuito.

\textbf{(c) Diagonal principal --- o caso mais simétrico.} Injetando corrente $I$
em A, ela se divide igualmente pelas 3 arestas que saem de A ($I/3$ em cada). Em
seguida, cada uma se divide em duas ($I/6$ nas 6 arestas intermediárias) e,
finalmente, recombina-se em 3 arestas ($I/3$) chegando em B. Somando as quedas ao
longo de um caminho:
\[
U = \frac{I}{3}R+\frac{I}{6}R+\frac{I}{3}R
= IR\left(\frac13+\frac16+\frac13\right)=\frac{5}{6}IR,
\]
logo $R_{AB}=\dfrac{5R}{6}=\dfrac{5\cdot12}{6}=\SI{10}{\ohm}$.

\textbf{(b) Diagonal de face:} por raciocínio análogo,
$R = \dfrac{3R}{4}=\dfrac{3\cdot12}{4}=\SI{9}{\ohm}$.

\textbf{(a) Aresta:} $R = \dfrac{7R}{12}=\dfrac{7\cdot12}{12}=\SI{7}{\ohm}$.

\medskip
\begin{center}\small\sffamily
\begin{tabular}{@{}l c c@{}}
\toprule
\textbf{Par de vértices} & \textbf{Fórmula} & \textbf{Valor ($R=\SI{12}{\ohm}$)}\\
\midrule
Aresta            & $\tfrac{7}{12}R$ & \SI{7}{\ohm}\\
Diagonal de face  & $\tfrac{3}{4}R$  & \SI{9}{\ohm}\\
Diagonal principal& $\tfrac{5}{6}R$  & \SI{10}{\ohm}\\
\bottomrule
\end{tabular}
\end{center}
\textbf{Note a ordenação:} quanto mais "distantes" os vértices, maior a
resistência --- mas a diferença é modesta (7 a 10 $\Omega$), porque há sempre
muitos caminhos paralelos. \emph{A simetria é a ferramenta central: sem ela,
seria preciso resolver um sistema de 7 equações nodais.}}
\end{questao}

\begin{questao}[3]{}
No circuito da figura, todos os resistores valem $\SI{6}{\ohm}$, exceto o
resistor central da ponte, de $\SI{5}{\ohm}$. Determine $R_{AB}$ e justifique por
que o valor do resistor central é irrelevante.

\circuitoq{%
\begin{circuitikz}[scale=1,transform shape,line width=0.8pt]
 \draw (0,0) node[left,font=\bfseries]{A} coordinate (A);
 \draw (5,0) node[right,font=\bfseries]{B} coordinate (B);
 \coordinate (C) at (2.5,1.6); \coordinate (D) at (2.5,-1.6);
 \draw (A) to[R,l={$\SI{4}{\ohm}$}] (C);
 \draw (C) to[R,l={$\SI{8}{\ohm}$}] (B);
 \draw (A) to[R,l_={$\SI{4}{\ohm}$}] (D);
 \draw (D) to[R,l_={$\SI{8}{\ohm}$}] (B);
 \draw (C) to[R,l={$\SI{5}{\ohm}$},color=Vcol] (D);
\end{circuitikz}}

\resposta{$R_{AB}=\SI{6}{\ohm}$; a ponte está equilibrada}
\resolucao{\textbf{Teste de equilíbrio:} os braços opostos satisfazem
\[
\frac{R_{AC}}{R_{CB}}=\frac{4}{8}=\frac12,
\qquad
\frac{R_{AD}}{R_{DB}}=\frac{4}{8}=\frac12.
\]
As razões são \emph{iguais}, logo a ponte está \textbf{equilibrada}: os nós C e D
estão no \emph{mesmo potencial}. Sem diferença de potencial, não circula corrente
pelo resistor central --- ele pode ser \emph{removido} (ou substituído por
qualquer valor, inclusive um curto ou um circuito aberto) sem alterar nada.

Restam dois ramos em série, associados em paralelo:
\[
R_{AB}=(4+8)\parallel(4+8)=\frac{12}{2}=\SI{6}{\ohm}.
\]
\textbf{Consequência prática --- a ponte como instrumento:} é exatamente essa
insensibilidade que torna a ponte de Wheatstone tão precisa. No equilíbrio, o
detector central indica zero \emph{independentemente} de sua própria resistência
e da tensão da fonte. A medição depende apenas das \emph{razões} entre
resistores conhecidos.}
\end{questao}

\begin{questao}[3]{}
Dispõe-se de um número ilimitado de resistores de $\SI{10}{\ohm}$. Projete
associações que resultem em (a) $\SI{15}{\ohm}$; (b) $\SI{25}{\ohm}$;
(c) $\SI{4}{\ohm}$, usando o \textbf{menor número possível} de resistores em cada
caso.
\resposta{(a) 3 resistores; (b) 4 resistores; (c) 4 resistores}
\resolucao{A estratégia é decompor o valor-alvo em somas de múltiplos
($nR$: série) e submúltiplos ($R/n$: paralelo) de $R = \SI{10}{\ohm}$.

\medskip
\textbf{(a) \SI{15}{\ohm} --- 3 resistores.}
\[
15 = 10 + 5 = R + \frac{R}{2}.
\]
Um resistor em série com dois em paralelo:
$10 + \dfrac{10\cdot10}{20} = 10+5 = \SI{15}{\ohm}$.

\medskip
\textbf{(b) \SI{25}{\ohm} --- 4 resistores.}
\[
25 = 20 + 5 = 2R + \frac{R}{2}.
\]
Dois em série ($\SI{20}{\ohm}$) somados a dois em paralelo ($\SI{5}{\ohm}$).

\medskip
\textbf{(c) \SI{4}{\ohm} --- 4 resistores.} Aqui a soma direta não funciona; é
preciso um \emph{paralelo de dois grupos}:
\[
\SI{4}{\ohm} = 20 \parallel 5 = \frac{20\cdot5}{25},
\]
ou seja, \textbf{dois em série} ($\SI{20}{\ohm}$) em paralelo com \textbf{dois em
paralelo} ($\SI{5}{\ohm}$). Verificando:
$\dfrac{20\cdot5}{20+5} = \dfrac{100}{25} = \SI{4}{\ohm}$.

\medskip
\textbf{Método geral.} Uma busca exaustiva sobre todas as combinações série e
paralelo confirma que 3, 4 e 4 são de fato os mínimos. O raciocínio de projeto é:
\begin{enumerate}[leftmargin=1.7em,itemsep=1pt,topsep=2pt]
\item se o alvo for \emph{maior} que $R$, comece somando (série);
\item se for \emph{menor} que $R$, comece dividindo (paralelo);
\item se não fechar, combine \emph{grupos} em paralelo --- como no item (c), onde
      $4 < 10$ exigiu paralelo de blocos, e não de resistores individuais.
\end{enumerate}
Essa habilidade é usada na prática quando um valor comercial não existe e precisa
ser sintetizado com os valores disponíveis na série E12/E24.}
\end{questao}

\blocodesafio

\begin{questao}[4]{}
A figura mostra uma rede em escada \textbf{infinita}, em que todos os resistores
valem $R$. Determine a resistência equivalente $\Req$ entre os terminais A e B.

\circuitoq{%
\begin{circuitikz}[scale=0.9,transform shape,line width=0.8pt]
 \draw (0,2) node[left,font=\bfseries]{A} to[short,o-] (0.8,2);
 \draw (0.8,2) to[R,l=$R$] (2.8,2) coordinate(a1);
 \draw (a1) to[R,l=$R$] (2.8,0) coordinate(b1);
 \draw (a1) to[R,l=$R$] (4.8,2) coordinate(a2);
 \draw (a2) to[R,l=$R$] (4.8,0) coordinate(b2);
 \draw (a2) to[R,l=$R$] (6.8,2) coordinate(a3);
 \draw (a3) to[R,l=$R$] (6.8,0) coordinate(b3);
 \draw[dotted,line width=1.2pt] (6.9,2) -- (8.0,2);
 \draw[dotted,line width=1.2pt] (6.9,0) -- (8.0,0);
 \node[font=\footnotesize] at (8.4,1) {$\cdots\infty$};
 \draw (0.8,0) to[short,o-] (0,0) node[left,font=\bfseries]{B};
 \draw (0.8,2)--(0.8,2); \draw (0.8,0)--(2.8,0);
 \draw (2.8,0)--(4.8,0); \draw (4.8,0)--(6.8,0);
\end{circuitikz}}

\resposta{$\Req = \dfrac{R(1+\sqrt5)}{2} \approx 1{,}618\,R$}
\resolucao{\textbf{Passo 1 --- explorar a autossimilaridade.} Se a rede é
infinita, retirar a primeira célula (um $R$ em série e um $R$ em paralelo) deixa
\emph{exatamente a mesma rede}. Logo, o bloco à direita do primeiro par também
vale $\Req$:
\[
\Req = R + \big(R \parallel \Req\big) = R + \frac{R\,\Req}{R+\Req}.
\]
\textbf{Passo 2 --- resolver a equação.} Multiplicando por $(R+\Req)$:
\[
\Req(R+\Req) = R(R+\Req) + R\,\Req
\;\Rightarrow\;
\Req^{\,2} - R\,\Req - R^{2} = 0.
\]
\textbf{Passo 3 --- raiz positiva.} Pela fórmula quadrática,
\[
\Req = \frac{R \pm \sqrt{R^2+4R^2}}{2} = \frac{R(1\pm\sqrt5)}{2}.
\]
Descarta-se a raiz negativa (resistência é positiva):
\[
\boxed{\ \Req = \frac{R(1+\sqrt5)}{2} \approx 1{,}618\,R\ }
\]
\textbf{Observação notável:} o fator $\dfrac{1+\sqrt5}{2}$ é a \textbf{razão
áurea} $\varphi$ --- a mesma constante da sequência de Fibonacci e da geometria
clássica. Ela aparece aqui porque a recorrência do circuito tem a mesma forma
algébrica $x = 1 + 1/x$.}
\end{questao}

\begin{questao}[4]{}
Quantas seções da escada da questão anterior são necessárias para que a
resistência de entrada difira do valor infinito em menos de
$\SI{0.1}{\percent}$? Considere $R = \SI{1}{\ohm}$ e use a recorrência
$\Req^{(n)} = R + \big(R \parallel \Req^{(n-1)}\big)$, com
$\Req^{(1)} = R$.

\begin{table}[H]
\centering\small\sffamily
\begin{tabular}{@{}c c c@{}}
\toprule
\textbf{$n$ (seções)} & \textbf{$\Req$ ($\Omega$)} & \textbf{erro relativo}\\
\midrule
1  & 1,000000 & \SI{38}{\percent}\\
2  & 1,500000 & \SI{7.3}{\percent}\\
3  & 1,600000 & \SI{1.1}{\percent}\\
5  & 1,617647 & \SI{0.024}{\percent}\\
10 & 1,618034 & $<\SI{0.001}{\percent}$\\
\bottomrule
\end{tabular}
\caption{Convergência da escada finita para o valor infinito.}
\end{table}

\resposta{$n = 5$ seções já bastam}
\resolucao{Aplicando a recorrência iterativamente com $R=\SI{1}{\ohm}$:
\[
\Req^{(1)}=1,\quad
\Req^{(2)}=1+\frac{1\cdot1}{2}=1{,}5,\quad
\Req^{(3)}=1+\frac{1{,}5}{2{,}5}=1{,}6,\ \dots
\]
Comparando com o limite $\varphi \approx 1{,}618034$, o erro cai
\emph{geometricamente}: já em $n=5$ o desvio é de \SI{0.024}{\percent}, abaixo do
critério de \SI{0.1}{\percent}.\\
\textbf{Lição de engenharia:} redes "infinitas" são uma \emph{idealização útil}.
Na prática, meia dúzia de seções já reproduz o comportamento assintótico dentro
da tolerância dos componentes reais (resistores comerciais têm
\SI{1}{\percent} a \SI{5}{\percent} de tolerância). É por isso que atenuadores e
conversores R-2R usam poucos estágios.}
\end{questao}

\begin{questao}[4]{}
Considere agora uma escada infinita em que os resistores \emph{em série} valem $R$
e os \emph{em paralelo} valem $2R$ --- a chamada rede \textbf{R-2R}, usada em
conversores digital-analógico (DAC).

\circuitoq{%
\begin{circuitikz}[scale=0.9,transform shape,line width=0.8pt]
 \draw (0,2) node[left,font=\bfseries]{A} to[short,o-] (0.8,2);
 \draw (0.8,2) to[R,l=$R$] (2.8,2) coordinate(a1);
 \draw (a1) to[R,l=$2R$,color=Rcol] (2.8,0);
 \draw (a1) to[R,l=$R$] (4.8,2) coordinate(a2);
 \draw (a2) to[R,l=$2R$,color=Rcol] (4.8,0);
 \draw (a2) to[R,l=$R$] (6.8,2) coordinate(a3);
 \draw (a3) to[R,l=$2R$,color=Rcol] (6.8,0);
 \draw[dotted,line width=1.2pt] (6.9,2) -- (8.0,2);
 \draw[dotted,line width=1.2pt] (6.9,0) -- (8.0,0);
 \node[font=\footnotesize] at (8.4,1) {$\cdots\infty$};
 \draw (0.8,0) to[short,o-] (0,0) node[left,font=\bfseries]{B};
 \draw (0.8,0)--(6.8,0);
\end{circuitikz}}

\begin{itensq}
\item Determine $\Req$ entre A e B.
\item Explique por que esse resultado torna a rede R-2R especialmente útil.
\end{itensq}
\resposta{(a) $\Req = 2R$ (exato); (b) ver comentário}
\resolucao{(a) Pela autossimilaridade,
\[
\Req = R + \big(2R \parallel \Req\big) = R + \frac{2R\,\Req}{2R+\Req}.
\]
Multiplicando por $(2R+\Req)$:
\[
\Req(2R+\Req) = R(2R+\Req) + 2R\,\Req
\;\Rightarrow\;
\Req^{\,2} - R\,\Req - 2R^{2}=0.
\]
Fatorando: $(\Req - 2R)(\Req + R) = 0$, cuja raiz positiva é
\[
\boxed{\ \Req = 2R\ }
\]
--- um valor \textbf{exato e inteiro}, sem irracionais.\\
(b) Essa é a propriedade notável da rede R-2R: a impedância vista de qualquer
ponto da escada é sempre $2R$, o que a torna \emph{perfeitamente casada} em todos
os estágios. Consequência: a corrente se divide exatamente ao meio a cada nó,
gerando os pesos binários $\tfrac12, \tfrac14, \tfrac18,\dots$ --- exatamente o
que um \textbf{conversor D/A} precisa. Além disso, o circuito usa apenas dois
valores de resistor, o que é decisivo para a precisão em circuitos integrados
(é muito mais fácil casar razões de resistores iguais do que valores absolutos).}
\end{questao}

\begin{questao}[4]{}
Generalize: para uma escada infinita com resistores $R_s$ em série e $R_p$ em
paralelo, obtenha a expressão de $\Req$ e determine a condição sobre
$R_p$ para que $\Req$ seja um múltiplo \emph{racional} de $R_s$.
\resposta{$\Req = \dfrac{R_s + \sqrt{R_s^2+4R_sR_p}}{2}$; racional se
$R_s^2+4R_sR_p$ for quadrado perfeito}
\resolucao{Da recorrência $\Req = R_s + \dfrac{R_p\Req}{R_p+\Req}$, multiplicando
e reorganizando:
\[
\Req^{\,2} - R_s\,\Req - R_s R_p = 0
\quad\Rightarrow\quad
\Req = \frac{R_s + \sqrt{R_s^{2}+4R_sR_p}}{2}.
\]
O resultado é racional quando o discriminante $\Delta = R_s^2+4R_sR_p$ é um
\emph{quadrado perfeito}. Casos particulares:
\begin{itemize}[leftmargin=1.6em,itemsep=1pt,topsep=2pt]
\item $R_p = R_s$: $\Delta = 5R_s^2 \Rightarrow \Req = \varphi R_s$ (irracional);
\item $R_p = 2R_s$: $\Delta = 9R_s^2 \Rightarrow \Req = 2R_s$ (racional);
\item $R_p = 6R_s$: $\Delta = 25R_s^2 \Rightarrow \Req = 3R_s$ (racional).
\end{itemize}
De modo geral, $R_p = \dfrac{k(k-1)}{1}\,R_s$ com $k$ inteiro fornece
$\Req = k\,R_s$. Verifique para $k=2$: $R_p = 2R_s \Rightarrow \Req = 2R_s$.
\textbf{Este é o tipo de raciocínio que separa o técnico do projetista}: em vez de
calcular um caso, obtém-se a \emph{família} de soluções e escolhe-se a que atende
ao requisito.}
\end{questao}

\begin{questao}[4]{}
A figura mostra uma rede \textbf{infinita} em que cada estágio $n$ é formado por
$(n+1)$ resistores \emph{em paralelo}, ligados entre dois barramentos
consecutivos. Os valores seguem o padrão do \textbf{triângulo de Pascal}: no
estágio $n$, os resistores valem $\dfrac{1}{\binom{n}{k}}\,\si{\ohm}$, com
$k = 0,1,\dots,n$.

\circuitoq{%
\begin{circuitikz}[scale=0.82,transform shape,line width=0.75pt,font=\footnotesize]
 % barramentos verticais
 \foreach \x/\h in {0/0.6, 2.4/1.2, 4.8/1.8, 7.2/2.4}
   \draw[cinza] (\x,-\h) -- (\x,\h);
 \draw[cinza] (9.6,-3.0) -- (9.6,3.0);
 \node[left,font=\bfseries\normalsize] at (-0.35,0) {A};
 \draw (-0.35,0)--(0,0);
 % estagio 1: 2 resistores (1,1)
 \foreach \y in {-0.55,0.55}
   \draw (0,\y) to[R,l={$1$}] (2.4,\y);
 % estagio 2: 3 resistores (1, 1/2, 1)
 \foreach \y/\lab in {-1.1/{$1$}, 0/{$\frac12$}, 1.1/{$1$}}
   \draw (2.4,\y) to[R,l={\lab}] (4.8,\y);
 % estagio 3: 4 resistores (1, 1/3, 1/3, 1)
 \foreach \y/\lab in {-1.65/{$1$}, -0.55/{$\frac13$}, 0.55/{$\frac13$}, 1.65/{$1$}}
   \draw (4.8,\y) to[R,l={\lab}] (7.2,\y);
 % estagio 4: 5 resistores (1, 1/4, 1/6, 1/4, 1)
 \foreach \y/\lab in {-2.2/{$1$}, -1.1/{$\frac14$}, 0/{$\frac16$}, 1.1/{$\frac14$}, 2.2/{$1$}}
   \draw (7.2,\y) to[R,l={\lab}] (9.6,\y);
 % reticencias
 \foreach \y in {-2.6,-1.3,0,1.3,2.6}
   \draw[dotted,line width=1pt] (9.7,\y) -- (10.9,\y);
 \node[font=\normalsize] at (11.4,0) {$\infty$};
 \node[cinza,font=\scriptsize] at (1.2,-1.7) {estágio 1};
 \node[cinza,font=\scriptsize] at (3.6,-2.2) {2};
 \node[cinza,font=\scriptsize] at (6.0,-2.6) {3};
 \node[cinza,font=\scriptsize] at (8.4,-3.0) {4};
\end{circuitikz}}

Determine a resistência equivalente $R_{A\infty}$ vista do terminal A.
\resposta{$R_{A\infty} = \SI{1}{\ohm}$ (exatamente)}
\resolucao{\textbf{Passo 1 --- reconhecer a estrutura.} Cada estágio é um
conjunto de resistores \emph{em paralelo} entre dois barramentos. Estágios
consecutivos ficam \emph{em série}. Logo:
\[
R_{A\infty}=\sum_{n=1}^{\infty} R_n,
\qquad\text{onde } R_n \text{ é o paralelo do estágio } n.
\]

\textbf{Passo 2 --- calcular cada estágio pela condutância.} Em paralelo, somam-se
as \emph{condutâncias}. Como o resistor $k$ do estágio $n$ vale
$\dfrac{1}{\binom{n}{k}}$, sua condutância é exatamente $\binom{n}{k}$:
\[
G_n=\sum_{k=0}^{n}\binom{n}{k}.
\]

\textbf{Passo 3 --- a identidade decisiva.} A soma de uma linha inteira do
triângulo de Pascal é uma potência de dois:
\[
\sum_{k=0}^{n}\binom{n}{k}=2^{\,n}
\]
(consequência do binômio de Newton com $x=y=1$: $(1+1)^n = 2^n$). Portanto
\[
G_n = 2^{\,n}
\qquad\Longrightarrow\qquad
R_n=\frac{1}{2^{\,n}}\ \si{\ohm}.
\]
Conferindo os primeiros estágios:
\begin{center}\small\sffamily
\begin{tabular}{@{}c l c c@{}}
\toprule
$n$ & \textbf{Condutâncias} & $G_n$ & $R_n$ ($\Omega$)\\
\midrule
1 & $1,1$              & $2$  & $1/2$\\
2 & $1,2,1$            & $4$  & $1/4$\\
3 & $1,3,3,1$          & $8$  & $1/8$\\
4 & $1,4,6,4,1$        & $16$ & $1/16$\\
5 & $1,5,10,10,5,1$    & $32$ & $1/32$\\
\bottomrule
\end{tabular}
\end{center}

\textbf{Passo 4 --- somar a série.} Os estágios estão em série, então basta somar:
\[
R_{A\infty}=\sum_{n=1}^{\infty}\frac{1}{2^{\,n}}
=\frac12+\frac14+\frac18+\frac{1}{16}+\cdots
\]
É uma \textbf{série geométrica} de razão $q=\tfrac12$ e primeiro termo
$a_1=\tfrac12$:
\[
R_{A\infty}=\frac{a_1}{1-q}=\frac{1/2}{1-1/2}=1.
\]
\[
\boxed{\;R_{A\infty}=\SI{1}{\ohm}\;}
\]

\textbf{Convergência.} Somando apenas os primeiros estágios:
\[
0{,}5 \to 0{,}75 \to 0{,}875 \to 0{,}9375 \to 0{,}96875 \to \cdots \to 1.
\]
Com 8 estágios já se atinge $\SI{99.6}{\percent}$ do valor final --- o erro cai
pela metade a cada estágio acrescentado.

\medskip
\textbf{Por que esta questão é notável.} Ela reúne, em um único problema:
\begin{itemize}[leftmargin=1.7em,itemsep=1pt,topsep=2pt]
\item associação \emph{paralelo} (soma de condutâncias) e \emph{série}
      (soma de resistências), nas duas direções do circuito;
\item uma identidade combinatória (soma da linha do triângulo de Pascal $=2^n$);
\item a convergência de uma série geométrica infinita.
\end{itemize}
Um circuito de complexidade infinita colapsa em uma resposta \textbf{exata e
unitária} --- resultado que só se enxerga reconhecendo o padrão, jamais
calculando estágio a estágio.}
\end{questao}

% BEGIN BANCO COMPLEMENTAR Circuitos mistos
% =====================================================================
%  Banco complementar --- Circuitos Mistos  (REVISADO)
%  Substitui a versao gerada por template (13 clones em N2, 9 em N3,
%  5 questoes conceituais indevidamente marcadas como N4).
%  Sem sobreposicao com o cap03, que ja traz: redes em escada finita e
%  infinita, R-2R, escada generalizada, cubo de resistores (duas
%  versoes), ponte com resistor central e simetria, projeto de valores
%  a partir de resistores de 10 ohm e reducao com fio ideal.
%  Sem transformacao estrela-triangulo (banco proprio).
%  Distribuicao mantida: 8 N1 + 13 N2 + 9 N3 + 5 N4 = 35
% =====================================================================

\subsubsection*{Banco complementar --- Circuitos mistos}

\begin{atencao}[Organização do banco]
Toda rede mista se resolve com o mesmo procedimento: identificar os blocos
puramente série ou puramente paralelo, reduzir \textbf{dos mais internos para
os mais externos} até chegar a $R_{eq}$, e depois \emph{expandir} de volta para
recuperar tensões e correntes de cada elemento. O N3 trata de carregamento,
falhas e saturação térmica; o N4 aborda projeto sob restrição, enumeração de
topologias, diagnóstico e redução por simetria.
\end{atencao}

\blocofixacao

\begin{questao}[1]{}
$R_1=\SI{10}{\ohm}$ está em série com o paralelo de $R_2=\SI{20}{\ohm}$ e
$R_3=\SI{20}{\ohm}$. Determine $R_{eq}$.
\resposta{$R_{eq}=\SI{20}{\ohm}$}
\resolucao{Reduza primeiro o bloco \emph{interno} (o paralelo):
\[
R_{23}=\frac{20\cdot20}{40}=\SI{10}{\ohm}.
\]
Depois o bloco externo (a série):
\[
R_{eq}=10+10=\SI{20}{\ohm}.
\]
A ordem importa: somar os três de uma vez daria $\SI{50}{\ohm}$, resultado sem
sentido físico.}
\end{questao}

\begin{questao}[1]{}
Calcule $R_{eq}$ de $(\SI{10}{\ohm}\parallel\SI{10}{\ohm})$ em série com
$(\SI{30}{\ohm}\parallel\SI{60}{\ohm})$.
\resposta{$R_{eq}=\SI{25}{\ohm}$}
\resolucao{Dois blocos paralelos independentes, depois somados:
\[
R_{a}=\frac{10}{2}=\SI{5}{\ohm},
\qquad
R_{b}=\frac{30\cdot60}{90}=\SI{20}{\ohm},
\]
\[
R_{eq}=5+20=\SI{25}{\ohm}.
\]
Blocos que não compartilham nós podem ser reduzidos em qualquer ordem entre si
--- só a hierarquia interna de cada um é obrigatória.}
\end{questao}

\begin{questao}[1]{}
Dois resistores estão ligados \emph{aos mesmos dois nós} do circuito. Eles
estão:
\begin{alternativas}[2]
\item em série
\item em paralelo
\item em série ou paralelo, conforme o desenho
\item nem em série nem em paralelo
\end{alternativas}
\resposta{(b)}
\resolucao{O critério é \textbf{topológico}, não gráfico: dois elementos ligados
ao mesmo par de nós estão submetidos à mesma tensão --- definição de paralelo.
Já a série exige que compartilhem \emph{um} nó \emph{e} que nenhum outro
elemento se conecte a ele, garantindo a mesma corrente.\\
Não se guie pelo desenho: resistores que "parecem" em série no esquema podem
estar em paralelo se um fio os ligar pelos dois lados.}
\end{questao}

\begin{questao}[1]{}
A estratégia correta para reduzir uma rede mista é:
\begin{alternativas}[2]
\item começar sempre pelos resistores mais próximos da fonte
\item reduzir dos blocos mais internos para os mais externos
\item somar todas as resistências e depois corrigir
\item reduzir os paralelos antes de qualquer série, sempre
\end{alternativas}
\resposta{(b)}
\resolucao{Reduz-se sempre \emph{de dentro para fora}, isto é, partindo do
bloco mais distante da fonte (ou mais aninhado) em direção aos terminais. Não
existe regra fixa de "paralelo antes de série": a ordem é ditada pela
hierarquia da própria rede.\\
A alternativa (d) falha, por exemplo, em $R_1+(R_2+R_3)\parallel R_4$, onde a
\emph{série} $R_2+R_3$ precisa ser feita antes do paralelo.}
\end{questao}

\begin{questao}[1]{}
Determine $R_{eq}$ de $(\SI{10}{\ohm}+\SI{20}{\ohm})$ em paralelo com
$(\SI{15}{\ohm}+\SI{15}{\ohm})$.
\resposta{$R_{eq}=\SI{15}{\ohm}$}
\resolucao{Cada ramo é reduzido primeiro:
$R_a=10+20=\SI{30}{\ohm}$ e $R_b=15+15=\SI{30}{\ohm}$. Depois o paralelo:
\[
R_{eq}=\frac{30}{2}=\SI{15}{\ohm}.
\]
Aqui a série \emph{obrigatoriamente} vem antes: só depois de reduzidos os dois
ramos é que eles ficam ligados ao mesmo par de nós.}
\end{questao}

\begin{questao}[1]{}
$R_1=\SI{12}{\ohm}$ está em série com o paralelo entre $R_2=\SI{6}{\ohm}$ e a
série $R_3=\SI{4}{\ohm}$ com $R_4=\SI{8}{\ohm}$. Determine $R_{eq}$.
\resposta{$R_{eq}=\SI{16}{\ohm}$}
\resolucao{Três níveis, do mais interno ao mais externo:
\[
R_{34}=4+8=\SI{12}{\ohm};
\qquad
R_{234}=\frac{6\cdot12}{18}=\SI{4}{\ohm};
\qquad
R_{eq}=12+4=\SI{16}{\ohm}.
\]}
\end{questao}

\begin{questao}[1]{}
Uma fonte de $\SI{24}{\volt}$ alimenta a rede
$\SI{10}{\ohm}+(\SI{20}{\ohm}\parallel\SI{20}{\ohm})$. Determine a corrente
total e a tensão sobre o bloco paralelo.
\resposta{$I=\SI{1.2}{\ampere}$; $U_{par}=\SI{12}{\volt}$}
\resolucao{Já se sabe que $R_{eq}=\SI{20}{\ohm}$, logo
\[
I=\frac{24}{20}=\SI{1.2}{\ampere}.
\]
Toda a corrente atravessa $R_1$, e depois o bloco paralelo:
\[
U_{par}=R_{23}\,I=10\cdot1{,}2=\SI{12}{\volt}.
\]
\textbf{Conferência (LKT):} $U_1=10\cdot1{,}2=\SI{12}{\volt}$, e
$12+12=\SI{24}{\volt}$ \checkmark.}
\end{questao}

\begin{questao}[1]{}
Em uma rede mista, a corrente que atravessa o resistor imediatamente ligado ao
terminal da fonte é:
\begin{alternativas}[2]
\item igual à corrente total
\item metade da corrente total
\item dividida entre os ramos internos
\item indeterminada sem mais dados
\end{alternativas}
\resposta{(a)}
\resolucao{Se esse resistor está \emph{em série} com a fonte --- isto é, se não
há bifurcação antes dele --- toda a corrente que sai da fonte passa por ele.
Por isso ele é o primeiro a ser recuperado na etapa de expansão, e costuma ser
também o mais solicitado termicamente.\\
Note a ressalva: se houver um nó entre a fonte e o resistor, a alternativa (c)
passaria a valer. O critério é sempre topológico.}
\end{questao}

\blocoaplicacao

\begin{questao}[2]{}
Uma fonte de $\SI{24}{\volt}$ alimenta $R_1=\SI{10}{\ohm}$ em série com
$R_2=\SI{20}{\ohm}$ e $R_3=\SI{20}{\ohm}$ em paralelo. Determine todas as
correntes e tensões.
\resposta{$I_t=\SI{1.2}{\ampere}$; $U_1=\SI{12}{\volt}$;
$U_{23}=\SI{12}{\volt}$; $I_2=I_3=\SI{0.6}{\ampere}$}
\resolucao{\textbf{Redução:} $R_{23}=\SI{10}{\ohm}$, $R_{eq}=\SI{20}{\ohm}$,
$I_t=24/20=\SI{1.2}{\ampere}$.\\
\textbf{Expansão:} $U_1=10(1{,}2)=\SI{12}{\volt}$;
$U_{23}=24-12=\SI{12}{\volt}$;
\[
I_2=I_3=\frac{12}{20}=\SI{0.6}{\ampere}.
\]
\textbf{Verificação (LKC):} $0{,}6+0{,}6=\SI{1.2}{\ampere}$ \checkmark.
Como $R_2=R_3$, a corrente se divide igualmente --- resultado que a simetria
antecipa sem nenhum cálculo.}
\end{questao}

\begin{questao}[2]{}
No circuito anterior, $R_3$ é trocado por $\SI{60}{\ohm}$. Recalcule
$R_{eq}$, a corrente total e as correntes de ramo.
\resposta{$R_{eq}=\SI{25}{\ohm}$; $I_t=\SI{0.96}{\ampere}$;
$I_2=\SI{0.72}{\ampere}$, $I_3=\SI{0.24}{\ampere}$}
\resolucao{$R_{23}=\dfrac{20\cdot60}{80}=\SI{15}{\ohm}$ e
$R_{eq}=10+15=\SI{25}{\ohm}$, logo $I_t=24/25=\SI{0.96}{\ampere}$.\\
$U_{23}=15(0{,}96)=\SI{14.4}{\volt}$, e daí
\[
I_2=\frac{14{,}4}{20}=\SI{0.72}{\ampere},
\qquad
I_3=\frac{14{,}4}{60}=\SI{0.24}{\ampere}.
\]
\textbf{Divisor de corrente (alternativa):}
$I_2=0{,}96\cdot\frac{60}{80}=\SI{0.72}{\ampere}$ \checkmark. Note que a razão
$I_2/I_3=3$ é exatamente $R_3/R_2$ --- inversa das resistências.}
\end{questao}

\begin{questao}[2]{}
Na rede $\SI{10}{\ohm}+(\SI{20}{\ohm}\parallel\SI{20}{\ohm})$ sob
$\SI{24}{\volt}$, calcule a potência de cada resistor e confronte com a
potência da fonte.
\resposta{$\SI{14.4}{\watt}$; $\SI{7.2}{\watt}$; $\SI{7.2}{\watt}$; total
$\SI{28.8}{\watt}$}
\resolucao{
\[
P_1=R_1I_t^{2}=10(1{,}2)^{2}=\SI{14.4}{\watt};
\qquad
P_2=P_3=\frac{U_{23}^{2}}{20}=\frac{144}{20}=\SI{7.2}{\watt}.
\]
Fonte: $P=24(1{,}2)=\SI{28.8}{\watt}=14{,}4+7{,}2+7{,}2$ \checkmark.\\
\textbf{Observação:} $R_1$ dissipa exatamente metade de toda a potência, embora
seja o menor resistor da rede --- porque é o único que conduz a corrente
inteira. Em redes mistas, "quem aquece mais" depende da posição, não apenas do
valor.}
\end{questao}

\begin{questao}[2]{}
Determine $R_{eq}$ de $(\SI{5}{\ohm}+\SI{15}{\ohm})$ em paralelo com
$(\SI{10}{\ohm}+\SI{10}{\ohm})$, e a corrente de cada ramo sob
$\SI{30}{\volt}$.
\resposta{$R_{eq}=\SI{10}{\ohm}$; $I_t=\SI{3}{\ampere}$; $\SI{1.5}{\ampere}$
em cada ramo}
\resolucao{Ramos: $R_a=\SI{20}{\ohm}$ e $R_b=\SI{20}{\ohm}$, logo
$R_{eq}=\SI{10}{\ohm}$ e $I_t=30/10=\SI{3}{\ampere}$.\\
Ambos os ramos estão sob os mesmos $\SI{30}{\volt}$:
$I_a=I_b=30/20=\SI{1.5}{\ampere}$.\\
\textbf{Cuidado:} apesar de $R_a=R_b$, as \emph{quedas internas} diferem. No
ramo $a$: $\SI{7.5}{\volt}$ e $\SI{22.5}{\volt}$; no ramo $b$:
$\SI{15}{\volt}$ e $\SI{15}{\volt}$. Ramos equivalentes nos terminais não são
equivalentes internamente.}
\end{questao}

\begin{questao}[2]{}
Uma rede é formada por $R_1=\SI{8}{\ohm}$ em série com o paralelo de
$\SI{24}{\ohm}$ e $R$. Sabendo que $R_{eq}=\SI{24}{\ohm}$, determine $R$.
\resposta{$R=\SI{48}{\ohm}$}
\resolucao{O bloco paralelo deve valer $24-8=\SI{16}{\ohm}$. Em condutâncias:
\[
\frac{1}{R}=\frac{1}{16}-\frac{1}{24}=\frac{3-2}{48}
\;\Longrightarrow\;
R=\SI{48}{\ohm}.
\]
\textbf{Conferência:} $24\parallel48=\dfrac{24\cdot48}{72}=\SI{16}{\ohm}$ e
$8+16=\SI{24}{\ohm}$ \checkmark.\\
Note que a subtração foi feita em \emph{condutâncias}: tentar subtrair
resistências ($16-24$) daria valor negativo, sinal claro de método errado.}
\end{questao}

\begin{questao}[2]{}
Determine $R_{eq}$ de $R_1=\SI{10}{\ohm}$ em série com o paralelo entre
$(\SI{20}{\ohm}+\SI{20}{\ohm})$ e $\SI{10}{\ohm}$.
\resposta{$R_{eq}=\SI{18}{\ohm}$}
\resolucao{
\[
R_{23}=20+20=\SI{40}{\ohm};
\qquad
R_{234}=\frac{40\cdot10}{50}=\SI{8}{\ohm};
\qquad
R_{eq}=10+8=\SI{18}{\ohm}.
\]
O bloco paralelo resultou em $\SI{8}{\ohm}$ --- menor que o menor dos dois
($\SI{10}{\ohm}$), como sempre deve ocorrer. Esse teste de sanidade a cada
etapa evita que um erro se propague por toda a redução.}
\end{questao}

\begin{questao}[2]{}
Na rede $\SI{10}{\ohm}+(\SI{20}{\ohm}\parallel\SI{20}{\ohm})$ sob
$\SI{24}{\volt}$, o nó entre $R_1$ e o bloco paralelo é chamado $A$, e o
terminal negativo da fonte é o terra. Determine $V_A$.
\resposta{$V_A=\SI{12}{\volt}$}
\resolucao{O potencial de $A$ em relação ao terra é a tensão sobre o bloco
paralelo, pois é ele que separa $A$ do terra:
\[
V_A=U_{23}=\SI{12}{\volt}.
\]
\textbf{Alternativamente,} descendo da fonte: $V_A=24-U_1=24-12=\SI{12}{\volt}$
\checkmark.\\
Os dois caminhos precisam concordar --- é a LKT em ação. Se discordassem,
haveria erro na redução.}
\end{questao}

\begin{questao}[2]{}
Na mesma rede, um dos resistores de $\SI{20}{\ohm}$ sofre
\textbf{curto-circuito}. Recalcule $R_{eq}$, a corrente total e a potência de
$R_1$.
\resposta{$R_{eq}=\SI{10}{\ohm}$; $I_t=\SI{2.4}{\ampere}$;
$P_1=\SI{57.6}{\watt}$}
\resolucao{Um curto no bloco paralelo anula o bloco inteiro: o outro resistor
de $\SI{20}{\ohm}$ fica em paralelo com $\SI{0}{\ohm}$ e deixa de conduzir.
\[
R_{eq}=10+0=\SI{10}{\ohm},
\qquad
I_t=\frac{24}{10}=\SI{2.4}{\ampere},
\qquad
P_1=10(2{,}4)^{2}=\SI{57.6}{\watt}.
\]
$P_1$ saltou de $14{,}4$ para $\SI{57.6}{\watt}$ --- fator \textbf{quatro},
pois a corrente dobrou. Se $R_1$ fosse especificado para
$\SI{25}{\watt}$, queimaria em seguida.\\
\textbf{Padrão de falha:} o curto em um ramo interno transfere toda a
solicitação para o elemento em série --- e o defeito "migra" da causa para a
vítima, dificultando o diagnóstico.}
\end{questao}

\begin{questao}[2]{}
Na mesma rede, agora um dos resistores de $\SI{20}{\ohm}$ \textbf{abre}.
Recalcule $R_{eq}$, a corrente total e a potência do resistor remanescente.
\resposta{$R_{eq}=\SI{30}{\ohm}$; $I_t=\SI{0.8}{\ampere}$;
$P=\SI{12.8}{\watt}$}
\resolucao{Sem um dos ramos, o bloco paralelo passa a valer $\SI{20}{\ohm}$:
\[
R_{eq}=10+20=\SI{30}{\ohm},
\qquad
I_t=\frac{24}{30}=\SI{0.8}{\ampere},
\qquad
U_{par}=20(0{,}8)=\SI{16}{\volt},
\]
\[
P=\frac{16^{2}}{20}=\SI{12.8}{\watt}.
\]
\textbf{Contraste com o curto:} aqui a corrente total \emph{caiu} (de $1{,}2$
para $\SI{0.8}{\ampere}$), mas o resistor sobrevivente passou a dissipar
$12{,}8$ em vez de $\SI{7.2}{\watt}$ --- $78\%$ mais, porque agora assume
sozinho a corrente antes dividida. Falhas de abertura em blocos paralelos
\emph{sobrecarregam} os ramos irmãos.}
\end{questao}

\begin{questao}[2]{}
Duas redes têm $R_{eq}=\SI{20}{\ohm}$: (A) $\SI{10}{\ohm}$ em série com
$\SI{20}{\ohm}\parallel\SI{20}{\ohm}$; (B) $\SI{30}{\ohm}\parallel\SI{60}{\ohm}$.
Sob $\SI{60}{\volt}$, compare a potência total e a do componente mais
solicitado.
\resposta{Total $\SI{180}{\watt}$ em ambas; pior componente:
$\SI{90}{\watt}$ em (A) e $\SI{120}{\watt}$ em (B)}
\resolucao{$I_t=60/20=\SI{3}{\ampere}$ e $P_{tot}=60(3)=\SI{180}{\watt}$ nas
duas.\\
\textbf{(A)} $P_{10}=10(3)^{2}=\SI{90}{\watt}$; bloco sob
$\SI{30}{\volt}$, cada $\SI{20}{\ohm}$ com $\SI{45}{\watt}$. Soma:
$90+45+45=\SI{180}{\watt}$ \checkmark.\\
\textbf{(B)} $P_{30}=\dfrac{3600}{30}=\SI{120}{\watt}$;
$P_{60}=\dfrac{3600}{60}=\SI{60}{\watt}$. Soma $\SI{180}{\watt}$ \checkmark.\\
\textbf{Conclusão:} resistência equivalente idêntica, potência total idêntica,
mas o ponto quente difere em $33\%$. A escolha da topologia é decisão térmica,
não elétrica.}
\end{questao}

\begin{questao}[2]{}
Uma fonte de $\EMF=\SI{24}{\volt}$ e $r=\SI{2}{\ohm}$ alimenta a rede mista de
$\SI{10}{\ohm}$. Determine a corrente, a tensão nos terminais e a potência
entregue à rede.
\resposta{$I=\SI{2}{\ampere}$; $V=\SI{20}{\volt}$; $P=\SI{40}{\watt}$}
\resolucao{
\[
I=\frac{24}{2+10}=\SI{2}{\ampere};
\qquad
V=24-2(2)=\SI{20}{\volt};
\qquad
P_{rede}=20(2)=\SI{40}{\watt}.
\]
A fonte gera $24(2)=\SI{48}{\watt}$, dos quais $2(2)^{2}=\SI{8}{\watt}$ se
perdem internamente. Rendimento: $40/48=83{,}3\%$.\\
\textbf{Regra útil:} o rendimento é $R_{rede}/(R_{rede}+r)$ --- aqui
$10/12$. Ele depende apenas dessa razão, não da topologia interna da rede.}
\end{questao}

\begin{questao}[2]{}
Em uma rede mista de $R_{eq}=\SI{16}{\ohm}$ sob $\SI{32}{\volt}$, verificou-se
$P_{total}=\SI{64}{\watt}$. A medição é consistente? Justifique de duas
maneiras.
\resposta{Sim: $P=V^{2}/R_{eq}=VI=\SI{64}{\watt}$}
\resolucao{\textbf{Primeiro caminho:}
$P=\dfrac{V^{2}}{R_{eq}}=\dfrac{32^{2}}{16}=\SI{64}{\watt}$ \checkmark.\\
\textbf{Segundo caminho:} $I=32/16=\SI{2}{\ampere}$ e
$P=VI=32(2)=\SI{64}{\watt}$ \checkmark.\\
\textbf{Terceiro (se os elementos forem conhecidos):} somar $R_kI_k^{2}$ de
cada resistor deve devolver os mesmos $\SI{64}{\watt}$. Essa terceira
verificação é a mais forte, porque testa a \emph{expansão} da rede --- as duas
primeiras só testam $R_{eq}$ e passariam mesmo se a distribuição interna
estivesse errada.}
\end{questao}

\begin{questao}[2]{}
Na rede $R_1=\SI{12}{\ohm}$ em série com $R_2=\SI{6}{\ohm}\parallel
(R_3=\SI{4}{\ohm}+R_4=\SI{8}{\ohm})$, sob $\SI{32}{\volt}$, determine a
corrente em $R_3$.
\resposta{$I_3=\SI{0.667}{\ampere}$}
\resolucao{\textbf{Redução:} $R_{34}=\SI{12}{\ohm}$;
$R_{234}=6\parallel12=\SI{4}{\ohm}$; $R_{eq}=12+4=\SI{16}{\ohm}$.\\
\textbf{Expansão:} $I_t=32/16=\SI{2}{\ampere}$;
$U_{234}=4(2)=\SI{8}{\volt}$.\\
Esse é o mesmo $\SI{8}{\volt}$ aplicado ao ramo $R_3+R_4$:
\[
I_3=\frac{8}{12}=\SI{0.667}{\ampere}.
\]
\textbf{Conferência pelo divisor:}
$I_3=I_t\cdot\dfrac{6}{6+12}=2\cdot\dfrac13=\SI{0.667}{\ampere}$ \checkmark. E
$I_2=8/6=\SI{1.333}{\ampere}$, com $0{,}667+1{,}333=\SI{2}{\ampere}$
\checkmark.}
\end{questao}

\blocoaprofundamento

\begin{questao}[3]{}
Uma rede consiste em $(R_1+R_2)$ em paralelo com $(R_3+R_4)$.
\begin{itensq}
\item Obtenha a expressão de $R_{eq}$.
\item Avalie para $R_1=\SI{8}{\ohm}$, $R_2=\SI{16}{\ohm}$,
      $R_3=R_4=\SI{12}{\ohm}$.
\item Analise o caso-limite $R_3+R_4\gg R_1+R_2$.
\end{itensq}
\resposta{(a) $R_{eq}=\dfrac{(R_1+R_2)(R_3+R_4)}{R_1+R_2+R_3+R_4}$;
(b) $\SI{12}{\ohm}$; (c) $R_{eq}\to R_1+R_2$}
\resolucao{(a) Reduzindo cada ramo e aplicando produto pela soma:
\[
R_{eq}=\frac{(R_1+R_2)(R_3+R_4)}{(R_1+R_2)+(R_3+R_4)}.
\]
(b) $R_a=\SI{24}{\ohm}$, $R_b=\SI{24}{\ohm}$, logo
$R_{eq}=\dfrac{24\cdot24}{48}=\SI{12}{\ohm}$.\\
(c) Se $R_b\gg R_a$, divida numerador e denominador por $R_b$:
\[
R_{eq}=\frac{R_a}{\dfrac{R_a}{R_b}+1}\longrightarrow R_a=R_1+R_2.
\]
O ramo de alta resistência comporta-se como circuito aberto e não participa.
\textbf{Uso prático:} é essa aproximação que justifica desprezar a resistência
de um voltímetro ligado em paralelo a um trecho de circuito --- desde que ela
seja, digamos, cem vezes maior que o trecho medido.}
\end{questao}

\begin{questao}[3]{}
Um divisor formado por $R_1=\SI{140}{\ohm}$ e $R_2=\SI{100}{\ohm}$ é alimentado
por $\SI{12}{\volt}$.
\begin{itensq}
\item Determine a tensão de saída em vazio.
\item Recalcule com uma carga de $\SI{1}{\kilo\ohm}$ conectada à saída.
\item Determine o erro percentual introduzido pela carga.
\end{itensq}
\resposta{(a) $\SI{5.00}{\volt}$; (b) $\SI{4.72}{\volt}$; (c) $-5{,}5\%$}
\resolucao{(a) $U_{2}=12\cdot\dfrac{100}{240}=\SI{5.00}{\volt}$.\\
(b) A carga entra em \textbf{paralelo} com $R_2$:
\[
R_2'=\frac{100\cdot1000}{1100}=\SI{90.91}{\ohm},
\qquad
U_2=12\cdot\frac{90{,}91}{140+90{,}91}=\SI{4.72}{\volt}.
\]
(c) Erro $=(4{,}72-5{,}00)/5{,}00=-5{,}5\%$.\\
\textbf{Diagnóstico:} a carga de $\SI{1}{\kilo\ohm}$ é apenas $10\times$ maior
que $R_2$ --- insuficiente. A regra de bolso exige
$R_L\ge10R_2$ para erro em torno de $5\%$, e $R_L\ge100R_2$ para erro em torno
de $0{,}5\%$. Um divisor resistivo não é fonte de tensão: ele tem resistência
interna $R_1\parallel R_2=\SI{58.3}{\ohm}$, e é ela que produz a queda ao
ser carregada. O projeto para uma carga \emph{conhecida} é feito no bloco de
desafio.}
\end{questao}

\begin{questao}[3]{}
Uma rede tem $R_1=\SI{6}{\ohm}$ em série com $(\SI{12}{\ohm}\parallel
\SI{12}{\ohm})$, e este em série com $(\SI{8}{\ohm}\parallel\SI{8}{\ohm})$, sob
$\SI{32}{\volt}$. Determine a potência de cada um dos cinco resistores.
\resposta{$\SI{24}{\watt}$ em $R_1$; $\SI{12}{\watt}$ em cada
$\SI{12}{\ohm}$; $\SI{8}{\watt}$ em cada $\SI{8}{\ohm}$}
\resolucao{\textbf{Redução:} $12\parallel12=\SI{6}{\ohm}$;
$8\parallel8=\SI{4}{\ohm}$; $R_{eq}=6+6+4=\SI{16}{\ohm}$ e
$I_t=32/16=\SI{2}{\ampere}$.\\
\textbf{Expansão:}
\begin{itemize}
\item $R_1$: conduz $\SI{2}{\ampere}$ $\Rightarrow$
      $P=6(4)=\SI{24}{\watt}$.
\item Bloco de $\SI{12}{\ohm}$: $U=6(2)=\SI{12}{\volt}$; cada ramo com
      $\SI{1}{\ampere}$ $\Rightarrow$ $P=12(1)=\SI{12}{\watt}$.
\item Bloco de $\SI{8}{\ohm}$: $U=4(2)=\SI{8}{\volt}$; cada ramo com
      $\SI{1}{\ampere}$ $\Rightarrow$ $P=8(1)=\SI{8}{\watt}$.
\end{itemize}
\textbf{Balanço:} $24+12+12+8+8=\SI{64}{\watt}=32(2)$ \checkmark.\\
Observe que $R_1$, sendo o menor resistor da rede, é o que mais dissipa ---
porque conduz o dobro da corrente dos demais e a potência vai com $I^{2}$.}
\end{questao}

\begin{questao}[3]{}
Um reostato de $\SI{0}{\ohm}$ a $\SI{60}{\ohm}$ está em paralelo com
$\SI{30}{\ohm}$, e o conjunto em série com $\SI{20}{\ohm}$, sob
$\SI{60}{\volt}$.
\begin{itensq}
\item Determine a faixa de $R_{eq}$.
\item Determine a faixa da corrente total.
\item Explique por que a faixa é limitada.
\end{itensq}
\resposta{(a) de $\SI{20}{\ohm}$ a $\SI{40}{\ohm}$; (b) de $\SI{1.5}{\ampere}$
a $\SI{3}{\ampere}$}
\resolucao{(a) Com o reostato em $\SI{0}{\ohm}$, o bloco paralelo é
curto-circuitado e vale $\SI{0}{\ohm}$: $R_{eq}=\SI{20}{\ohm}$.\\
Com o reostato em $\SI{60}{\ohm}$: $30\parallel60=\SI{20}{\ohm}$ e
$R_{eq}=\SI{40}{\ohm}$.\\
(b) $I_{\max}=60/20=\SI{3}{\ampere}$ e
$I_{\min}=60/40=\SI{1.5}{\ampere}$: faixa de $2{:}1$.\\
(c) Dois limites atuam. O resistor \emph{série} de $\SI{20}{\ohm}$ impõe um
teto à corrente ($\SI{3}{\ampere}$), por mais que o bloco paralelo seja
anulado. E o resistor de $\SI{30}{\ohm}$ \emph{em paralelo} com o reostato
impõe um piso a $R_{eq}$: mesmo com o reostato em $\infty$, o bloco não
passaria de $\SI{30}{\ohm}$, e $R_{eq}$ de $\SI{50}{\ohm}$.\\
\textbf{Consequência de projeto:} para ampliar a faixa de controle, é preciso
reduzir o resistor série \emph{e} aumentar o resistor em paralelo --- alterar
apenas o reostato tem efeito saturante.}
\end{questao}

\begin{questao}[3]{}
Na rede $\SI{10}{\ohm}+(\SI{20}{\ohm}\parallel\SI{20}{\ohm})$, todos os
resistores têm potência nominal de $\SI{10}{\watt}$. Determine a tensão máxima
aplicável e identifique o componente crítico.
\resposta{$V_{\max}=\SI{20}{\volt}$; o crítico é $R_1$}
\resolucao{Expresse cada potência em função da corrente total $I$:
\[
P_1=10I^{2};
\qquad
P_2=P_3=20\left(\frac{I}{2}\right)^{2}=5I^{2}.
\]
$R_1$ dissipa \textbf{o dobro} de cada ramo paralelo --- é ele que satura
primeiro:
\[
10I^{2}=10 \;\Longrightarrow\; I=\SI{1}{\ampere}
\;\Longrightarrow\;
V_{\max}=R_{eq}I=20(1)=\SI{20}{\volt}.
\]
Nessa condição os ramos paralelos operam a $\SI{5}{\watt}$ --- metade do
nominal, ou seja, subutilizados.\\
\textbf{Melhoria possível:} trocar $R_1$ por dois de $\SI{20}{\ohm}$ em
paralelo (mesmos $\SI{10}{\ohm}$, mas $\SI{20}{\watt}$ de capacidade) elevaria
$V_{\max}$ para o limite imposto pelos ramos:
$5I^{2}=10\Rightarrow I=\SI{1.41}{\ampere}$ e
$V_{\max}=\SI{28.3}{\volt}$. Identificar o componente crítico é o primeiro
passo para melhorar qualquer rede.}
\end{questao}

\begin{questao}[3]{}
Obtenha $\SI{24}{\ohm}$ dispondo apenas de resistores de $\SI{36}{\ohm}$.
Determine a associação de \textbf{menor} número de unidades e prove que não há
solução com duas.
\resposta{Três unidades: $(\SI{36}{\ohm}+\SI{36}{\ohm})\parallel\SI{36}{\ohm}$}
\resolucao{\textbf{Com duas unidades} só existem dois valores possíveis:
$36+36=\SI{72}{\ohm}$ e $36\parallel36=\SI{18}{\ohm}$. Nenhum é
$\SI{24}{\ohm}$ --- logo duas são insuficientes.\\
\textbf{Com três:}
\[
(36+36)\parallel36=\frac{72\cdot36}{108}=\SI{24}{\ohm}.\ \checkmark
\]
\textbf{Verificação por outro caminho:} três em paralelo dariam
$\SI{12}{\ohm}$; três em série, $\SI{108}{\ohm}$; e
$36+(36\parallel36)=\SI{54}{\ohm}$. Das quatro topologias possíveis com três
unidades, apenas a indicada resolve.\\
\textbf{Observação:} $24=\frac23\cdot36$, e a razão $2/3$ é característica da
combinação "dois em série contra um em paralelo" --- padrão útil de memorizar,
pois $\dfrac{2R\cdot R}{3R}=\dfrac{2R}{3}$ para qualquer $R$.}
\end{questao}

\begin{questao}[3]{}
Em uma rede mista alimentada por fonte ideal, um resistor de um ramo paralelo
abre. Descreva um procedimento sistemático para localizar a falha usando apenas
um voltímetro, e aplique-o à rede
$\SI{10}{\ohm}+(\SI{20}{\ohm}\parallel\SI{20}{\ohm})$ sob $\SI{24}{\volt}$.
\resposta{Comparar as tensões medidas com as previstas; o bloco defeituoso
apresenta tensão \emph{maior} que a esperada}
\resolucao{\textbf{Procedimento.} \emph{(i)} Calcule as tensões esperadas em
operação normal. \emph{(ii)} Meça, bloco a bloco, a partir da fonte.
\emph{(iii)} Compare: um aumento de tensão indica bloco com resistência
\emph{maior} que a prevista (abertura); uma queda indica resistência
\emph{menor} (curto).\\
\textbf{Aplicação.} Previsto: $U_1=\SI{12}{\volt}$ e
$U_{par}=\SI{12}{\volt}$. Com um ramo aberto,
$R_{eq}=\SI{30}{\ohm}$, $I=\SI{0.8}{\ampere}$ e:
\[
U_1=10(0{,}8)=\SI{8}{\volt}\ (\downarrow),
\qquad
U_{par}=20(0{,}8)=\SI{16}{\volt}\ (\uparrow).
\]
O voltímetro acusa $\SI{16}{\volt}$ onde se esperavam $\SI{12}{\volt}$ ---
o bloco paralelo é o suspeito.\\
\textbf{Confirmação:} isolando o bloco, um ohmímetro leria
$\SI{20}{\ohm}$ em vez de $\SI{10}{\ohm}$. E \emph{qual} dos dois ramos
abriu se identifica medindo cada um separadamente --- a tensão sozinha não
distingue ramos idênticos em paralelo.\\
\textbf{Por que $U_1$ cai:} a corrente total diminuiu. Note a assinatura
característica: em uma abertura, as tensões se \emph{redistribuem} com soma
constante ($8+16=24$), enquanto em um curto elas colapsam localmente.}
\end{questao}

\begin{questao}[3]{}
Uma rede mista de $R_{eq}=\SI{20}{\ohm}$ é alimentada por fonte de
$\EMF=\SI{60}{\volt}$ com $r=\SI{5}{\ohm}$.
\begin{itensq}
\item Determine a corrente, a tensão na rede e o rendimento.
\item Determine o valor de $R_{eq}$ que maximizaria a potência na rede.
\item Compare o rendimento nas duas situações.
\end{itensq}
\resposta{(a) $\SI{2.4}{\ampere}$, $\SI{48}{\volt}$, $80\%$;
(b) $R_{eq}=\SI{5}{\ohm}$; (c) $80\%$ contra $50\%$}
\resolucao{(a) $I=\dfrac{60}{25}=\SI{2.4}{\ampere}$;
$V=20(2{,}4)=\SI{48}{\volt}$;
$P_{rede}=48(2{,}4)=\SI{115.2}{\watt}$; rendimento
$=20/25=80\%$.\\
(b) A potência na rede é máxima quando $R_{eq}=r=\SI{5}{\ohm}$. Então
$I=60/10=\SI{6}{\ampere}$, $V=\SI{30}{\volt}$ e
$P=\SI{180}{\watt}$.\\
(c) Rendimento na condição de casamento: $5/10=50\%$. A fonte passa a
gerar $\SI{360}{\watt}$, dos quais $\SI{180}{\watt}$ se perdem internamente.\\
\textbf{Conflito de objetivos:} entregar \emph{mais} potência
($180$ contra $\SI{115.2}{\watt}$) custa despachar $\SI{360}{\watt}$ em vez de
$\SI{144}{\watt}$ --- e aquecer a fonte com $\SI{180}{\watt}$. Em eletrônica de
sinal, o casamento vale a pena; em eletrotécnica de potência, jamais.}
\end{questao}

\begin{questao}[3]{}
Uma fonte de \textbf{corrente} de $\SI{3}{\ampere}$ alimenta $R_1=\SI{20}{\ohm}$
em paralelo com o ramo formado por $R_2=\SI{10}{\ohm}$ em série com
$(\SI{20}{\ohm}\parallel\SI{20}{\ohm})$.
\begin{itensq}
\item Determine $R_{eq}$ e a tensão da associação.
\item Obtenha a corrente em cada um dos cinco resistores.
\item Verifique o balanço de potência.
\end{itensq}
\resposta{(a) $\SI{10}{\ohm}$ e $\SI{30}{\volt}$;
(b) $\SI{1.5}{\ampere}$ em $R_1$ e em $R_2$, $\SI{0.75}{\ampere}$ em cada
$\SI{20}{\ohm}$ do bloco; (c) $\SI{90}{\watt}$}
\resolucao{(a) O ramo vale $10+(20\parallel20)=10+10=\SI{20}{\ohm}$, e
\[
R_{eq}=20\parallel20=\SI{10}{\ohm}
\;\Longrightarrow\;
U=I_tR_{eq}=3(10)=\SI{30}{\volt}.
\]
\textbf{Diferença essencial em relação à fonte de tensão:} aqui $U$ é
\emph{resultado}, não dado. Se a rede mudar, a tensão muda --- a corrente é que
permanece fixa.\\
(b) Os dois ramos estão sob $\SI{30}{\volt}$ e têm $\SI{20}{\ohm}$ cada:
$I_{R_1}=I_{ramo}=\SI{1.5}{\ampere}$. Dentro do ramo,
$U_{R_2}=10(1{,}5)=\SI{15}{\volt}$ e restam $\SI{15}{\volt}$ para o bloco
paralelo:
\[
I=\frac{15}{20}=\SI{0.75}{\ampere}\ \text{em cada}.
\]
\textbf{LKC:} $0{,}75+0{,}75=\SI{1.5}{\ampere}$ \checkmark; e
$1{,}5+1{,}5=\SI{3}{\ampere}$ \checkmark.\\
(c) $P_1=\dfrac{30^{2}}{20}=45$; $P_2=10(1{,}5)^{2}=22{,}5$; e
$20(0{,}75)^{2}=\SI{11.25}{\watt}$ em cada resistor do bloco:
\[
45+22{,}5+11{,}25+11{,}25=\SI{90}{\watt}=30(3)\ \checkmark
\]
\textbf{Observação:} a potência da fonte de corrente só pôde ser calculada
\emph{depois} de obtida a tensão. Com fonte de tensão ocorre o inverso --- é a
corrente que precisa ser encontrada primeiro.}
\end{questao}

\blocodesafio

\begin{questao}[4]{}
\textbf{(Projeto de divisor carregado.)} Deseja-se obter $\SI{5}{\volt}$ a
partir de $\SI{12}{\volt}$ para alimentar uma carga fixa de
$\SI{1}{\kilo\ohm}$.
\begin{itensq}
\item Deduza o critério de "corrente de sangria" e determine $R_1,R_2$ para
      erro inferior a $2\%$.
\item Projete o divisor \emph{compensado}, que entregue exatamente
      $\SI{5}{\volt}$ com a carga presente.
\item Compare as duas soluções em potência dissipada e em robustez.
\end{itensq}
\resposta{(a) $R_2\approx\SI{25}{\ohm}$, $R_1=\SI{35}{\ohm}$ (erro $1{,}4\%$);
(b) $R_1=\SI{140}{\ohm}$, $R_2=\SI{111}{\ohm}$;
(c) a compensada dissipa muito menos, mas só é exata para aquela carga}
\resolucao{(a) A \textbf{corrente de sangria} é a que circula por $R_2$ sem
passar pela carga. Se ela for muito maior que a da carga, conectar a carga
perturba pouco o divisor. Com $I_L=5/1000=\SI{5}{\milli\ampere}$, exigir
sangria $\ge20I_L=\SI{100}{\milli\ampere}$ dá
$R_2\le5/0{,}1=\SI{50}{\ohm}$. Testando valores (com
$R_1=1{,}4R_2$ para manter a razão $5/12$ em vazio):
\begin{center}
\begin{tabular}{cccc}
\hline
$R_2$ & $R_1$ & $U_{saída}$ & erro\\
\hline
$\SI{100}{\ohm}$ & $\SI{140}{\ohm}$ & $\SI{4.72}{\volt}$ & $-5{,}5\%$\\
$\SI{50}{\ohm}$  & $\SI{70}{\ohm}$  & $\SI{4.86}{\volt}$ & $-2{,}8\%$\\
$\SI{25}{\ohm}$  & $\SI{35}{\ohm}$  & $\SI{4.93}{\volt}$ & $-1{,}4\%$\\
\hline
\end{tabular}
\end{center}
Escolha $R_2=\SI{25}{\ohm}$, $R_1=\SI{35}{\ohm}$.\\
(b) \textbf{Compensação:} projete com a carga \emph{já incluída}. Exigindo
$12\cdot\dfrac{x}{R_1+x}=5$ com $x=R_2\parallel1000$, vem $x=\dfrac{5R_1}{7}$.
Tomando $R_1=\SI{140}{\ohm}$: $x=\SI{100}{\ohm}$ e
\[
\frac{1}{R_2}=\frac{1}{100}-\frac{1}{1000}
\;\Longrightarrow\;
R_2=\SI{111}{\ohm}.
\]
Verificação: $111\parallel1000=\SI{100}{\ohm}$ e
$12\cdot100/240=\SI{5.00}{\volt}$ \checkmark.\\
(c) \emph{Potência:} na solução (a), a corrente do divisor é
$\approx\SI{200}{\milli\ampere}$ e a dissipação total
$\approx\SI{2.4}{\watt}$; na (b), $\approx\SI{50}{\milli\ampere}$ e
$\approx\SI{0.6}{\watt}$ --- quatro vezes menos.\\
\emph{Robustez:} a solução (b) é exata \textbf{apenas} para
$R_L=\SI{1}{\kilo\ohm}$; se a carga for retirada, a saída sobe para
$12\cdot111/251=\SI{5.31}{\volt}$ ($+6\%$). A solução (a) é menos precisa,
porém quase indiferente à carga. \textbf{A escolha depende de a carga ser fixa
e conhecida ou variável} --- e, se for variável, nenhuma das duas serve:
o caso pede um regulador.}
\end{questao}

\begin{questao}[4]{}
\textbf{(Enumeração de topologias.)} Dispõe-se de quatro resistores iguais de
$\SI{100}{\ohm}$.
\begin{itensq}
\item Determine todos os valores distintos de $R_{eq}$ obteníveis usando os
      quatro.
\item Identifique os dois arranjos que resultam em $\SI{100}{\ohm}$ e
      compare a potência do componente mais solicitado sob
      $\SI{100}{\volt}$.
\item Justifique por que os extremos são $\SI{25}{\ohm}$ e
      $\SI{400}{\ohm}$.
\end{itensq}
\resposta{(a) $25$, $40$, $60$, $75$, $100$, $133{,}3$, $166{,}7$, $250$ e
$\SI{400}{\ohm}$; (b) $\SI{25}{\watt}$ em ambos}
\resolucao{(a) Enumerando as topologias possíveis:
\begin{center}
\begin{tabular}{ll}
\hline
Arranjo & $R_{eq}$\\
\hline
4 em paralelo & $\SI{25}{\ohm}$\\
$(2\text{ série})\parallel(2\text{ paralelo})$ & $\SI{40}{\ohm}$\\
$[(2\parallel)+1]\parallel 1$ & $\SI{60}{\ohm}$\\
$(3\text{ série})\parallel 1$ & $\SI{75}{\ohm}$\\
$(2\text{ série})\parallel(2\text{ série})$ & $\SI{100}{\ohm}$\\
$(2\parallel)+(2\parallel)$ & $\SI{100}{\ohm}$\\
$(3\parallel)+1$ & $\SI{133.3}{\ohm}$\\
$1+[(2\text{ série})\parallel 1]$ & $\SI{166.7}{\ohm}$\\
$(2\parallel)+2\text{ série}$ & $\SI{250}{\ohm}$\\
4 em série & $\SI{400}{\ohm}$\\
\hline
\end{tabular}
\end{center}
Nove valores distintos.\\
(b) Sob $\SI{100}{\volt}$, ambos os arranjos de $\SI{100}{\ohm}$ dão
$I_t=\SI{1}{\ampere}$ e $P_{tot}=\SI{100}{\watt}$.\\
\emph{Dois ramos de dois em série:} cada ramo conduz $\SI{0.5}{\ampere}$
$\Rightarrow$ $P=100(0{,}25)=\SI{25}{\watt}$ por resistor.\\
\emph{Dois blocos paralelos em série:} cada bloco sob $\SI{50}{\volt}$, cada
resistor com $\SI{0.5}{\ampere}$ $\Rightarrow$ também
$\SI{25}{\watt}$.\\
Os dois arranjos são \textbf{topologicamente distintos, mas termicamente
equivalentes} --- ambos distribuem a dissipação em quatro partes iguais. É o
melhor aproveitamento possível dos componentes.\\
(c) O mínimo é o paralelo total, pois cada resistor acrescentado em paralelo só
pode \emph{reduzir} $R_{eq}$; o máximo é a série total, pela razão simétrica.
Qualquer topologia mista fica necessariamente entre $R/n$ e $nR$.}
\end{questao}

\begin{questao}[4]{}
\textbf{(Projeto sob duas restrições.)} Uma rede deve apresentar
$R_{eq}=\SI{20}{\ohm}$ e dissipar $\SI{45}{\watt}$.
\begin{itensq}
\item Determine a tensão e a corrente nominais.
\item Compare três topologias quanto à potência do componente mais solicitado.
\item Escolha a melhor e justifique.
\end{itensq}
\resposta{(a) $\SI{30}{\volt}$ e $\SI{1.5}{\ampere}$; (b) pior componente:
$\SI{30}{\watt}$, $\SI{22.5}{\watt}$ e $\SI{11.25}{\watt}$;
(c) a matriz $2\times2$ de $\SI{20}{\ohm}$}
\resolucao{(a) De $P=V^{2}/R_{eq}$ e $P=RI^{2}$:
\[
V=\sqrt{PR_{eq}}=\sqrt{45\cdot20}=\SI{30}{\volt},
\qquad
I=\sqrt{P/R_{eq}}=\sqrt{2{,}25}=\SI{1.5}{\ampere}.
\]
(b) Três topologias que dão $\SI{20}{\ohm}$:
\begin{center}
\begin{tabular}{lll}
\hline
Topologia & Potências & Pior\\
\hline
$\SI{30}{\ohm}\parallel\SI{60}{\ohm}$ & $30$; $\SI{15}{\watt}$
  & $\SI{30}{\watt}$\\
$\SI{10}{\ohm}+(\SI{20}{\ohm}\parallel\SI{20}{\ohm})$
  & $22{,}5$; $11{,}25$; $\SI{11.25}{\watt}$ & $\SI{22.5}{\watt}$\\
$(\SI{20}{\ohm}+\SI{20}{\ohm})\parallel(\SI{20}{\ohm}+\SI{20}{\ohm})$
  & $11{,}25\times4$ & $\SI{11.25}{\watt}$\\
\hline
\end{tabular}
\end{center}
Em todas, o total é $\SI{45}{\watt}$ --- como tem de ser, pois $R_{eq}$ e $V$
são os mesmos.\\
(c) \textbf{A matriz $2\times2$} de resistores de $\SI{20}{\ohm}$. Ela reparte
a dissipação em quatro parcelas \emph{iguais} de $\SI{11.25}{\watt}$,
permitindo componentes de $\SI{25}{\watt}$ (fator $2$ de folga) e um único
valor de estoque. As alternativas exigiriam um componente de $\SI{50}{\watt}$
ou de $\SI{50}{\watt}$/$\SI{25}{\watt}$ mistos, com ponto quente concentrado.\\
\textbf{Princípio geral:} para uma dada $R_{eq}$ e potência total, a topologia
ótima é a que \textbf{equaliza} a dissipação --- matrizes $n\times n$ de
resistores idênticos são o caso limite dessa ideia.}
\end{questao}

\begin{questao}[4]{}
\textbf{(Diagnóstico quantitativo.)} Uma rede
$R_1=\SI{10}{\ohm}+(R_2=\SI{20}{\ohm}\parallel R_3=\SI{20}{\ohm})$ opera sob
$\SI{24}{\volt}$. Mediu-se $I_t=\SI{2.4}{\ampere}$.
\begin{itensq}
\item Mostre que a rede está defeituosa e identifique o tipo de falha.
\item Localize o componente e prove que a localização é única.
\item Proponha uma medição de confirmação.
\end{itensq}
\resposta{$R_{eq}$ medida $=\SI{10}{\ohm}$ contra $\SI{20}{\ohm}$ prevista;
o bloco paralelo está em curto}
\resolucao{(a) $R_{eq}^{med}=24/2{,}4=\SI{10}{\ohm}$, contra
$\SI{20}{\ohm}$ previstos. Resistência \emph{menor} que a nominal
$\Rightarrow$ falha do tipo \textbf{curto}. (Se fosse maior, seria abertura.)\\
(b) A diferença é de exatamente $\SI{10}{\ohm}$, que é o valor do bloco
paralelo --- ele está anulado. \emph{Unicidade:} $R_1$ não pode estar em curto,
pois isso daria $R_{eq}=\SI{10}{\ohm}$ também... \textbf{e aqui está o ponto:}
as duas hipóteses são numericamente idênticas! Um curto em $R_1$ deixaria
$R_{eq}=0+10=\SI{10}{\ohm}$, exatamente o medido.\\
Logo a medição de corrente \textbf{não} distingue as duas falhas --- e afirmar
localização única com esse único dado seria erro de método. O item (c) é
obrigatório, não opcional.\\
(c) \textbf{Medição decisiva:} um voltímetro sobre $R_1$.
\begin{itemize}
\item Se o bloco paralelo está em curto: $U_1=10(2{,}4)=\SI{24}{\volt}$
      (toda a fonte).
\item Se $R_1$ está em curto: $U_1=\SI{0}{\volt}$, e os
      $\SI{24}{\volt}$ aparecem sobre o bloco paralelo.
\end{itemize}
\textbf{Lição de diagnóstico:} uma única medição escalar raramente identifica
uma falha em rede mista, porque topologias diferentes podem produzir a mesma
$R_{eq}$. É preciso pelo menos uma medida de \emph{distribuição} --- uma
tensão interna --- para quebrar a ambiguidade.}
\end{questao}

\begin{questao}[4]{}
\textbf{(Redução por simetria.)} Seis resistores de $\SI{6}{\ohm}$ formam as
arestas de um \textbf{tetraedro} regular (quatro vértices, todos ligados dois a
dois).
\begin{itensq}
\item Determine $R_{AB}$ entre dois vértices quaisquer, usando simetria.
\item Determine a corrente em cada aresta para $I_t=\SI{4}{\ampere}$.
\item Recalcule $R_{AB}$ se a aresta $AB$ for removida.
\end{itensq}
\resposta{(a) $R_{AB}=\SI{3}{\ohm}$; (c) $\SI{6}{\ohm}$}
\resolucao{(a) Sejam $A$ e $B$ os terminais e $C$, $D$ os outros dois vértices.
Como $C$ e $D$ ocupam posições \textbf{simétricas} em relação ao par $A$--$B$
(cada um ligado a $A$, a $B$ e ao outro, com resistências iguais), tem-se
$V_C=V_D$. A aresta $CD$ não conduz e pode ser \textbf{removida}.\\
Restam três caminhos em paralelo: a aresta direta $AB$ ($\SI{6}{\ohm}$), o
percurso $A$--$C$--$B$ ($\SI{12}{\ohm}$) e o percurso $A$--$D$--$B$
($\SI{12}{\ohm}$):
\[
R_{AB}=6\parallel12\parallel12=6\parallel6=\SI{3}{\ohm}=\frac{R}{2}.
\]
(b) Com $I_t=\SI{4}{\ampere}$, $U_{AB}=3(4)=\SI{12}{\volt}$:
\begin{itemize}
\item aresta $AB$: $12/6=\SI{2}{\ampere}$;
\item cada percurso lateral: $12/12=\SI{1}{\ampere}$;
\item aresta $CD$: $\SI{0}{\ampere}$.
\end{itemize}
Soma: $2+1+1=\SI{4}{\ampere}$ \checkmark. A aresta direta leva metade da
corrente total.\\
(c) Sem a aresta $AB$, a simetria $C\leftrightarrow D$ \emph{persiste} --- ela
não dependia daquela aresta. Logo $V_C=V_D$ ainda, e $CD$ continua sem
conduzir:
\[
R_{AB}=12\parallel12=\SI{6}{\ohm}=R.
\]
\textbf{Resultado elegante:} remover uma aresta dobra a resistência
($3\to\SI{6}{\ohm}$), porque justamente essa aresta carregava metade da
corrente.\\
\textbf{Por que a simetria funciona:} a solução de uma rede resistiva é
\emph{única}; se permutar $C$ com $D$ devolve o mesmo circuito, a solução
permutada é a mesma solução --- e isso exige $V_C=V_D$. O argumento é o mesmo
que resolve a ponte equilibrada e o cubo de resistores, e dispensa qualquer
sistema de equações.}
\end{questao}

% END BANCO COMPLEMENTAR

\gabarito[3.3 Circuitos mistos]

% =====================================================================
\subsection{Transformação estrela-triângulo}
% =====================================================================

\begin{conceito}[Quando usar]
Use a transformação $\Delta\leftrightarrow Y$ quando o circuito \textbf{não}
puder ser reduzido apenas por série e paralelo --- o caso típico é a
\emph{ponte de Wheatstone} desequilibrada. Converta um dos triângulos (ou uma das
estrelas) para "abrir" o circuito e revelar novas associações série/paralelo.
\end{conceito}

\legendaniveis

\blocofixacao

\begin{questao}[1]{}
Considere uma configuração \textbf{estrela} com $R_1=\SI{3}{\ohm}$,
$R_2=\SI{6}{\ohm}$ e $R_3=\SI{9}{\ohm}$. Determine os resistores da configuração
\textbf{triângulo} equivalente.
\resposta{$R_{AB}=\SI{11}{\ohm}$, $R_{BC}=\SI{33}{\ohm}$, $R_{CA}=\SI{16.5}{\ohm}$}
\resolucao{Com $P = R_1R_2+R_2R_3+R_3R_1 = 18+54+27 = \SI{99}{\ohm\squared}$, cada
resistor do triângulo é $P$ dividido pelo resistor \emph{oposto} da estrela:
\[
R_{AB}=\frac{P}{R_3}=\frac{99}{9}=\SI{11}{\ohm},\quad
R_{BC}=\frac{P}{R_1}=\frac{99}{3}=\SI{33}{\ohm},\quad
R_{CA}=\frac{P}{R_2}=\frac{99}{6}=\SI{16.5}{\ohm}.
\]}
\end{questao}

\begin{questao}[1]{}
Um circuito \textbf{triângulo} possui $R_{AB}=\SI{10}{\ohm}$,
$R_{BC}=\SI{5}{\ohm}$ e $R_{CA}=\SI{15}{\ohm}$. Calcule os resistores da
\textbf{estrela} equivalente.
\resposta{$R_A=\SI{5}{\ohm}$, $R_B=\SI{1.67}{\ohm}$, $R_C=\SI{2.5}{\ohm}$}
\resolucao{Com $S = R_{AB}+R_{BC}+R_{CA} = 10+5+15 = \SI{30}{\ohm}$, cada resistor
da estrela é o produto dos dois lados \emph{adjacentes} ao seu vértice, dividido
por $S$:
\[
R_A=\frac{R_{AB}R_{CA}}{S}=\frac{10\cdot15}{30}=\SI{5}{\ohm},\quad
R_B=\frac{R_{AB}R_{BC}}{S}=\frac{10\cdot5}{30}\approx\SI{1.67}{\ohm},\quad
R_C=\frac{R_{BC}R_{CA}}{S}=\frac{5\cdot15}{30}=\SI{2.5}{\ohm}.
\]}
\end{questao}

\begin{questao}[1]{}
Explique, com suas palavras, por que a transformação estrela-triângulo é útil e
cite um exemplo prático de circuito em que ela é necessária.
\resposta{Resposta dissertativa --- ver comentário}
\resolucao{Ela permite substituir um agrupamento de três resistores em triângulo
por três em estrela (ou vice-versa) \emph{sem alterar} o comportamento elétrico
visto pelos três terminais externos. Isso "desmancha" configurações que não são
nem série nem paralelo, revelando novas associações simples. O exemplo clássico é
a \textbf{ponte de Wheatstone desequilibrada}: convertendo um dos triângulos em
estrela, os ramos passam a ficar em série/paralelo e o circuito se resolve sem
precisar montar um sistema de equações. Outra aplicação é a análise de redes
trifásicas, onde cargas ligadas em $\Delta$ e em $Y$ são convertidas entre si.}
\end{questao}

\blocoaplicacao

\begin{questao}[2]{}
No triângulo da figura, com $R_{AB}=\SI{9}{\ohm}$ (lado direto entre A e B),
$R_{AC}=\SI{6}{\ohm}$ e $R_{CB}=\SI{3}{\ohm}$, determine a resistência
equivalente entre A e B.

\circuitoq{%
\begin{circuitikz}[scale=1,transform shape,line width=0.8pt]
 \draw (0,0) node[below left,font=\bfseries]{A} coordinate(A)
       to[R,l=$\SI{6}{\ohm}$] (2,2.4) node[above,font=\bfseries]{C} coordinate(C)
       to[R,l=$\SI{3}{\ohm}$] (4,0) node[below right,font=\bfseries]{B} coordinate(B);
 \draw (A) to[R,l_=$\SI{9}{\ohm}$] (B);
\end{circuitikz}}

\resposta{$R_{AB} = \SI{4.5}{\ohm}$}
\resolucao{O caminho superior A--C--B está em \emph{série}:
$6+3 = \SI{9}{\ohm}$. Esse ramo está em \emph{paralelo} com o lado direto
$R_{AB} = \SI{9}{\ohm}$:
\[
R_{AB}=9\parallel 9=\frac{9}{2}=\SI{4.5}{\ohm}.
\]
Aqui a redução série/paralelo bastou. Guarde a transformação $\Delta$-$Y$ para
quando há um resistor \emph{atravessando} o triângulo (a ponte), como nas
questões seguintes.}
\end{questao}

\begin{questao}[2]{}
Três resistores iguais de $\SI{2}{\ohm}$ estão ligados em \textbf{estrela}.
Convertendo-os em \textbf{triângulo} e aplicando uma fonte de $\SI{12}{\volt}$
entre dois dos terminais (A e B), determine a corrente total fornecida pela fonte.
\resposta{$I = \SI{3}{\ampere}$}
\resolucao{A estrela de três resistores iguais $R$ vira um triângulo de
resistores $3R = \SI{6}{\ohm}$ cada. Entre A e B temos o lado direto de
\SI{6}{\ohm} em paralelo com o caminho A--C--B ($6+6 = \SI{12}{\ohm}$):
\[
R_{AB}=6\parallel 12=\frac{6\cdot12}{18}=\SI{4}{\ohm}
\quad\Rightarrow\quad
I=\frac{12}{4}=\SI{3}{\ampere}.
\]
\textbf{Confira:} entre A e B, a estrela original também dá o mesmo resultado ---
$R_A + (R_B \parallel$ com o caminho por C$)$ ---, como esperado, já que as duas
representações são equivalentes vistas dos terminais.}
\end{questao}

\begin{questao}[2]{}
Um triângulo entre A, B e C tem $R_{AB}=\SI{12}{\ohm}$, $R_{BC}=\SI{6}{\ohm}$ e
$R_{CA}=\SI{6}{\ohm}$.
\begin{itensq}
\item Converta para estrela.
\item Calcule a resistência equivalente entre A e B após a conversão.
\end{itensq}
\resposta{(a) $R_A=\SI{3}{\ohm}$, $R_B=\SI{3}{\ohm}$, $R_C=\SI{1.5}{\ohm}$;
(b) $R_{AB}=\SI{6}{\ohm}$}
\resolucao{(a) $S = 12+6+6 = \SI{24}{\ohm}$:
\[
R_A=\frac{R_{AB}R_{CA}}{S}=\frac{12\cdot6}{24}=\SI{3}{\ohm},\quad
R_B=\frac{R_{AB}R_{BC}}{S}=\frac{12\cdot6}{24}=\SI{3}{\ohm},\quad
R_C=\frac{R_{BC}R_{CA}}{S}=\frac{6\cdot6}{24}=\SI{1.5}{\ohm}.
\]
(b) Na estrela, entre A e B a corrente percorre $R_A$ e $R_B$ em série (o braço
$R_C$ fica em aberto, pois C não é terminal):
$R_{AB} = R_A + R_B = 3+3 = \SI{6}{\ohm}$.}
\end{questao}

\blocoaprofundamento

\begin{questao}[3]{}
A figura mostra uma \textbf{ponte de Wheatstone}. Verifique se ela está
equilibrada e, caso esteja, determine a resistência equivalente entre A e B.

\circuitoq{%
\begin{circuitikz}[scale=1,transform shape,line width=0.8pt]
 \draw (0,0) node[left,font=\bfseries]{A} coordinate(A);
 \draw (5,0) node[right,font=\bfseries]{B} coordinate(B);
 \coordinate (C) at (2.5,1.7);
 \coordinate (D) at (2.5,-1.7);
 \draw (A) to[R,l=$\SI{10}{\ohm}$] (C);
 \draw (C) to[R,l=$\SI{20}{\ohm}$] (B);
 \draw (A) to[R,l_=$\SI{15}{\ohm}$] (D);
 \draw (D) to[R,l_=$\SI{30}{\ohm}$] (B);
 \draw (C) to[R,l=$\SI{8}{\ohm}$] (D);
 \node[font=\footnotesize] at (3.1,-0.27) {galv.};
\end{circuitikz}}

\resposta{Equilibrada; $R_{AB} = \SI{18}{\ohm}$}
\resolucao{A ponte está \textbf{equilibrada} quando os produtos dos braços opostos
são iguais:
\[
\frac{R_{AC}}{R_{CB}}=\frac{10}{20}=\frac12,
\qquad
\frac{R_{AD}}{R_{DB}}=\frac{15}{30}=\frac12.
\]
As razões são iguais, logo C e D estão no \emph{mesmo potencial}: o resistor
central de \SI{8}{\ohm} não conduz e pode ser removido. Restam dois ramos em série,
associados em paralelo:
\[
R_{AB}=(10+20)\parallel(15+30)=30\parallel45=\frac{30\cdot45}{75}=\SI{18}{\ohm}.
\]}
\end{questao}

\begin{questao}[3]{}
Na ponte da figura, $R_{AC}=\SI{6}{\ohm}$, $R_{CB}=\SI{18}{\ohm}$,
$R_{AD}=\SI{18}{\ohm}$, $R_{DB}=\SI{6}{\ohm}$ e o resistor central (ponte) vale
$R_{CD}=\SI{6}{\ohm}$. Determine a resistência equivalente entre A e B.

\circuitoq{%
\begin{circuitikz}[scale=1,transform shape,line width=0.8pt]
 \draw (0,0) node[left,font=\bfseries]{A} coordinate(A);
 \draw (5,0) node[right,font=\bfseries]{B} coordinate(B);
 \coordinate (C) at (2.5,1.7);
 \coordinate (D) at (2.5,-1.7);
 \draw (A) to[R,l=$\SI{6}{\ohm}$] (C);
 \draw (C) to[R,l=$\SI{18}{\ohm}$] (B);
 \draw (A) to[R,l_=$\SI{18}{\ohm}$] (D);
 \draw (D) to[R,l_=$\SI{6}{\ohm}$] (B);
 \draw (C) to[R,l=$\SI{6}{\ohm}$] (D);
\end{circuitikz}}

\resposta{$R_{AB} = \SI{10}{\ohm}$}
\resolucao{Como $\dfrac{6}{18}\neq\dfrac{18}{6}$, a ponte está
\textbf{desequilibrada} e o resistor central conduz: não dá para reduzir só por
série/paralelo. Convertemos o triângulo superior A--C--D
($R_{AC}=6$, $R_{CD}=6$, $R_{AD}=18$) em estrela, com $S = 6+6+18 = \SI{30}{\ohm}$:
\[
R_{A}=\frac{6\cdot18}{30}=\SI{3.6}{\ohm},\quad
R_{C}=\frac{6\cdot6}{30}=\SI{1.2}{\ohm},\quad
R_{D}=\frac{6\cdot18}{30}=\SI{3.6}{\ohm}.
\]
Chamando de N o centro da estrela, os caminhos até B ficam:
$N$--$C$--$B = 1{,}2+18 = \SI{19.2}{\ohm}$ e
$N$--$D$--$B = 3{,}6+6 = \SI{9.6}{\ohm}$, em paralelo:
\[
19{,}2\parallel 9{,}6 = \frac{19{,}2\cdot9{,}6}{28{,}8}=\SI{6.4}{\ohm}.
\]
Somando o braço $R_A$ que liga A ao centro:
\[
R_{AB}=3{,}6+6{,}4=\SI{10}{\ohm}.
\]}
\end{questao}

\begin{questao}[3]{}
Três resistores formam um triângulo entre os nós X, Y e Z, com
$R_{XY}=\SI{9}{\ohm}$, $R_{YZ}=\SI{12}{\ohm}$ e $R_{ZX}=\SI{15}{\ohm}$. Além
disso, cada nó é ligado ao terra por um resistor de $\SI{6}{\ohm}$. Descreva o
procedimento para obter a resistência entre o nó X e o terra e monte a expressão
final (sem exigir o valor numérico fechado).
\resposta{Ver roteiro comentado}
\resolucao{\textbf{Passo 1.} Converta o triângulo XYZ em estrela, com centro N e
$S = 9+12+15 = \SI{36}{\ohm}$:
\[
R_X=\frac{R_{XY}R_{ZX}}{S}=\frac{9\cdot15}{36}=\SI{3.75}{\ohm},\quad
R_Y=\frac{R_{XY}R_{YZ}}{S}=\frac{9\cdot12}{36}=\SI{3}{\ohm},\quad
R_Z=\frac{R_{YZ}R_{ZX}}{S}=\frac{12\cdot15}{36}=\SI{5}{\ohm}.
\]
\textbf{Passo 2.} Agora cada terminal externo Y e Z leva ao terra por
$R_Y+6$ e $R_Z+6$, ambos em paralelo (pois se encontram no terra):
$(3+6)\parallel(5+6) = 9\parallel 11$. \textbf{Passo 3.} Esse resultado fica em
série com $R_X$, e todo o conjunto em paralelo com o resistor de \SI{6}{\ohm} que
liga X diretamente ao terra:
\[
R_{X\to\text{terra}}
=\Big[\,R_X+\big(9\parallel11\big)\Big]\ \big\|\ 6
=\Big[\,3{,}75+\tfrac{99}{20}\Big]\ \big\|\ 6 .
\]
O que se pretende com esta questão é o \emph{método}: transformar o triângulo em
estrela desmonta o emaranhado e transforma tudo em série/paralelo.}
\end{questao}

\begin{questao}[3]{}
Uma ponte de Wheatstone é usada para medir uma resistência desconhecida $R_x$.
Os braços conhecidos são $R_1 = \SI{100}{\ohm}$, $R_2 = \SI{200}{\ohm}$ e
$R_3$ é uma década ajustável. O equilíbrio é obtido com
$R_3 = \SI{347}{\ohm}$.
\begin{itensq}
\item Determine $R_x$.
\item Se $R_1$ e $R_2$ têm tolerância de $\SI{0.1}{\percent}$ cada e $R_3$ de
      $\SI{0.05}{\percent}$, estime a incerteza relativa de $R_x$.
\item Por que a ponte é mais precisa que medir $R_x$ com um ohmímetro comum?
\end{itensq}
\resposta{(a) $R_x = \SI{694}{\ohm}$; (b) $\approx\SI{0.25}{\percent}$;
(c) ver comentário}
\resolucao{(a) No equilíbrio, os produtos cruzados se igualam:
\[
\frac{R_1}{R_2}=\frac{R_3}{R_x}
\;\Rightarrow\;
R_x = R_3\,\frac{R_2}{R_1} = 347\cdot\frac{200}{100} = \SI{694}{\ohm}.
\]
(b) Para um produto/quociente, as incertezas \emph{relativas} se somam (no pior
caso):
\[
\frac{\delta R_x}{R_x} = \frac{\delta R_3}{R_3}+\frac{\delta R_2}{R_2}
+\frac{\delta R_1}{R_1} = 0{,}05\% + 0{,}1\% + 0{,}1\% = \SI{0.25}{\percent},
\]
ou seja, $R_x = \SI{694}{\ohm} \pm \SI{1.7}{\ohm}$.
(Somando em quadratura, para erros independentes, obtém-se
$\sqrt{0{,}05^2+0{,}1^2+0{,}1^2}\approx\SI{0.15}{\percent}$ --- estimativa mais
realista.)\\
(c) Porque a ponte é um método de \textbf{zero} (ou de detecção nula): no
equilíbrio, o galvanômetro indica \emph{zero}, e a medida não depende da precisão
do próprio galvanômetro nem da tensão da fonte --- depende apenas das
\emph{razões} entre resistores conhecidos. Um ohmímetro comum, ao contrário,
depende da estabilidade da sua fonte interna e da calibração do seu indicador.
\textbf{Este é um princípio geral da metrologia:} comparar contra um padrão é
sempre mais preciso do que medir em escala absoluta.}
\end{questao}

% BEGIN BANCO COMPLEMENTAR Transformação estrela-triângulo
% =====================================================================
%  Banco complementar --- Transformacao Estrela-Triangulo  (REVISADO)
%  Substitui a versao gerada por template (9 clones em N2, 8 em N3 e
%  8 questoes conceituais marcadas como N4).
%  O cap03 ja cobre: estrela 3/6/9 -> triangulo, triangulo 10/5/15 ->
%  estrela, "explique por que a transformacao e util", triangulo
%  9/6/3 com figura, tres de 2 ohm em estrela com fonte, triangulo
%  12/6/6 com itens, e o triangulo XYZ 9/12/15. As conversoes basicas,
%  portanto, ja estao exaustivamente cobertas la.
%  Este banco explora: reversibilidade e verificacao, casos-limite
%  (lado nulo, lado infinito), ponte desequilibrada por dois caminhos,
%  condicionamento numerico, valores comerciais, e as DEDUCOES das
%  duas formulas a partir das resistencias entre terminais.
%  Distribuicao mantida: 5 N1 + 9 N2 + 8 N3 + 8 N4 = 30
% =====================================================================

\subsubsection*{Banco complementar --- Transformação estrela-triângulo}

\begin{atencao}[Organização do banco]
As duas fórmulas têm estruturas espelhadas. \textbf{Triângulo para estrela:}
cada braço é o \emph{produto} dos dois lados que tocam aquele terminal, dividido
pela \emph{soma} dos três lados. \textbf{Estrela para triângulo:} cada lado é a
soma dos três produtos de braços, dividida pelo braço \emph{oposto}. O N3 trata
de pontes, casos-limite e condicionamento; o N4 exige as deduções, provas de
positividade e a análise de valores comerciais.
\end{atencao}

\blocofixacao

\begin{questao}[1]{}
Um triângulo equilátero tem três lados de $\SI{30}{\ohm}$. Determine a estrela
equivalente.
\resposta{Três braços de $\SI{10}{\ohm}$}
\resolucao{Para o caso equilibrado, a fórmula geral colapsa:
\[
R_Y=\frac{R_\Delta\cdot R_\Delta}{3R_\Delta}=\frac{R_\Delta}{3}
=\frac{30}{3}=\SI{10}{\ohm}.
\]
\textbf{Regra rápida:} $R_Y=R_\Delta/3$. Guarde o sentido --- a estrela é
sempre \emph{menor}, porque seus dois braços em série ($2R_Y=\SI{20}{\ohm}$)
devem igualar o paralelo do triângulo
($30\parallel60=\SI{20}{\ohm}$) \checkmark.}
\end{questao}

\begin{questao}[1]{}
Uma estrela equilibrada tem três braços de $\SI{5}{\ohm}$. Determine o
triângulo equivalente.
\resposta{Três lados de $\SI{15}{\ohm}$}
\resolucao{
\[
R_\Delta=\frac{3R_Y^{2}}{R_Y}=3R_Y=3(5)=\SI{15}{\ohm}.
\]
\textbf{Regra rápida:} $R_\Delta=3R_Y$ --- inversa da anterior. Se
$R_Y=R_\Delta/3$ e $R_\Delta=3R_Y$ não fossem consistentes, uma das fórmulas
estaria trocada. É a conferência mais barata que existe para o caso
equilibrado.}
\end{questao}

\begin{questao}[1]{}
Escreva, \textbf{sem calcular}, a expressão do braço $R_B$ de uma estrela
equivalente a um triângulo de lados $R_{AB}$, $R_{BC}$ e $R_{CA}$.
\resposta{$R_B=\dfrac{R_{AB}\,R_{BC}}{R_{AB}+R_{BC}+R_{CA}}$}
\resolucao{A regra é posicional: o braço ligado ao terminal $B$ é o
\textbf{produto dos dois lados que tocam $B$} --- que são $R_{AB}$ e
$R_{BC}$ --- dividido pela \textbf{soma dos três lados}.\\
\textbf{Como não errar:} identifique o terminal, olhe quais dois lados chegam
nele e multiplique. O lado \emph{oposto} ($R_{CA}$) aparece apenas no
denominador, junto com os outros dois.}
\end{questao}

\begin{questao}[1]{}
Compare estrela e triângulo quanto ao número de elementos e de nós.
\begin{alternativas}[2]
\item ambos têm 3 resistores; a estrela tem um nó interno adicional
\item a estrela tem 3 e o triângulo tem 4 resistores
\item ambos têm o mesmo número de nós
\item o triângulo tem um nó interno adicional
\end{alternativas}
\resposta{(a)}
\resolucao{Ambas as configurações usam \textbf{três} resistores e possuem os
mesmos três terminais externos ($A$, $B$, $C$). A diferença é o \textbf{nó
central} da estrela, que não tem correspondente no triângulo.\\
\textbf{Consequência importante:} a equivalência vale apenas para o que se
observa nos \emph{três terminais externos}. O potencial do nó central da estrela
não corresponde a ponto algum do triângulo --- e nenhuma corrente interna é
preservada pela transformação.}
\end{questao}

\begin{questao}[1]{}
Identifique qual das expressões é a conversão triângulo-estrela:
\begin{alternativas}[2]
\item $R=\dfrac{\text{produto de dois}}{\text{soma dos três}}$
\item $R=\dfrac{\text{soma dos produtos}}{\text{oposto}}$
\item $R=\dfrac{\text{soma dos três}}{3}$
\item $R=\dfrac{\text{produto dos três}}{\text{soma dos três}}$
\end{alternativas}
\resposta{(a)}
\resolucao{\textbf{Triângulo $\to$ estrela:} produto de dois sobre a soma dos
três --- alternativa (a). O resultado tem dimensão de $\si{\ohm}$
($\si{\ohm^2}/\si{\ohm}$) \checkmark.\\
\textbf{Estrela $\to$ triângulo:} soma dos produtos sobre o braço oposto ---
alternativa (b). Também $\si{\ohm}$ \checkmark.\\
\textbf{Teste de sanidade:} a alternativa (d) daria $\si{\ohm^{2}}$ --- errada
por análise dimensional. E, no caso equilibrado, (a) devolve $R/3$ e (b)
devolve $3R$, o que confirma qual é qual.}
\end{questao}

\blocoaplicacao

\begin{questao}[2]{}
Um triângulo tem $R_{AB}=\SI{20}{\ohm}$, $R_{BC}=\SI{30}{\ohm}$ e
$R_{CA}=\SI{50}{\ohm}$. Determine a estrela equivalente.
\resposta{$R_A=\SI{10}{\ohm}$, $R_B=\SI{6}{\ohm}$, $R_C=\SI{15}{\ohm}$}
\resolucao{A soma dos lados é $20+30+50=\SI{100}{\ohm}$.
\[
R_A=\frac{R_{AB}R_{CA}}{100}=\frac{20(50)}{100}=\SI{10}{\ohm};
\]
\[
R_B=\frac{R_{AB}R_{BC}}{100}=\frac{20(30)}{100}=\SI{6}{\ohm};
\qquad
R_C=\frac{R_{BC}R_{CA}}{100}=\frac{30(50)}{100}=\SI{15}{\ohm}.
\]
\textbf{Conferência de plausibilidade:} cada braço é menor que os lados que o
originaram, e o maior braço ($R_C$) está no terminal tocado pelos dois maiores
lados. Coerente.}
\end{questao}

\begin{questao}[2]{}
Converta de volta a estrela obtida ($10$, $6$ e $\SI{15}{\ohm}$) para
triângulo e verifique a reversibilidade.
\resposta{$20$, $30$ e $\SI{50}{\ohm}$ --- os valores originais}
\resolucao{Soma dos produtos:
\[
\Sigma=R_AR_B+R_BR_C+R_CR_A=60+90+150=\SI{300}{\ohm\squared}.
\]
Cada lado é $\Sigma$ dividido pelo braço \emph{oposto}:
\[
R_{AB}=\frac{300}{R_C}=\frac{300}{15}=\SI{20}{\ohm};
\quad
R_{BC}=\frac{300}{R_A}=\SI{30}{\ohm};
\quad
R_{CA}=\frac{300}{R_B}=\SI{50}{\ohm}.
\]
\textbf{Reversibilidade confirmada.} As duas transformações são inversas exatas
--- fato garantido pela construção, já que ambas derivam do mesmo sistema de
três equações. Fazer a ida e a volta é a verificação mais completa possível de
uma conversão.}
\end{questao}

\begin{questao}[2]{}
Para o triângulo $20/30/\SI{50}{\ohm}$ e sua estrela equivalente, calcule a
resistência entre $A$ e $B$ (com $C$ aberto) nas duas configurações.
\resposta{$\SI{16}{\ohm}$ em ambas}
\resolucao{\textbf{No triângulo:} o lado direto $R_{AB}$ em paralelo com o
caminho $A$--$C$--$B$:
\[
R_{AB}^{eq}=20\parallel(50+30)=\frac{20(80)}{100}=\SI{16}{\ohm}.
\]
\textbf{Na estrela:} com $C$ aberto, o braço $R_C$ não conduz; sobram
$R_A$ e $R_B$ em série:
\[
R_{AB}^{eq}=10+6=\SI{16}{\ohm}.\ \checkmark
\]
\textbf{É exatamente esta a definição de equivalência:} a resistência vista
entre cada par de terminais, com o terceiro aberto, deve coincidir. As três
igualdades desse tipo constituem o sistema de onde as fórmulas são deduzidas ---
como se mostra no bloco de desafio.}
\end{questao}

\begin{questao}[2]{}
Uma estrela tem $R_A=\SI{4}{\ohm}$, $R_B=\SI{8}{\ohm}$ e
$R_C=\SI{12}{\ohm}$. Determine o triângulo equivalente e identifique o maior
lado.
\resposta{$R_{AB}=\SI{14.67}{\ohm}$, $R_{BC}=\SI{44}{\ohm}$,
$R_{CA}=\SI{22}{\ohm}$; o maior é $R_{BC}$}
\resolucao{
\[
\Sigma=4(8)+8(12)+12(4)=32+96+48=\SI{176}{\ohm\squared}.
\]
\[
R_{AB}=\frac{176}{12}=\SI{14.67}{\ohm};
\qquad
R_{BC}=\frac{176}{4}=\SI{44}{\ohm};
\qquad
R_{CA}=\frac{176}{8}=\SI{22}{\ohm}.
\]
\textbf{Padrão:} o maior lado do triângulo é o \emph{oposto} ao menor braço da
estrela. Como $R_A=\SI{4}{\ohm}$ é o menor braço, $R_{BC}$ é o maior lado ---
a divisão por um número pequeno produz um resultado grande.\\
\textbf{Interpretação:} um braço pequeno "aproxima" seu terminal do centro,
o que equivale, no triângulo, a afastar os outros dois entre si.}
\end{questao}

\begin{questao}[2]{}
Em um triângulo, o lado $R_{BC}$ é substituído por um \textbf{circuito aberto}.
Determine a estrela equivalente e interprete.
\resposta{$R_A=\dfrac{R_{AB}R_{CA}}{R_{AB}+R_{CA}}$, $R_B=R_C=0$ --- ou seja,
apenas dois resistores em série}
\resolucao{Com $R_{BC}\to\infty$, a soma dos lados também diverge. Nos numeradores:
\[
R_A=\frac{R_{AB}R_{CA}}{R_{AB}+R_{BC}+R_{CA}}\to0
\]
--- mas os braços $R_B$ e $R_C$ contêm $R_{BC}$ no numerador e o limite exige
cuidado:
\[
R_B=\frac{R_{AB}R_{BC}}{R_{AB}+R_{BC}+R_{CA}}\to R_{AB};
\qquad
R_C=\frac{R_{BC}R_{CA}}{R_{AB}+R_{BC}+R_{CA}}\to R_{CA}.
\]
\textbf{Resultado:} $R_A=0$, $R_B=R_{AB}$ e $R_C=R_{CA}$. A estrela degenera em
$R_{AB}$ e $R_{CA}$ ligados diretamente em série pelo terminal $A$ --- que é
exatamente o circuito restante. \checkmark\\
\textbf{Lição:} os casos-limite não quebram as fórmulas; eles as confirmam. Um
lado ausente do triângulo apenas devolve a associação série trivial.}
\end{questao}

\begin{questao}[2]{}
Em um triângulo, o lado $R_{AB}$ é substituído por um \textbf{curto-circuito}.
Determine a estrela equivalente.
\resposta{$R_A=R_B=0$ e $R_C=\dfrac{R_{BC}R_{CA}}{R_{BC}+R_{CA}}$}
\resolucao{Com $R_{AB}=0$, a soma vale $R_{BC}+R_{CA}$ e:
\[
R_A=\frac{0\cdot R_{CA}}{\Sigma}=0;
\qquad
R_B=\frac{0\cdot R_{BC}}{\Sigma}=0;
\qquad
R_C=\frac{R_{BC}R_{CA}}{R_{BC}+R_{CA}}.
\]
\textbf{Interpretação:} os terminais $A$ e $B$ estão em curto, logo ambos
coincidem com o nó central --- daí $R_A=R_B=0$. E $R_C$ é o paralelo de
$R_{BC}$ com $R_{CA}$, o que é correto: com $A$ e $B$ unidos, os dois resistores
ficam mesmo em paralelo entre $C$ e esse nó comum. \checkmark\\
\textbf{Observação:} a conversão \emph{inversa} não existe aqui --- uma estrela
com dois braços nulos daria $\Sigma=0$ e divisão por zero. A transformação
$Y\to\Delta$ exige todos os braços estritamente positivos.}
\end{questao}

\begin{questao}[2]{}
Verifique numericamente, para o triângulo $20/30/\SI{50}{\ohm}$, a identidade
$R_{AB}R_{BC}+R_{BC}R_{CA}+R_{CA}R_{AB}
=\left(\sum R_\Delta\right)\left(\sum R_Y\right)$.
\resposta{$3100=100\times31$ \checkmark}
\resolucao{\textbf{Lado esquerdo:}
\[
20(30)+30(50)+50(20)=600+1500+1000=3100.
\]
\textbf{Lado direito:} $\sum R_\Delta=100$ e
$\sum R_Y=10+6+15=31$, logo $100(31)=3100$ \checkmark.\\
\textbf{Utilidade:} a identidade permite obter a soma dos braços da estrela
\emph{sem calcular cada um}:
\[
\sum R_Y=\frac{\sum_{\text{pares}}R_\Delta R_\Delta}{\sum R_\Delta}
=\frac{3100}{100}=\SI{31}{\ohm}.
\]
É uma conferência rápida do conjunto das três conversões --- se a soma não
bater, ao menos um braço está errado.}
\end{questao}

\begin{questao}[2]{}
Em uma rede, três resistores formam um triângulo entre os nós $P$, $Q$ e $R$,
e há mais resistores ligados a esses três nós. Explique por que convém
converter o triângulo em estrela, e não o contrário.
\resposta{A estrela cria um nó interno isolado, liberando associações
série/paralelo}
\resolucao{\textbf{O problema do triângulo:} cada um de seus lados liga dois nós
que \emph{também} recebem outros elementos. Nenhum lado fica em série ou em
paralelo com nada --- a redução trava.\\
\textbf{O que a estrela resolve:} ela substitui os três lados por três braços
que se encontram em um nó \textbf{novo}, ao qual nada mais está ligado. Cada
braço passa a estar em série com o que já existia em seu terminal, e essas
séries podem então ser combinadas em paralelo.\\
\textbf{Quando fazer o contrário:} se a rede contiver uma \emph{estrela} cujo
nó central recebe outros elementos, converta-a em triângulo --- isso elimina o
nó problemático. \textbf{A regra é sempre a mesma:} transforme na direção que
\emph{elimina} o nó que está impedindo a redução.}
\end{questao}

\begin{questao}[2]{}
Uma estrela tem braços $R_A=R_B=\SI{6}{\ohm}$ e $R_C=\SI{3}{\ohm}$. Determine o
triângulo equivalente e verifique pela resistência entre $A$ e $B$.
\resposta{$R_{AB}=\SI{24}{\ohm}$, $R_{BC}=R_{CA}=\SI{12}{\ohm}$}
\resolucao{
\[
\Sigma=6(6)+6(3)+3(6)=36+18+18=\SI{72}{\ohm\squared};
\]
\[
R_{AB}=\frac{72}{3}=\SI{24}{\ohm};
\qquad
R_{BC}=R_{CA}=\frac{72}{6}=\SI{12}{\ohm}.
\]
\textbf{Verificação (entre $A$ e $B$, com $C$ aberto):}\\
\emph{Estrela:} $R_A+R_B=\SI{12}{\ohm}$.\\
\emph{Triângulo:} $24\parallel(12+12)=\dfrac{24(24)}{48}=\SI{12}{\ohm}$
\checkmark.\\
\textbf{Simetria preservada:} como $R_A=R_B$, o triângulo também sai simétrico
em relação ao par $A$--$B$. Qualquer simetria da estrela reaparece no
triângulo --- propriedade útil para conferir resultados rapidamente.}
\end{questao}

\blocoaprofundamento

\begin{questao}[3]{}
Uma ponte é alimentada por $\SI{84}{\volt}$ entre o nó $S$ e o terra. Do nó
$S$ partem $\SI{6}{\ohm}$ até $A$ e $\SI{12}{\ohm}$ até $B$; de $A$ sai
$\SI{12}{\ohm}$ ao terra e de $B$ sai $\SI{6}{\ohm}$ ao terra; entre $A$ e
$B$ há $\SI{6}{\ohm}$.
\begin{itensq}
\item Verifique que a ponte está desequilibrada.
\item Converta o triângulo $S$--$A$--$B$ em estrela e determine $R_{eq}$.
\item Determine a corrente total.
\end{itensq}
\resposta{(b) $R_{eq}=\SI{8.4}{\ohm}$; (c) $I=\SI{10}{\ampere}$}
\resolucao{(a) Equilíbrio exigiria $\dfrac{6}{12}=\dfrac{12}{6}$, ou seja
$0{,}5=2$ --- falso. \textbf{Desequilibrada}, e o resistor central conduz.
Nenhuma associação série/paralelo é possível.\\
(b) O triângulo formado por $S$, $A$ e $B$ tem lados
$R_{SA}=6$, $R_{SB}=12$ e $R_{AB}=\SI{6}{\ohm}$, com soma
$\SI{24}{\ohm}$. Convertendo em estrela de centro $N$:
\[
R_S=\frac{6(12)}{24}=\SI{3}{\ohm};
\quad
R_A=\frac{6(6)}{24}=\SI{1.5}{\ohm};
\quad
R_B=\frac{12(6)}{24}=\SI{3}{\ohm}.
\]
Agora a rede é redutível: de $N$ partem dois caminhos ao terra,
\[
(1{,}5+12)=\SI{13.5}{\ohm}
\quad\text{e}\quad
(3+6)=\SI{9}{\ohm},
\]
em paralelo: $\dfrac{13{,}5(9)}{22{,}5}=\SI{5.4}{\ohm}$. Somando $R_S$:
\[
R_{eq}=3+5{,}4=\SI{8.4}{\ohm}.
\]
(c) $I=\dfrac{84}{8{,}4}=\SI{10}{\ampere}$.}
\end{questao}

\begin{questao}[3]{}
Resolva a mesma ponte por \textbf{análise nodal} e confronte com o resultado da
conversão.
\resposta{$V_A=\SI{48}{\volt}$, $V_B=\SI{36}{\volt}$; $R_{eq}=\SI{8.4}{\ohm}$
--- confere}
\resolucao{Com $V_S=\SI{84}{\volt}$ conhecido:
\[
\begin{cases}
V_A\left(\tfrac16+\tfrac1{12}+\tfrac16\right)-\tfrac16V_B=\tfrac{84}{6}=14\\[4pt]
-\tfrac16V_A+V_B\left(\tfrac1{12}+\tfrac16+\tfrac16\right)=\tfrac{84}{12}=7
\end{cases}
\]
Resolvendo: $V_A=\SI{48}{\volt}$ e $V_B=\SI{36}{\volt}$.\\
\textbf{Corrente total} (a que sai de $S$):
\[
I=\frac{84-48}{6}+\frac{84-36}{12}=6+4=\SI{10}{\ampere}
\;\Longrightarrow\;
R_{eq}=\frac{84}{10}=\SI{8.4}{\ohm}.\ \checkmark
\]
\textbf{Corrente no ramo central:} $(48-36)/6=\SI{2}{\ampere}$, de $A$ para
$B$ --- confirmando o desequilíbrio.\\
\textbf{Vantagem da conversão:} ela entrega $R_{eq}$ sem resolver sistema
algum. \textbf{Vantagem da nodal:} ela entrega \emph{também} os potenciais
internos e a corrente do ramo central --- que a conversão destrói. Escolha
conforme o que se precisa saber.}
\end{questao}

\begin{questao}[3]{}
Um triângulo tem um lado muito menor que os outros dois:
$R_{AB}=\SI{0.1}{\ohm}$, $R_{BC}=R_{CA}=\SI{100}{\ohm}$.
\begin{itensq}
\item Determine a estrela equivalente.
\item Compare $R_A+R_B$ com $R_{AB}$ e interprete.
\end{itensq}
\resposta{$R_A=R_B=\SI{49.98}{\milli\ohm}$, $R_C=\SI{49.98}{\ohm}$;
$R_A+R_B\approx R_{AB}$}
\resolucao{(a) $\Sigma=0{,}1+100+100=\SI{200.1}{\ohm}$:
\[
R_A=R_B=\frac{0{,}1(100)}{200{,}1}=\SI{49.98}{\milli\ohm};
\qquad
R_C=\frac{100(100)}{200{,}1}=\SI{49.98}{\ohm}.
\]
(b) $R_A+R_B=\SI{99.95}{\milli\ohm}\approx R_{AB}=\SI{0.1}{\ohm}$.\\
\textbf{Por quê:} a resistência entre $A$ e $B$ com $C$ aberto vale
$R_{AB}\parallel(R_{BC}+R_{CA})=0{,}1\parallel200\approx\SI{0.1}{\ohm}$, e na
estrela ela é exatamente $R_A+R_B$. O lado pequeno domina o paralelo, e a soma
dos dois braços apenas reproduz esse domínio.\\
\textbf{Consequência prática:} quando um lado do triângulo é muito menor que os
outros, a estrela resultante tem \emph{dois braços minúsculos e um braço
médio}. Os dois braços pequenos costumam ser desprezíveis diante do restante da
rede --- mas verifique antes de descartá-los, porque em redes de baixa
resistência eles podem ser exatamente o que importa.}
\end{questao}

\begin{questao}[3]{}
Uma estrela tem um braço muito pequeno: $R_A=R_B=\SI{10}{\ohm}$ e
$R_C=\SI{0.1}{\ohm}$.
\begin{itensq}
\item Determine o triângulo equivalente.
\item Analise o comportamento quando $R_C\to0$.
\end{itensq}
\resposta{$R_{AB}=\SI{1020}{\ohm}$, $R_{BC}=R_{CA}=\SI{10.2}{\ohm}$;
$R_{AB}\to\infty$}
\resolucao{(a) $\Sigma=10(10)+10(0{,}1)+0{,}1(10)=100+1+1
=\SI{102}{\ohm\squared}$:
\[
R_{AB}=\frac{102}{0{,}1}=\SI{1020}{\ohm};
\qquad
R_{BC}=R_{CA}=\frac{102}{10}=\SI{10.2}{\ohm}.
\]
(b) Com $R_C\to0$, o numerador tende a $R_AR_B=100$ e
\[
R_{AB}=\frac{100}{R_C}\to\infty,
\]
enquanto $R_{BC}\to R_B$ e $R_{CA}\to R_A$.\\
\textbf{Interpretação:} $R_C=0$ significa que o terminal $C$ coincide com o nó
central. No triângulo, isso equivale a $A$ e $B$ se ligarem a $C$ diretamente
por $R_A$ e $R_B$, sem lado direto entre eles --- e $R_{AB}=\infty$ é
exatamente a ausência desse lado. \checkmark\\
\textbf{Nota numérica:} $\SI{1020}{\ohm}$ em paralelo com $\SI{10.2}{\ohm}$
contribui com menos de $1\%$. Nesses casos, o lado enorme pode ser descartado
sem prejuízo --- e vale conferir se descartá-lo simplifica a rede antes de
seguir com o valor exato.}
\end{questao}

\begin{questao}[3]{}
Discuta o condicionamento das fórmulas de conversão para resistores positivos.
\begin{itensq}
\item Identifique se há risco de cancelamento catastrófico.
\item Determine onde, na prática, a precisão realmente se perde.
\end{itensq}
\resposta{As fórmulas são bem condicionadas; a perda de precisão ocorre na
\emph{redução posterior} da rede}
\resolucao{(a) \textbf{Não há risco.} Ambas as fórmulas envolvem apenas
\emph{produtos} e \emph{somas} de quantidades positivas, seguidos de uma
divisão por um número positivo. Não ocorre subtração de valores próximos ---
que é a única fonte de cancelamento catastrófico. Erros relativos de entrada se
propagam com fator próximo de $1$.\\
\textbf{Verificação:} arredondando a estrela do exemplo anterior
($49{,}98\,\mathrm{m}$; $49{,}98\,\mathrm{m}$; $\SI{49.98}{\ohm}$) para três
algarismos e convertendo de volta, recupera-se
$R_{AB}=\SI{0.1000}{\ohm}$ --- erro inferior a $0{,}1\%$.\\
(b) \textbf{A precisão se perde depois.} Situações típicas:
\begin{itemize}
\item Um braço convertido, muito grande, fica em paralelo com um pequeno: sua
contribuição é desprezível, e carregar seus algarismos é inútil --- mas
\emph{descartá-lo cedo demais} pode mascarar um efeito real.
\item O resultado final é uma \emph{diferença} entre valores próximos --- por
exemplo, ao calcular a corrente de uma ponte quase equilibrada, em que
$V_A-V_B$ é pequeno diante de $V_A$.
\end{itemize}
\textbf{Regra prática:} conduza a conversão com todos os algarismos e arredonde
\emph{apenas no resultado final}. E desconfie sempre que a resposta for uma
pequena diferença de números grandes --- ali, e não na transformação, é onde a
precisão evapora.}
\end{questao}

\begin{questao}[3]{}
Uma conversão produz um triângulo de $36{,}67$, $110$ e $\SI{55}{\ohm}$.
Escolha valores comerciais E24 e avalie o erro.
\begin{itensq}
\item Indique os valores adotados.
\item Converta de volta e compare com a estrela original de $10$, $20$ e
      $\SI{30}{\ohm}$.
\end{itensq}
\resposta{$36$, $110$ e $\SI{56}{\ohm}$; erros de $-0{,}2\%$, $-2{,}0\%$ e
$+1{,}7\%$ nos braços}
\resolucao{(a) Valores E24 mais próximos: $\SI{36}{\ohm}$, $\SI{110}{\ohm}$ e
$\SI{56}{\ohm}$ (desvios de $-1{,}8\%$, $0\%$ e $+1{,}8\%$ nos lados).\\
(b) Convertendo o triângulo comercial de volta a estrela
($\Sigma=36+110+56=\SI{202}{\ohm}$):
\[
R_A=\frac{36(56)}{202}=\SI{9.98}{\ohm};
\quad
R_B=\frac{36(110)}{202}=\SI{19.60}{\ohm};
\quad
R_C=\frac{110(56)}{202}=\SI{30.50}{\ohm}.
\]
Erros: $-0{,}2\%$, $-2{,}0\%$ e $+1{,}7\%$ em relação a $10$, $20$ e
$\SI{30}{\ohm}$.\\
\textbf{Observação importante:} os erros dos braços \emph{não} reproduzem os
erros dos lados --- eles se recombinam. Um lado com erro nulo
($\SI{110}{\ohm}$) contribuiu para dois braços que erraram $2\%$.\\
\textbf{Método correto de avaliação:} nunca julgue o erro pelos componentes
isolados. Converta de volta e compare com o alvo, ou --- melhor ainda --- calcule
a grandeza que realmente interessa (a resistência entre os terminais usados) nos
dois casos.}
\end{questao}

\begin{questao}[3]{}
Compare o esforço de resolver a ponte desequilibrada por conversão
estrela-triângulo e por análise nodal.
\begin{itensq}
\item Enumere as etapas de cada método.
\item Indique o critério de escolha.
\end{itensq}
\resposta{Conversão: 3 fórmulas e reduções; nodal: sistema $2\times2$. A
conversão dá só $R_{eq}$; a nodal dá tudo}
\resolucao{(a) \textbf{Conversão:} identificar o triângulo (1); aplicar três
fórmulas (2); somar duas séries (3); um paralelo (4); uma soma final (5). Sem
sistemas de equações.\\
\textbf{Nodal:} montar a matriz $2\times2$ (1); resolver (2); calcular a
corrente total a partir dos potenciais (3).\\
(b) \textbf{Critério.}
\begin{itemize}
\item Se a pergunta é \textbf{só} $R_{eq}$ --- ou a corrente total --- a
conversão é mais rápida e não exige álgebra linear.
\item Se são necessários potenciais internos, correntes de ramo ou potências
individuais, use \textbf{nodal}. A conversão destrói a estrutura interna: o nó
central da estrela não existe no circuito real, e nenhuma corrente calculada
nele tem significado físico.
\item Para pontes com \emph{mais} de um triângulo irredutível, a conversão
precisa ser aplicada repetidamente e perde a vantagem; a nodal escala melhor.
\end{itemize}
\textbf{Estratégia mista frequente:} converter para obter $R_{eq}$ e a corrente
total, e depois voltar ao circuito \emph{original} para distribuir essa corrente
--- aproveitando o melhor dos dois.}
\end{questao}

\begin{questao}[3]{}
Uma rede de três terminais é formada por resistores positivos em uma topologia
qualquer. Discuta se ela sempre admite equivalente em estrela.
\begin{itensq}
\item Estabeleça a condição.
\item Dê um exemplo em que a estrela equivalente teria braço negativo.
\end{itensq}
\resposta{A condição é a desigualdade triangular sobre as resistências entre
pares}
\resolucao{(a) Sejam $R_{ab}$, $R_{bc}$ e $R_{ca}$ as resistências medidas
entre cada par de terminais (com o terceiro aberto). Da estrela, valem
$R_{ab}=R_A+R_B$, $R_{bc}=R_B+R_C$ e $R_{ca}=R_C+R_A$, de onde
\[
R_A=\frac{R_{ab}+R_{ca}-R_{bc}}{2},
\]
e analogamente para $R_B$ e $R_C$.\\
\textbf{Condição:} os três braços são não negativos se e somente se cada soma
de duas resistências medidas for maior ou igual à terceira --- exatamente a
\textbf{desigualdade triangular}.\\
(b) \textbf{Contraexemplo.} Suponha
$R_{ab}=\SI{1}{\ohm}$, $R_{bc}=\SI{10}{\ohm}$ e
$R_{ca}=\SI{1}{\ohm}$. Então
\[
R_B=\frac{1+10-1}{2}=\SI{5}{\ohm};
\qquad
R_A=\frac{1+1-10}{2}=\SI{-4}{\ohm}.
\]
Braço negativo --- fisicamente irrealizável com resistores passivos.\\
\textbf{Mas atenção:} essa combinação de medidas também não é realizável por
uma rede passiva de três terminais. A desigualdade triangular é uma
\emph{propriedade} de tais redes, não uma restrição adicional. Braços negativos
só surgem como artifício de cálculo --- em transformações intermediárias ou em
redes com elementos ativos.}
\end{questao}

\blocodesafio

\begin{questao}[4]{}
\textbf{(Dedução --- triângulo para estrela.)} Deduza as fórmulas de conversão a
partir das resistências entre terminais.
\begin{itensq}
\item Escreva as três equações de equivalência.
\item Resolva o sistema para $R_A$, $R_B$ e $R_C$.
\item Obtenha a forma final em função dos lados do triângulo.
\end{itensq}
\resposta{$R_A=\dfrac{R_{AB}R_{CA}}{R_{AB}+R_{BC}+R_{CA}}$ e análogas}
\resolucao{(a) A equivalência exige que a resistência entre cada par de
terminais, com o terceiro \emph{aberto}, seja a mesma nas duas configurações:
\[
\begin{cases}
R_A+R_B=R_{AB}\parallel(R_{BC}+R_{CA})\\[3pt]
R_B+R_C=R_{BC}\parallel(R_{CA}+R_{AB})\\[3pt]
R_C+R_A=R_{CA}\parallel(R_{AB}+R_{BC})
\end{cases}
\]
(b) Somando as três e dividindo por $2$, obtém-se $R_A+R_B+R_C$. Subtraindo
dessa soma a segunda equação, resta
\[
R_A=\frac{(R_A+R_B)+(R_C+R_A)-(R_B+R_C)}{2}.
\]
(c) Substituindo os paralelos, com $\Sigma=R_{AB}+R_{BC}+R_{CA}$:
\[
R_A+R_B=\frac{R_{AB}(R_{BC}+R_{CA})}{\Sigma};
\qquad
R_C+R_A=\frac{R_{CA}(R_{AB}+R_{BC})}{\Sigma};
\]
\[
R_B+R_C=\frac{R_{BC}(R_{CA}+R_{AB})}{\Sigma}.
\]
Levando à expressão de $R_A$ e simplificando o numerador:
\[
R_{AB}R_{BC}+R_{AB}R_{CA}+R_{CA}R_{AB}+R_{CA}R_{BC}
-R_{BC}R_{CA}-R_{BC}R_{AB}=2R_{AB}R_{CA},
\]
donde
\[
R_A=\frac{2R_{AB}R_{CA}}{2\Sigma}=\frac{R_{AB}R_{CA}}{\Sigma}.
\]
\textbf{A regra mnemônica agora tem origem:} o produto dos dois lados que tocam
o terminal aparece porque os demais termos se cancelam aos pares na combinação
$(R_{ab}+R_{ca}-R_{bc})/2$.}
\end{questao}

\begin{questao}[4]{}
\textbf{(Dedução --- estrela para triângulo.)} Obtenha a transformação inversa.
\begin{itensq}
\item Deduza $R_{AB}$ em função dos braços.
\item Verifique que a composição das duas transformações é a identidade.
\item Explique por que o denominador é o braço \emph{oposto}.
\end{itensq}
\resposta{$R_{AB}=\dfrac{R_AR_B+R_BR_C+R_CR_A}{R_C}$}
\resolucao{(a) Tomando as três fórmulas diretas e formando a soma dos produtos
dois a dois:
\[
R_AR_B+R_BR_C+R_CR_A
=\frac{R_{AB}R_{BC}R_{CA}\left(R_{AB}+R_{BC}+R_{CA}\right)}{\Sigma^{2}}
=\frac{R_{AB}R_{BC}R_{CA}}{\Sigma}.
\]
Chamando esse valor de $P$ e observando que
$R_C=\dfrac{R_{BC}R_{CA}}{\Sigma}$, tem-se
\[
\frac{P}{R_C}
=\frac{R_{AB}R_{BC}R_{CA}/\Sigma}{R_{BC}R_{CA}/\Sigma}=R_{AB}.
\]
Logo
\[
R_{AB}=\frac{R_AR_B+R_BR_C+R_CR_A}{R_C}.
\]
(b) A dedução \emph{é} a verificação: partiu-se das fórmulas diretas e
chegou-se de volta aos lados originais. Numericamente, o exemplo
$20/30/\SI{50}{\ohm}\to10/6/\SI{15}{\ohm}\to20/30/\SI{50}{\ohm}$ confirma.\\
(c) \textbf{Por que o oposto.} O braço $R_C$ é o único que \emph{não} contém
$R_{AB}$ em seu numerador --- ele é construído a partir de $R_{BC}$ e
$R_{CA}$. Ao dividir $P$ (que contém o produto dos três lados) por $R_C$,
cancelam-se exatamente $R_{BC}$ e $R_{CA}$, isolando $R_{AB}$. É por isso que
o denominador tem de ser o braço oposto, e não outro.}
\end{questao}

\begin{questao}[4]{}
\textbf{(Positividade.)} Prove que as transformações preservam a positividade
dos elementos.
\begin{itensq}
\item Mostre que $\Delta\to Y$ com lados positivos gera braços positivos.
\item Mostre que $Y\to\Delta$ com braços positivos gera lados positivos.
\item Determine o que ocorre quando um elemento tende a zero ou a infinito.
\end{itensq}
\resposta{Ambas as fórmulas são quocientes de quantidades positivas ---
positividade garantida}
\resolucao{(a) $R_A=\dfrac{R_{AB}R_{CA}}{R_{AB}+R_{BC}+R_{CA}}$. Com todos os
lados positivos, numerador e denominador são positivos, logo $R_A>0$. O mesmo
para $R_B$ e $R_C$.\\
\textbf{Além disso, há um limite superior:} como
$\Sigma>R_{CA}$, segue $R_A<R_{AB}$; e como $\Sigma>R_{AB}$, segue
$R_A<R_{CA}$. Ou seja, \textbf{cada braço é menor que os dois lados que o
originaram} --- teste de sanidade imediato para qualquer conversão.\\
(b) $R_{AB}=\dfrac{R_AR_B+R_BR_C+R_CR_A}{R_C}$: soma de três produtos positivos
dividida por um positivo, logo $R_{AB}>0$.\\
\textbf{E há um limite inferior:} $R_{AB}>\dfrac{R_AR_C+R_BR_C}{R_C}=R_A+R_B$
--- \textbf{cada lado é maior que a soma dos dois braços correspondentes}.
Segundo teste de sanidade, dual do anterior.\\
(c) \textbf{Casos-limite.}
\begin{itemize}
\item $R_{AB}\to0$ (lado em curto) $\Rightarrow$ $R_A,R_B\to0$: os terminais
$A$ e $B$ colapsam sobre o nó central.
\item $R_{AB}\to\infty$ (lado aberto) $\Rightarrow$ a estrela degenera em uma
série simples.
\item $R_C\to0$ na estrela $\Rightarrow$ $R_{AB}\to\infty$: o lado direto
desaparece.
\item $R_C\to\infty$ $\Rightarrow$ $R_{AB}\to R_A+R_B$, e os outros dois lados
divergem.
\end{itemize}
Em todos os casos, o limite reproduz a interpretação física esperada --- as
fórmulas nunca produzem resultados sem sentido para entradas positivas.}
\end{questao}

\begin{questao}[4]{}
\textbf{(Alcance da equivalência.)} Analise \emph{o que} a transformação
preserva e o que ela destrói.
\begin{itensq}
\item Prove que as resistências entre os três terminais são preservadas.
\item Determine se correntes e potências internas são preservadas.
\item Estabeleça o procedimento correto quando grandezas internas são
      necessárias.
\end{itensq}
\resposta{Preserva o comportamento nos três terminais; \emph{não} preserva nada
do interior}
\resolucao{(a) Por construção, as três equações de equivalência impõem
igualdade das resistências entre pares. Como a rede de três terminais é linear e
tem apenas três acessos, seu comportamento externo é completamente descrito por
esses três valores --- qualquer excitação aplicada aos terminais produz a mesma
resposta nas duas configurações.\\
\emph{(Formalmente: a matriz de condutâncias de três terminais tem três
parâmetros independentes, e as três equações os fixam todos.)}\\
(b) \textbf{Não.} A estrela possui um nó central inexistente no triângulo, e o
triângulo possui um lado direto inexistente na estrela. Consequências:
\begin{itemize}
\item As correntes dos elementos são \emph{completamente diferentes}.
\item As potências individuais também --- embora a potência \emph{total} seja a
mesma, pois o comportamento externo é idêntico.
\item O potencial do nó central não corresponde a ponto algum do triângulo.
\end{itemize}
\textbf{Exemplo:} na ponte resolvida anteriormente, o resistor central de
$\SI{6}{\ohm}$ conduz $\SI{2}{\ampere}$. Após a conversão, esse resistor
\emph{não existe} --- e nenhum braço da estrela conduz $\SI{2}{\ampere}$.\\
(c) \textbf{Procedimento correto:} use a conversão para obter as grandezas
\emph{externas} (resistência equivalente, corrente total, tensões nos
terminais). Depois \textbf{volte ao circuito original} e, com essas grandezas já
conhecidas, distribua correntes e tensões pelos elementos reais.\\
Se muitas grandezas internas forem necessárias, a análise nodal costuma ser
mais direta desde o início.}
\end{questao}

\begin{questao}[4]{}
\textbf{(Projeto com valores comerciais.)} Uma estrela de $10$, $20$ e
$\SI{30}{\ohm}$ deve ser substituída por um triângulo montado com resistores
E24.
\begin{itensq}
\item Determine o triângulo exato.
\item Escolha valores comerciais e avalie o erro na resistência entre $A$ e
      $B$.
\item Proponha uma estratégia para reduzir o erro.
\end{itensq}
\resposta{Exato: $36{,}67$, $110$ e $\SI{55}{\ohm}$; com $36/110/56$, o erro em
$R_{AB}$ é de cerca de $1\%$}
\resolucao{(a) $\Sigma=10(20)+20(30)+30(10)=200+600+300=\SI{1100}{\ohm\squared}$:
\[
R_{AB}=\frac{1100}{30}=\SI{36.67}{\ohm};
\quad
R_{BC}=\frac{1100}{10}=\SI{110}{\ohm};
\quad
R_{CA}=\frac{1100}{20}=\SI{55}{\ohm}.
\]
(b) E24: $\SI{36}{\ohm}$, $\SI{110}{\ohm}$, $\SI{56}{\ohm}$.\\
\textbf{Grandeza que interessa} --- resistência entre $A$ e $B$ com $C$ aberto:
\[
\text{alvo (estrela)}: R_A+R_B=\SI{30}{\ohm};
\]
\[
\text{obtido}: 36\parallel(110+56)=\frac{36(166)}{202}=\SI{29.58}{\ohm}.
\]
Erro de $-1{,}4\%$ --- aceitável para a maioria das aplicações, e menor que a
tolerância de $5\%$ dos próprios resistores.\\
(c) \textbf{Estratégias, em ordem de custo.}
\begin{itemize}
\item \textbf{Testar combinações:} $39/110/56$ daria
$39\parallel166=\SI{31.6}{\ohm}$ ($+5{,}3\%$) --- pior. O par
$36/110/56$ já é bom.
\item \textbf{Série E96} ($1\%$): contém $36{,}5$ e $54{,}9$, reduzindo o erro
a frações de por cento.
\item \textbf{Associação:} $\SI{36.67}{\ohm}$ sai de $\SI{33}{\ohm}$ em série
com $\SI{3.6}{\ohm}$, ou de $\SI{110}{\ohm}$ em paralelo com
$\SI{55}{\ohm}$ --- e note que esses dois valores já estão na lista!
\end{itemize}
\textbf{Critério de avaliação:} sempre meça o erro na \emph{grandeza de
interesse}, não nos componentes. Erros de lados individuais podem se compensar
ou se somar, e só o cálculo da resistência terminal revela o efeito real.}
\end{questao}

\begin{questao}[4]{}
\textbf{(Ponte quase equilibrada.)} Uma ponte tem $R_1=\SI{100}{\ohm}$,
$R_2=\SI{100}{\ohm}$, $R_3=\SI{100}{\ohm}$ e
$R_4=\SI{101}{\ohm}$, com $R_g=\SI{100}{\ohm}$ e alimentação de
$\SI{10}{\volt}$.
\begin{itensq}
\item Estime a corrente do galvanômetro por análise nodal.
\item Discuta a viabilidade de obter esse valor por conversão
      estrela-triângulo.
\end{itensq}
\resposta{$I_g\approx\SI{-124}{\micro\ampere}$ (de $B$ para $A$); a conversão é
inadequada para esta grandeza}
\resolucao{(a) Com $V_S=\SI{10}{\volt}$ e o terra na base, o sistema nodal
\[
\begin{cases}
V_A\left(\tfrac{1}{100}+\tfrac{1}{100}+\tfrac{1}{100}\right)
-\tfrac{V_B}{100}=\tfrac{10}{100}\\[4pt]
-\tfrac{V_A}{100}
+V_B\left(\tfrac{1}{100}+\tfrac{1}{101}+\tfrac{1}{100}\right)=\tfrac{10}{100}
\end{cases}
\]
dá $V_A=\SI{5.00621}{\volt}$ e $V_B=\SI{5.01863}{\volt}$, de onde
\[
I_g=\frac{V_A-V_B}{100}=\frac{-0{,}01242}{100}=\SI{-124}{\micro\ampere}.
\]
O sinal indica corrente de $B$ para $A$ --- coerente, pois $R_4>R_3$ eleva
$V_B$.\\
(b) \textbf{A conversão é inadequada, por duas razões independentes.}\\
\emph{(i) Ela não fornece a grandeza pedida.} Após converter o triângulo, o
resistor do galvanômetro deixa de existir como elemento --- e é justamente sua
corrente que se quer. Seria necessário converter, obter $R_{eq}$ e a corrente
total, voltar ao circuito original e só então distribuir --- caminho mais longo
que a nodal direta.\\
\emph{(ii) O problema é mal condicionado.} A resposta é a diferença
$V_A-V_B=\SI{-12.4}{\milli\volt}$ entre dois valores próximos de
$\SI{5}{\volt}$ --- uma diferença relativa de apenas $0{,}25\%$.
\begin{center}
\begin{tabular}{ccc}
\hline
Algarismos usados & $V_A-V_B$ obtido & erro em $I_g$\\
\hline
$6$ ($5{,}00621$; $5{,}01863$) & $\SI{-12.42}{\milli\volt}$ & $<0{,}1\%$\\
$4$ ($5{,}006$; $5{,}019$) & $\SI{-13}{\milli\volt}$ & $\approx5\%$\\
$3$ ($5{,}01$; $5{,}02$) & $\SI{-10}{\milli\volt}$ & $\approx20\%$\\
\hline
\end{tabular}
\end{center}
Arredondar cedo destrói o resultado.\\
\textbf{Método adequado:} tratar o desequilíbrio como \emph{perturbação}. Para
$\delta=\Delta R/R$ pequeno, $I_g$ é aproximadamente proporcional a $\delta$, e
a expressão linearizada evita completamente a subtração de valores próximos.\\
\textbf{Moral:} escolher o método é também escolher \emph{como} o erro numérico
se propaga. A conversão estrela-triângulo é excelente para $R_{eq}$ e péssima
para desequilíbrios pequenos.}
\end{questao}

\begin{questao}[4]{}
\textbf{(Síntese.)} Uma rede de três terminais deve apresentar
$R_{AB}=\SI{20}{\ohm}$, $R_{BC}=\SI{30}{\ohm}$ e $R_{CA}=\SI{40}{\ohm}$
(medidas com o terceiro terminal aberto).
\begin{itensq}
\item Verifique a realizabilidade e determine a estrela que atende.
\item Determine o triângulo equivalente.
\item Escolha valores comerciais e avalie qual das duas topologias erra menos.
\end{itensq}
\resposta{$R_A=\SI{15}{\ohm}$, $R_B=\SI{5}{\ohm}$, $R_C=\SI{25}{\ohm}$;
triângulo $23$, $38{,}33$ e $\SI{115}{\ohm}$; a estrela é preferível}
\resolucao{(a) \textbf{Realizabilidade.} A desigualdade triangular exige que
cada resistência medida seja menor que a soma das outras duas:
$20<70$, $30<60$, $40<50$ \checkmark. A rede é realizável com resistores
positivos.\\
Invertendo o sistema $R_A+R_B=R_{AB}$ e análogos:
\[
R_A=\frac{20+40-30}{2}=\SI{15}{\ohm};
\quad
R_B=\frac{20+30-40}{2}=\SI{5}{\ohm};
\quad
R_C=\frac{30+40-20}{2}=\SI{25}{\ohm}.
\]
\textbf{Conferência:} $15+5=20$; $5+25=30$; $25+15=40$ \checkmark.\\
(b) $\Sigma=15(5)+5(25)+25(15)=75+125+375=\SI{575}{\ohm\squared}$:
\[
R_{AB}=\frac{575}{25}=\SI{23}{\ohm};
\quad
R_{BC}=\frac{575}{15}=\SI{38.33}{\ohm};
\quad
R_{CA}=\frac{575}{5}=\SI{115}{\ohm}.
\]
(c) \textbf{Estrela com E24:} $\SI{15}{\ohm}$ e $\SI{24}{\ohm}$ existem, e o
mais próximo de $\SI{5}{\ohm}$ é $\SI{5.1}{\ohm}$. Com
$(15;\,5{,}1;\,24)$:
\[
R_{AB}=\SI{20.1}{\ohm}\ (+0{,}5\%);
\quad
R_{BC}=\SI{29.1}{\ohm}\ (-3{,}0\%);
\quad
R_{CA}=\SI{39.0}{\ohm}\ (-2{,}5\%).
\]
\textbf{Triângulo com E24:} $(24;\,39;\,110)$:
\[
R_{AB}=\SI{20.67}{\ohm}\ (+3{,}4\%);
\quad
R_{BC}=\SI{30.21}{\ohm}\ (+0{,}7\%);
\quad
R_{CA}=\SI{40.06}{\ohm}\ (+0{,}1\%).
\]
\textbf{Conclusão.} Os dois erram no mesmo patamar (pior caso de $3{,}0\%$
contra $3{,}4\%$), mas a \textbf{estrela é preferível} por três razões: seus
valores são menores (componentes mais baratos), o erro é dominado por um único
braço difícil ($\SI{5}{\ohm}$), e cada resistência medida depende de
\emph{apenas dois} braços --- o que torna o ajuste previsível.\\
No triângulo, cada medida envolve os \emph{três} lados através de um paralelo, e
os erros se recombinam de forma difícil de controlar.\\
\textbf{Se a precisão for crítica:} o braço de $\SI{5}{\ohm}$ pode ser obtido
associando $\SI{10}{\ohm}$ em paralelo com $\SI{10}{\ohm}$, eliminando o
único desvio significativo.}
\end{questao}

\begin{questao}[4]{}
\textbf{(Generalização.)} A transformação estrela-triângulo é o caso $n=3$ de
uma família mais ampla.
\begin{itensq}
\item Descreva a transformação estrela-polígono para $n$ terminais.
\item Determine quantos elementos há em cada configuração.
\item Explique por que a transformação inversa nem sempre existe para $n>3$.
\end{itensq}
\resposta{Estrela: $n$ braços; polígono: $n(n-1)/2$ lados. A inversa só existe
para $n=3$}
\resolucao{(a) Uma estrela de $n$ braços $R_1,\dots,R_n$ ligados a um nó comum
equivale a um polígono completo em que o lado entre os terminais $i$ e $j$ vale
\[
R_{ij}=R_iR_j\sum_{k=1}^{n}\frac{1}{R_k}.
\]
Para $n=3$ isso reproduz a fórmula conhecida:
$R_{12}=R_1R_2\left(\frac1{R_1}+\frac1{R_2}+\frac1{R_3}\right)
=\frac{R_1R_2+R_2R_3+R_3R_1}{R_3}$ \checkmark.\\
(b) \textbf{Estrela:} $n$ elementos. \textbf{Polígono completo:}
$\dbinom{n}{2}=\dfrac{n(n-1)}{2}$ elementos.
\begin{center}
\begin{tabular}{ccc}
\hline
$n$ & estrela & polígono\\
\hline
$3$ & $3$ & $3$\\
$4$ & $4$ & $6$\\
$5$ & $5$ & $10$\\
\hline
\end{tabular}
\end{center}
(c) \textbf{Contagem de graus de liberdade.} A transformação
$Y\to$ polígono sempre existe. A inversa exigiria determinar $n$ braços a
partir de $n(n-1)/2$ lados --- um sistema \textbf{sobredeterminado} para
$n>3$.\\
Só quando $n=3$ os dois números coincidem ($3=3$), e a correspondência é
biunívoca. Para $n=4$, seria preciso obter $4$ incógnitas de $6$ equações: um
polígono \emph{genérico} de quatro terminais simplesmente não admite estrela
equivalente.\\
\textbf{Consequência prática:} a conversão $\Delta\to Y$ é uma ferramenta
peculiar ao caso de três terminais. Redes de quatro ou mais acessos exigem
descrição matricial completa --- e é por isso que a análise nodal, que não
depende de topologias especiais, é o método geral.}
\end{questao}

% END BANCO COMPLEMENTAR

\gabarito[3.4 Transformação estrela-triângulo]
